\part{Shifted fibre products}
\section{Shifted fibre products of two six-spheres}\label{sec:shifted}

The same sphere fibrations give compact complex fivefolds by placing
their exceptional fibres over disjoint subsets of the base. Their topology
is governed by a torus-knot peripheral group and the integral cusp maps.

\begin{definition}[Shifted fibre product]\label{def:V:shifted-product}
Let $X$ and $X'$ be canonical marked members of the $(3,4,\infty)$
family, and choose $\varphi\in\operatorname{PGL}_2(\mathbb C)$ so that
the exceptional values of $X$ and $\varphi^*X'$ are disjoint.
Their \emph{shifted fibre product} is $P=X\times_B\varphi^*X'$.
\end{definition}

\begin{theorem}[A shifted second sphere]\hypertarget{tag-02MN}{}\label{maps:src:E.2.3}
\StatementRecord{maps:src:E.2.3}{proved}{02MN}{\SeriesRef{maps:src:E.2.2}, \eqref{maps:eq:E.9}}


Let $X'$ be a second canonical member and choose $\varphi\in
\operatorname{PGL}_2(\mathbb C)$ so that its three exceptional
values are disjoint from those of $X$.  The smooth fivefold
$P=X\times_B\varphi^*X'$ has Euler characteristic zero and
an infinite fundamental group with presentation
\[
\begin{split}
G=\langle a_1,a_2,b_1,b_2,h_1,h_2\mid\ &
 h_1,h_2\ \text{central},\quad
 a_1^3=h_1,\ a_2^4=h_1^{-1},\\
 &b_1^3=h_2,\ b_2^4=h_2^{-1},\quad
 a_1a_2b_1b_2=1\rangle .
\end{split}
\tag{E.16}\hypertarget{tag-029M}{}\label{maps:eq:E.16}
\]
Its abelianization is $\mathbb Z$, and its rational Betti numbers are
\[
 (1,1,1,20,63,88,63,20,1,1,1).
 \tag{E.17}\hypertarget{tag-029N}{}\label{maps:eq:E.17}
\]
For fixed $X,X'$, all admissible $\varphi$ give the same oriented
diffeomorphism type.

\end{theorem}


\begin{proof}
\begin{enumerate}[label={$\langle1\rangle$\arabic*.},leftmargin=*]
\item Theorem \SeriesRef{maps:src:E.2.2} proves smoothness.  At each special value the other
factor is a smooth four-torus.  The Euler characteristic of every
fibre is therefore zero, and Euler additivity gives $\chi(P)=0$.
We apply van Kampen with the cusp relations.
Over the six-punctured sphere the fibre group is
$\Lambda_X\oplus\Lambda_{X'}$ and the base is free on five
peripheral loops.  Each monodromy acts on its own direct summand.
Filling the first cusp kills its angular generators
$\widehat w,\widehat\delta$ and its cusp meridian.  Their normal
closure under $A_1$ also kills $\widehat u$, since
$A_1\widehat w=\widehat u-\widehat w$.  The same calculation
in the other summand leaves only $h_1=\widehat\gamma$ and
$h_2=\widehat\gamma'$.  The finite monodromies now act trivially
on them.  The four finite filling relations are precisely those
in \eqref{maps:eq:E.16}, since $\gamma(v_1)=1$ and $\gamma(v_2)=-1$.
Both cusp collars are untwisted; their meridians are killed.
The product of the remaining four meridians gives the last
relation.  These are all relations introduced by the six
fillings, so van Kampen gives exactly \eqref{maps:eq:E.16}.

\item Quotienting by $h_1,h_2$ gives the orbifold group of the sphere
with cone orders $3,4,3,4$.  Its orbifold Euler characteristic
is $-5/6$, so its universal orbifold cover is the disc and the
group is infinite.  Equivalently one can build the Seifert
$T^2$-bundle with these four primitive translation vectors; over
the disc it is a trivial torus bundle.  Its exact sequence shows
that the central subgroup in \eqref{maps:eq:E.16} is $\mathbb Z^2$.
The abelianized five-by-six relation matrix has Smith factors
$(1,1,1,1,1)$ by integral row and column operations, proving $G_{\rm ab}=\mathbb Z$.

\item For the rational homology there is no extra cusp kernel: each
cusp is the rank-two cusp of the sphere fibration times a smooth torus, and its
specialization is an isomorphism onto invariants over $\mathbb Q$.
The two monodromy groups act on independent summands.  In exterior
degrees $q=0,\ldots,8$, the $H^0$ and $H^2$ rows are both
\[
 (1,2,3,4,5,4,3,2,1),
\]
and the $H^1$ row is
\[
 (0,0,18,60,86,60,18,0,0).
\]
The local invariant polynomials of the first factor are
$A(z)=1+2z+2z^2+2z^3+z^4$ at a finite point and
$C(z)=A(z)+2z^2$ at its cusp.  Hence the Euler polynomial of
the direct-image rows is
\[
 -4(1+z)^8+(4A(z)+2C(z))(1+z)^4.
\]
Formula \eqref{maps:eq:E.9}, with six point strata, gives the displayed rows.

\item Each invariant algebra of a factor has one generator in degrees
zero through four.  Its invariant-to-coinvariant map is nonzero
in even degrees and zero in odd degrees.  In the marking used
above, this follows from the invariant form
$6\gamma\wedge\delta+u\wedge w$ and its exterior products:
the odd invariant classes pair to zero with every invariant
complementary class, and the even ones have nonzero square or
top pairing.  Its four transgressions are all nonzero, as the
factor is a sphere.  The derivation rule therefore sends a tensor
of bidegree $(i,j)$ into bidegrees $(i-1,j)$ and $(i,j-1)$,
with the first term nonzero exactly when $i>0$ and $j$ is even,
and the second exactly when $j>0$ and $i$ is even.  The nonzero
column supports are disjoint except for pairs in a single row,
so their scalar normalizations do not affect ranks.  This gives
\[
 (r_1,\ldots,r_8)=(1,2,2,3,3,2,2,1).
\]
Row reduction of these exterior-product matrices gives the displayed ranks.
There are no later differentials.  Their diagonal sums give
\eqref{maps:eq:E.17}, which also satisfies Poincar\'e duality and Euler zero.

\item Finally, the admissible set of $\varphi$ is the complement of
finitely many proper complex hypersurfaces in the connected
complex manifold $\operatorname{PGL}_2(\mathbb C)$, so it is
path connected.  The universal pullback over this set is a
proper smooth submersion by the pointwise transversality above.  Parallel transport along a path gives an oriented
diffeomorphism of its fibres.
\end{enumerate}
\end{proof}

\begin{proposition}[The shifted product's smooth type]\hypertarget{tag-02MO}{}\label{maps:src:E.2.4}
\StatementRecord{maps:src:E.2.4}{proved}{02MO}{\eqref{maps:eq:E.16}, \SeriesRef{maps:src:E.2.21}, \SeriesRef{maps:src:E.2.1}, \SeriesRef{maps:src:E.2.5}, \SeriesRef{maps:src:E.2.18}}


The shifted product in \eqref{maps:eq:E.16} is the closed oriented ten-manifold
$\mathcal H_{10}$ specified by the $2432$ ordinary framed handles
of \ref{maps:src:E.2.21}. Its regular-fibre complement is, relative to its entire
marked boundary, the $76$-handle six-manifold $\mathcal H_6$ of
that theorem. The full attaching embeddings, their ordered normal
coordinates, and every dual belt embedding are prescribed by
finite coordinate formulas. The identification follows by a
relative diffeomorphism of the cusp, core and finite caps,
and then by the handle-by-handle comparison of \ref{maps:src:E.2.21} using
Milnor's Theorems~1.4, 3.13 and~5.8. 

The smooth type is given by the explicit framed-handle presentation. \ref{maps:src:E.2.18}'s quadratic obstruction is consistent
with this presentation. 


\end{proposition}
\begin{proof}
\begin{enumerate}[label={$\langle1\rangle$\arabic*.},leftmargin=*]
\item The relative diffeomorphism of \cref{maps:src:E.2.20} identifies the cusp, flat core and both finite caps while fixing the full outer torus marking. Apply the ordered block expansion of \cref{maps:src:E.2.21} to this comparison: each full attaching embedding, including its normal coordinates, is carried to the displayed coordinate embedding. The compactly supported collar isotopies in that proof extend the comparison across each successive handle, with the protected boundary fixed. Taking its product with the two regular torus factors and gluing by the explicit normal-angle reversal gives the ten-dimensional comparison. The stated counts are the two expanded handle lists of \cref{maps:src:E.2.21}; every top handle uses its prescribed attaching map. The remaining assertions are exactly the already established results \cref{maps:src:E.2.1}, \SeriesRef{maps:src:E.2.2}, \cref{maps:src:E.2.3}, \cref{maps:src:E.2.5}, \cref{maps:src:E.2.18}.
\end{enumerate}
\end{proof}



\begin{proposition}[The first integral torsion of the shifted product]\hypertarget{tag-02MP}{}\label{maps:src:E.2.5}
\StatementRecord{maps:src:E.2.5}{proved}{02MP}{\SeriesRef{maps:src:E.2.3}, \SeriesRef{maps:src:E.1.3}, \SeriesRef{maps:src:D.1.2}, \eqref{maps:eq:E.I1}, \SeriesRef{maps:src:E.1.5}, \eqref{maps:eq:E.20}}


Let \(Z_\varphi=X\times_B\varphi^*X\) be the transverse
fivefold in \ref{maps:src:E.2.3}, with the two sets of three critical values
disjoint. Its second integral homology is
\[
 H_2(Z_\varphi;\mathbb Z)
   =\mathbb Z\oplus(\mathbb Z/12)^7.
\]
\end{proposition}


\begin{proof}
\begin{enumerate}[label={$\langle1\rangle$\arabic*.},leftmargin=*]
\item Use subscripts one and two for the two rank-four cohomology
local systems. At a critical value of one factor the other
factor is a smooth torus bundle over a disc, hence a product
there. Each local specialization is therefore the tensor
product of the original specialization with the identity on
the exterior algebra of \(\mathbb Z^4\). The original finite
reduced surfaces and the original cusp have free cohomology.
Their specialization maps are injective; at the cusp they
are integral isomorphisms onto invariants. For the latter,
apply the support-function construction of \ref{maps:src:E.1.3} in two
coordinates. It supplies the invariant deck classes
\(b_1,b_2\), the three classes
\(a_1b_1,a_2b_2,a_1b_2+a_2b_1\), and their products.
Together with \(b_1b_2\) these give an integral kernel
basis in degree two; their products give integral bases
in degrees three and four. Both the central-fibre and
invariant lattices have ranks \((1,2,4,2,1)\), as follows
also from the normalization matrix in \SeriesRef{maps:src:D.1.2} at \(d=1\).
Thus the specialization is onto between free groups of
equal ranks in every degree, proving the claimed isomorphism.
Thus the only difference between \(R^qg_*\mathbb Z\) and
\(\mathcal J_q=j_*\bigwedge^q(\mathcal V_1\oplus\mathcal V_2)\)
is their finite covering cokernel, with no point-supported kernel.

\item The finite Wang calculation \eqref{maps:eq:E.I1}, now with the original
three-dimensional vertical lattice, gives
\[
 C_3^1=\mathbb Z/3,\quad C_4^1=\mathbb Z/4,
 \qquad C_3^2=0,\quad C_4^2=\mathbb Z/2.
\]
The invariant covector \(\gamma\) generates the first two
cokernels, and the invariant two-form \(q_X\) generates the
last, by the same even-intersection argument as in \ref{maps:src:E.1.5}.
Consequently the degree-two covering cokernel at all four
finite values, combined, is
\[
 C^2=\left((\mathbb Z/3)^4\oplus
                 (\mathbb Z/4)^4\oplus\mathbb Z/2\right)^2.
\]
The two monodromy groups act independently on the two factors.
The three global invariant two-forms have integral basis
\(q_1,\gamma_1\gamma_2,q_2\). The first and third map
onto the respective \(\mathbb Z/2\) summands and vanish
in the other summands. The middle form maps to the primitive
\(\gamma\) coordinate of each of the four order-three or
order-four blocks. After the two order-two summands have
been removed, it therefore generates a direct cyclic summand
of order twelve in \((\mathbb Z/12)^8\). Hence
\[
 \operatorname{coker}\{H^0(B,\mathcal J_2)\to C^2\}
      =(\mathbb Z/12)^7.
 \tag{E.20}\hypertarget{tag-029P}{}\label{maps:eq:E.20}
\]
The eighteen-row integer presentation of these covering cokernels has
Smith form with eleven unit factors and seven factors twelve.

\item In degree two the global coboundary image is primitive.
Thus the parabolic group, a subgroup of the free-group
cohomology on the punctured base, is torsion free. Its rank
is eighteen: the six local invariant ranks are
\((16,16,18,16,16,18)\), and its global invariant and
coinvariant ranks are both three. The Euler formula on the
six-punctured sphere gives \(3+3+4\cdot28-100=18\).
These ranks and the primitivity are computed with exact exterior matrices. The sheaf sequence
\(0\to R^2g_*\mathbb Z\to\mathcal J_2\to C^2\to0\)
and \eqref{maps:eq:E.20} now give
\[
 H^1(B,R^2g_*\mathbb Z)
     =\mathbb Z^{18}\oplus(\mathbb Z/12)^7.
\]
Here the quotient by its finite subgroup is free, so this
abelian-group extension splits.

\item The group \(H^2(B,R^1g_*\mathbb Z)\) is the free
rank-two group of coinvariants, one primitive summand from
each factor. Pulling back the integral degree-two
transgression of the original sphere fibration through each
projection kills these two summands separately and
surjectively. Thus \(E_\infty^{2,1}=0\) in the Leray
sequence of \(g\). Its other middle group
\(E_\infty^{0,3}\) is free, since \(R^3g_*\mathbb Z\)
injects into \(\mathcal J_3\). There are no higher
differentials. Consequently
\(\operatorname{Tor}H^3(Z_\varphi;\mathbb Z)
=(\mathbb Z/12)^7\). The universal coefficient theorem
identifies this with \(\operatorname{Tor}H_2(Z_\varphi;\mathbb Z)\).
Finally \ref{maps:src:E.2.3} gives \(b_2=1\), proving the assertion.
\end{enumerate}
\end{proof}


\begin{theorem}[Full integral homology of the shifted product]\hypertarget{tag-02MQ}{}\label{maps:src:E.2.6}
\StatementRecord{maps:src:E.2.6}{proved}{02MQ}{\SeriesRef{maps:src:E.2.3}, \eqref{maps:eq:E.MV1}, \SeriesRef{maps:src:E.2.5}, \eqref{maps:eq:E.MV3}, \eqref{maps:eq:E.MV2}, \eqref{maps:eq:E.MV5}, \eqref{maps:eq:E.MV6}, \eqref{maps:eq:E.MV7}, \eqref{maps:eq:E.MV8}}


For the transverse shifted product $P$ of \ref{maps:src:E.2.3}, write
\[
 A=(\mathbb Z/12)^7,\qquad
 B=(\mathbb Z/12)^{10}\oplus(\mathbb Z/2)^7,\qquad
 C=(\mathbb Z/12)^6\oplus\mathbb Z/24\oplus(\mathbb Z/2)^{17}.
\]
Its complete integral homology is
\[
 \begin{array}{c|l}
 i&H_i(P;\mathbb Z)\\\hline
 0,1,8,9,10&\mathbb Z\\
 2&\mathbb Z\oplus A\\
 3&\mathbb Z^{20}\oplus B\\
 4&\mathbb Z^{63}\oplus C\\
 5&\mathbb Z^{88}\oplus C\\
 6&\mathbb Z^{63}\oplus B\\
 7&\mathbb Z^{20}\oplus A.
 \end{array}
\]
In particular the higher integral filtration extensions are nontrivial:
the middle torsion includes a cyclic factor of order $24$.
\end{theorem}


\begin{proof}
\begin{enumerate}[label={$\langle1\rangle$\arabic*.},leftmargin=*]
\item \emph{A decomposition with free groups on every piece.}
Choose a separating circle in $B=\mathbb P^1$ so that one closed
disc contains the three critical values of $X$ and the other
contains the three critical values of $X'$. The smooth torus
family of the other factor is trivial over each disc. Set
\[
 F=T^4,\qquad E=X\setminus\operatorname{int}(F\times D^2),
 \qquad E'=X'\setminus\operatorname{int}(F'\times D^2).
\]
The product trivializations give the decomposition
\[
 P=(E\times F')\ \cup_{F\times F'\times S^1}\ (F\times E').
 \tag{E.MV1}\hypertarget{tag-02BA}{}\label{maps:eq:E.MV1}
\]
The boundary identification preserves the two marked torus factors
and reverses the normal circle, because the two base discs have
opposite boundary orientations. The same regular-disc trivializations define both the
boundary identifications of $E,E'$ and \eqref{maps:eq:E.MV1}.

\item Since $X,X'$ are six-spheres, the Mayer--Vietoris sequence of
$S^6=E\cup(F\times D^2)$ gives, integrally,
\[
 H_0(E)=\mathbb Z,\qquad H_k(E)=H_{k-1}(F)\quad(1\leq k\leq4),
 \qquad H_k(E)=0\quad(k\geq5).
 \tag{E.MV2}\hypertarget{tag-02BB}{}\label{maps:eq:E.MV2}
\]
The isomorphism is the normal-circle sweep
$\alpha\mapsto[S^1]\times\alpha$. Indeed, in degrees one
through four the boundary map to $H_k(E)\oplus H_k(F)$ is
an isomorphism, and its second component is the identity on
the unswept fibre summand and zero on the swept summand.
In degree five the boundary fundamental class is the image
of the sphere fundamental class, so $H_5(E)=0$.
Thus all homology groups of the two pieces and their common
boundary in \eqref{maps:eq:E.MV1} are free abelian. The same statements hold
for $E'$.

\item \emph{Determining every boundary map integrally.}
Let $V=H^1(F;\mathbb Z)$, in the basis $(\gamma,u,w,\delta)$,
and use the positively ordered volume
$v=\gamma u w\delta$. The primitive global invariant classes
of the original torus fibration are
\[
 I_0=1,\quad I_1=\gamma,\quad
 I_2=q=uw+6\gamma\delta,\quad
 I_3=\gamma uw,\quad I_4=v.
\]
The images of $H^k(E)\to H^k(F)$ are exactly
\[
 n_k\mathbb Z I_k,\qquad
 (n_1,n_2,n_3,n_4)=(12,2,2,12).
 \tag{E.MV3}\hypertarget{tag-02BC}{}\label{maps:eq:E.MV3}
\]
We verify this integrally. On the disc the only local restrictions on a global
invariant class are the finite-cover pullback lattices;
the cusp specialization is an integral isomorphism onto its
invariant lattice by \ref{maps:src:E.2.5}. In degree one the finite images
contain $3\gamma$ and $4\gamma$, giving the intersection
$12\mathbb Z\gamma$. In degree two the order-three image
contains $q$, while the order-four image requires an even
coefficient of $q$, by the finite integral lattices computed in the proof of \cref{maps:src:E.1.5}.
In degree three use the perfect integral pairing downstairs:
an invariant class is a pullback precisely when its pairings
with every pulled-back degree-one class are divisible by the
covering degree. The other degree-one generators are
$2u+w+3\delta$ and $u+w+2\delta$. Their pairings with
$t\gamma uw$ are, up to signs, $3t$ and $2t$.
Thus the first cover imposes no further condition, and the
second requires $t$ even. In degree four the pullback degrees
are three and four, giving $12\mathbb Z v$.

\item These local restrictions are also sufficient for the image
from $E$. For a local system on a disc with punctures, the
restriction from its free-group $H^1$ to the direct sum of
the local coinvariants is onto: each puncture has its own
independent free generator. Consequently
$H^2(D,j_*\mathcal L)=0$. Replacing $j_*\mathcal L$ by the
direct-image sheaf changes it only by point-supported
cokernels, and so does not change this vanishing.
The Leray sequence over the disc therefore makes
$H^k(E)\to H^0(D,R^kf_*\mathbb Z)$ onto. This proves
\eqref{maps:eq:E.MV3} 

\item Write $\theta$ for the normal-circle covector. The dual of
\eqref{maps:eq:E.MV2} identifies $H^k(E)$ with $\bigwedge^{k-1}V$, and
its boundary restriction has the graph form
\[
 r_k(\alpha)=\phi_k(\alpha)+\theta\alpha,
 \qquad
 \phi_k:\bigwedge^{k-1}V\longrightarrow\bigwedge^kV.
 \tag{E.MV4}\hypertarget{tag-02BD}{}\label{maps:eq:E.MV4}
\]
The image of $\phi_k$ is exactly $n_k\mathbb Z I_k$ by
\eqref{maps:eq:E.MV3}. If $k+l=5$, classes extending to $E$ have zero
product evaluated on $\partial E$. Therefore
\[
 (-1)^k\phi_k(\alpha)\beta+
                    \alpha\phi_l(\beta)=0
 \quad\text{in }\bigwedge^4V.
 \tag{E.MV5}\hypertarget{tag-02BE}{}\label{maps:eq:E.MV5}
\]
Each $I_k$ is primitive, and wedge pairing between the two
complementary exterior powers is unimodular. Writing
$\phi_k=f_k I_k$, equation \eqref{maps:eq:E.MV5}, tested against a class
pairing to one with $I_k$, proves that $f_k$ is an integral
multiple of $\alpha\mapsto\langle\alpha I_{5-k},[F]\rangle$.
The exact image \eqref{maps:eq:E.MV3} fixes the absolute multiple. Thus
\[
 \begin{split}
 \phi_k(\alpha)&=c_k
       \langle\alpha I_{5-k},[F]\rangle I_k,\\
 (c_1,c_2,c_3,c_4)&=(12\varepsilon,2\eta,-2\eta,-12\varepsilon),
 \qquad \varepsilon,\eta\in\{1,-1\}.
 \end{split}
 \tag{E.MV6}\hypertarget{tag-02BF}{}\label{maps:eq:E.MV6}
\]
This proves every boundary map up to two signs. We allow those
signs independently for $E$ and $E'$ below; the answer is the
same for all sixteen possibilities.

\item \emph{The full integral Mayer--Vietoris matrices.}
Let $i_p:H_p(F)\to H_p(E)$ be the transpose of $\phi_p$ in
the sweep bases, with $i_0=1$; let $j_{p+1}$ denote the sweep
isomorphism for $0\leq p\leq3$, and zero for $p=4$.
The analogous maps of the second copy are $i'_q,j'_{q+1}$.
For a fibre tensor $\alpha\otimes\beta$, with degrees $p,q$,
the matrix $D_k$ of \eqref{maps:eq:E.MV1} is
\[
 \alpha\otimes\beta\longmapsto
  (i_p\alpha\otimes\beta,-\alpha\otimes i'_q\beta),
 \qquad p+q=k.
 \tag{E.MV7}\hypertarget{tag-02BG}{}\label{maps:eq:E.MV7}
\]
For a swept tensor, place the normal circle first. Then
\[
 \theta\times\alpha\times\beta\longmapsto
 \bigl(j_{p+1}\alpha\otimes\beta,
                 (-1)^p\alpha\otimes j'_{q+1}\beta\bigr),
 \qquad p+q=k-1.
 \tag{E.MV8}\hypertarget{tag-02BH}{}\label{maps:eq:E.MV8}
\]
The second sign contains the Koszul sign from moving the circle
past $\alpha$, the reversal of the two boundary-circle
orientations, and the second minus in Mayer--Vietoris. The last
two cancel. This specifies every integer entry of $D_k$.

\item In the exterior bases, exact Smith reduction of the four
adjoint maps \eqref{maps:eq:E.MV6} and the full matrices
\eqref{maps:eq:E.MV7}--\eqref{maps:eq:E.MV8} gives the following table. The source ranks are $\binom9k$; a zero binomial is
understood outside $0\leq k\leq9$.
\[
 \begin{array}{c|rrrr|l}
 k&\text{target rank}&\text{source rank}&\operatorname{rank}D_k&
       \text{unit factors}&\text{nonunit factors}\\\hline
 0&2&1&1&1&--\\
 1&10&9&9&9&--\\
 2&28&36&27&20&12^{[7]}\\
 3&64&84&53&36&2^{[7]},12^{[10]}\\
 4&114&126&82&58&2^{[17]},12^{[6]},24\\
 5&138&126&94&70&2^{[17]},12^{[6]},24\\
 6&104&84&73&56&2^{[7]},12^{[10]}\\
 7&44&36&35&28&12^{[7]}\\
 8&8&9&8&8&--\\
 9&0&1&0&0&--\\
 10&0&0&0&0&--
 \end{array}
\]
The notation $a^{[r]}$ denotes $r$ entries equal to $a$.
All sixteen choices of the two signs in each copy give precisely
this table.

\item Finally the exact homology sequence gives
\[
 0\longrightarrow\operatorname{coker}D_k
 \longrightarrow H_k(P)
 \longrightarrow\ker D_{k-1}\longrightarrow0.
\]
The last group is free, since the boundary homology is free.
This sequence therefore splits as abelian groups, so the Smith
factors in the table are the torsion of $H_k(P)$.
Adding the free kernel ranks gives exactly the rational Betti
vector in \ref{maps:src:E.2.3} and the integral groups asserted here.
The resulting torsion also satisfies the required ten-dimensional
duality $\operatorname{Tor}H_k(P)\cong
\operatorname{Tor}H_{9-k}(P)$.
\end{enumerate}
\end{proof}

\begin{proposition}[A unique transverse family type and its stable tangent bundle]\hypertarget{tag-02MR}{}\label{maps:src:E.2.7}
\StatementRecord{maps:src:E.2.7}{proved}{02MR}{\SeriesRef{maps:src:E.2.6}}


For fixed marked sphere fibrations, all M\"obius transformations
whose critical triples are disjoint give orientation-preservingly
diffeomorphic shifted products. The same conclusion holds while
either sphere varies in its proper smooth marked period family.
For every such fivefold $P$ there is an isomorphism of real vector
bundles
\[
 TP\oplus\mathbb R^4\simeq\mathbb R^{14}.
 \tag{E.MV9}\hypertarget{tag-02BI}{}\label{maps:eq:E.MV9}
\]
Thus all its positive-degree Stiefel--Whitney and Pontryagin
classes vanish. For its given complex tangent bundle,
$c_1(P)=c_2(P)=0$ integrally.
\end{proposition}


\begin{proof}
\begin{enumerate}[label={$\langle1\rangle$\arabic*.},leftmargin=*]
\item The allowed transformations form the complement in
$\operatorname{PGL}_2(\mathbb C)$ of the nine analytic
hypersurfaces imposing an equality of one critical value from
each triple. This complement is path connected: a piecewise
smooth path between two points can be perturbed relative to
its endpoints to avoid these sets, which have real codimension
two. Equivalently, use local complex coordinate balls and
perturb each segment there. Choose a smooth such path.

\item The corresponding fibre products form a compact smooth family
over its parameter interval. Indeed the vertical derivative of
$f_1(x)-f_2(y)$ is surjective at every point, since its two
factors cannot be critical at the same base value. The same
surjectivity solves the equation for a lift of the interval
coordinate vector field. A partition of unity supplies a
global smooth lift; compactness makes its flow exist for the
whole interval. This flow gives the claimed diffeomorphisms.
The complex orientations vary continuously and hence are
preserved. The identical argument applies to smooth proper
period families of either factor, keeping the critical triples
disjoint throughout.

\item View $P$ as the transverse inverse image of the diagonal in
$B\times B$ under $f_1\times f_2$. Its real oriented normal
two-plane bundle is the underlying bundle of $g^*T_B$,
where $g:P\to B$ is the common base map. It is trivial:
its first Chern class is
\[
 c_1(g^*T_B)=2g^*[B]^\vee
    =2\operatorname{pr}_1^*f_1^*[B]^\vee=0,
\]
because $H^2(X;\mathbb Z)=0$, and complex line bundles are
classified by their first Chern class. The tangent bundles
of the two six-spheres become trivial after adding one real
line each. Restriction of their product tangent bundle gives
\[
 (TP\oplus\mathbb R^2)\oplus\mathbb R^2
       \simeq\mathbb R^{14},
\]
proving \eqref{maps:eq:E.MV9} and the characteristic-class assertions for
the real tangent bundle. Finally the holomorphic tangent
sequence is
\[
 0\longrightarrow TP\longrightarrow
 \operatorname{pr}_1^*TX\oplus\operatorname{pr}_2^*TX'
 \longrightarrow g^*T_B\longrightarrow0.
\]
Both sphere factors have integral $c_1=c_2=0$, since their
cohomology in degrees two and four vanishes. The last line
has zero first Chern class, as just proved. Multiplying total
Chern classes in this exact sequence gives
$c_1(P)=c_2(P)=0$ integrally.
\end{enumerate}
\end{proof}

\begin{proposition}[The shifted product bounds with its natural stable framing]\hypertarget{tag-02MS}{}\label{maps:src:E.2.8}
\StatementRecord{maps:src:E.2.8}{proved}{02MS}{\eqref{maps:eq:E.MV9}, \SeriesRef{maps:src:B.2.1}, \eqref{maps:eq:E.B2}, \SeriesRef{maps:src:E.2.7}, \SeriesRef{maps:src:E.2.9}; \cite[Chapter 3, Section 2, Theorems 2.1 and 2.5]{rootHirsch76}}


Equip the two oriented six-spheres with their standard bounding
stable tangent framings. Trivialize the oriented two-plane bundle
$f_1^*T\mathbb P^1$ on the first sphere, and equip the transverse
shifted product $P$ with the induced stable framing in \eqref{maps:eq:E.MV9}.
Then there is a compact smooth oriented eleven-manifold $V$ with
\[
 \partial V=P,\qquad TV\oplus\mathbb R^3\simeq\mathbb R^{14},
 \qquad [P]=0\quad\hbox{in }\Omega^{\mathrm{fr}}_{10}.
 \tag{E.B1}\hypertarget{tag-029X}{}\label{maps:eq:E.B1}
\]
The boundary framing is the stated one. In particular its induced
spin structure bounds. 
\end{proposition}


\begin{proof}
\begin{enumerate}[label={$\langle1\rangle$\arabic*.},leftmargin=*]
\item Write $f_2:S^6\to\mathbb P^1$ for the shifted map of the second
factor. By \SeriesRef{maps:src:B.2.1} the first map $f_1:S^6\to\mathbb P^1$ is
null-homotopic. It therefore extends to a smooth map
$F:D^7\to\mathbb P^1$, which can be chosen to equal $f_1$ on
each sphere of a collar of the boundary. To see this with the
boundary fixed, first use the null-homotopy to obtain a continuous
extension constant in the collar direction. Smooth its interior
in coordinate charts, with cutoffs supported away from the collar;
on the overlap the extension is already smooth.

\item Let $C_2$ be the three critical values of $f_2$. The restriction
$f_1$ is transverse to each point of $C_2$, since the two critical
triples are disjoint. Perturb $F$ only in the interior to be
transverse to all these points. The applicable finite simultaneous
transversality statement is
\cite[Chapter 3, Section 2, Theorems 2.1 and 2.5]{rootHirsch76}.
For the relative assertion, the proof of that theorem uses local
perturbations multiplied by cutoff functions. Their supports may
be restricted to the complement of the collar, where
transversality has not already been secured; the open condition
on the collar persists. Compactness permits finitely many such
perturbations. 

\item Define the smooth fibre product
\[
 V=\{(x,y)\in D^7\times S^6:F(x)=f_2(y)\}.
 \tag{E.B2}\hypertarget{tag-029Y}{}\label{maps:eq:E.B2}
\]
At a point whose common value is outside $C_2$, the differential
$df_2$ is onto. At a point over $C_2$, the differential $dF$ is
onto. Consequently $dF-df_2$ is onto everywhere on \eqref{maps:eq:E.B2}.
The same argument holds on its boundary, using $df_1$ there.
The implicit-function theorem, including its boundary version,
makes $V$ a smooth eleven-manifold with boundary exactly $P$.
It is compact as a closed subset of $D^7\times S^6$.
The inverse-image orientation using the complex orientation of
$T\mathbb P^1$ induces the complex orientation of $P$ at the
boundary, because the collar keeps the first map fixed.

\item Let $g:V\to\mathbb P^1$ be the common map. The real normal
two-plane bundle of \eqref{maps:eq:E.B2} is $g^*T\mathbb P^1$, identified by
$dF-df_2$. This is the pullback from $D^7$ of $F^*T\mathbb P^1$,
which is trivial because the disc is contractible. A choice of
its oriented framing restricts on the boundary to the chosen
framing of $f_1^*T\mathbb P^1$ up to homotopy: two oriented
framings of this two-plane bundle differ by a map
$S^6\to SO(2)$, and $H^1(S^6;\mathbb Z)=0$.
The tangent exact sequence and the bounding stable framing of
the second sphere now give
\[
 \begin{split}
 TV\oplus\mathbb R^2&\simeq
       T(D^7\times S^6)|_V,\\
 TV\oplus\mathbb R^3&\simeq
       \mathbb R^7\oplus(TS^6\oplus\mathbb R)|_V
       \simeq\mathbb R^{14}.
 \end{split}
\]
At the boundary, the outward normal of $D^7$ supplies the
additional real line in $TV|_P=TP\oplus\mathbb R$.
The resulting framing is precisely the stable tangent framing
obtained from the two sphere framings and the normal framing
in \ref{maps:src:E.2.7}. This proves \eqref{maps:eq:E.B1} by the definition of framed bordism.
It also proves the assertion about the induced spin structure,
by forgetting the framing.

\end{enumerate}
\end{proof}


\begin{theorem}[Framed null-bordism over the fundamental group]\hypertarget{tag-02MT}{}\label{maps:src:E.2.9}
\StatementRecord{maps:src:E.2.9}{proved}{02MT}{\eqref{maps:eq:E.16}, \SeriesRef{maps:src:E.2.8}, \SeriesRef{maps:src:B.2.1}, \SeriesRef{maps:src:E.2.2}, \SeriesRef{maps:src:E.2.3}, \SeriesRef{maps:src:B.1.2}, \eqref{maps:eq:E.B7}; \cite[Theorems 2.3--2.4 and the Euler-number discussion in
Section 2]{rootScott83}; \cite[Definition 8.16 and Theorem 8.17]{rootDavisKirk01}; \cite[Corollary~1.3 and Table~1, pp.~112--113]{IsaksenWangXu2023}}


Let $G$ be the group in \eqref{maps:eq:E.16}. There is a closed smooth
aspherical four-manifold $N$ with
\[
 \pi_1(N)=G,\qquad \widetilde N\simeq\mathbb H^2\times\mathbb R^2,
 \qquad H_*(N;\mathbb Z)\cong H_*(S^1\times S^3;\mathbb Z),
 \tag{E.B3}\hypertarget{tag-029Z}{}\label{maps:eq:E.B3}
\]
and a smooth classifying map $c:P\to N$ inducing the displayed
identification of fundamental groups. With the natural stable
framing of \ref{maps:src:E.2.8}, one has
\[
                 [P,c]=0\quad\hbox{in }\Omega^{\mathrm{fr}}_{10}(N).
 \tag{E.B4}\hypertarget{tag-02A0}{}\label{maps:eq:E.B4}
\]
Thus $P$ has a compact stably framed eleven-dimensional filling
over $BG$. This is stronger than the filling of \ref{maps:src:E.2.8}.
\end{theorem}


\begin{proof}
\begin{enumerate}[label={$\langle1\rangle$\arabic*.},leftmargin=*]
\item \emph{The aspherical target.}
The quotient maps of \SeriesRef{mapsI:src:A.2.1} are smooth maps
$e_i:S^6\to S^3$ over the base, and \SeriesRef{maps:src:B.2.1} gives the exact
factorizations $f_i=h_i e_i$. Here $h_i:S^3\to\mathbb P^1$
is the weighted Hopf map of orders $3,4$, with the second
map shifted by the chosen M\"obius transformation. Define
\[
 N=\{(u,v)\in S^3\times S^3:h_1(u)=h_2(v)\},\qquad
 c(x,y)=(e_1(x),e_2(y)).
 \tag{E.B5}\hypertarget{tag-02A1}{}\label{maps:eq:E.B5}
\]
The critical values of each $h_i$ are only its two exceptional
circle values. They are disjoint, so the same transversality
argument as in \SeriesRef{maps:src:E.2.2} makes $N$ a compact smooth four-manifold.
The two circle actions on the weighted Hopf spheres give a
principal $T^2$ bundle in the orbifold sense over
\[
                  \mathcal B=S^2(3,4,3,4).
\]
At a cone of order $m$, its uniformizing chart is
$(D^2\times T^2)/C_m$, with the generator rotating the disc
and translating one torus coordinate by a primitive fraction
of denominator $m$. This local description is smooth and is
the same finite-filling description as \SeriesRef{mapsI:src:A.1.3}, with the unused circle kept.

\item The orbifold has Euler characteristic
$2-2(1-1/3)-2(1-1/4)=-5/6$. It is good and has hyperbolic
orbifold universal cover, by
\cite[Theorems 2.3--2.4 and the Euler-number discussion in
Section 2]{rootScott83}. Pulling the principal torus bundle to
$\mathbb H^2$ gives an ordinary principal $T^2$ bundle over a
contractible base, hence the trivial bundle. Its universal
cover is $\mathbb H^2\times\mathbb R^2$. The local chart above
shows that these orbifold pullbacks are ordinary covering maps
of total spaces. This proves asphericity and injectivity of
the fibre subgroup $\mathbb Z^2$.

\item Using the marked fibre generators $h_1,h_2$ and the four
positive peripheral generators gives exactly
\[
 \begin{split}
 \pi_1(N)=\langle a_1,a_2,b_1,b_2,h_1,h_2\mid {}&h_1,h_2
       \text{ central},\quad a_1^3=h_1,\quad a_2^4=h_1^{-1},\\
 &b_1^3=h_2,\quad b_2^4=h_2^{-1},\quad
 a_1a_2b_1b_2=1\rangle.
 \end{split}
\]
The map $c$ sends the surviving regular-fibre generators
$\widehat\gamma,\widehat\gamma'$ to $h_1,h_2$, respectively,
and each peripheral generator to the one with the same name.
This is because $e_i$ is the literal circle quotient on these
fibres and finite collars. Its cusp replacement extends over
a solid torus and kills precisely the cusp meridian in this
presentation. Comparing with the full van Kampen calculation
in \ref{maps:src:E.2.3} therefore shows that $c_*$ is the identity on the
displayed presentation, and in particular is an isomorphism.
As $N$ is aspherical, $c$ is a classifying map.

\item For completeness the abelian relation matrix has Smith
diagonal $(1,1,1,1,1)$; a primitive integer kernel vector in
the order $(a_1,a_2,b_1,b_2,h_1,h_2)$ is
$(4,-3,-4,3,12,-12)$. Thus $H_1(N)=\mathbb Z$.
Every stratum of its base has a torus fibre, so Euler
additivity gives $\chi(N)=0$. Integral Poincar\'e duality and
the universal coefficient theorem give $H_3(N)=\mathbb Z$
and no torsion in $H_2(N)$. Euler characteristic then forces
$H_2(N)=0$, proving the integral homology assertion in \eqref{maps:eq:E.B3}.

\item \emph{The bordism computation before taking the inverse image.}
For a based sphere, the stable splitting
$(S^3)_+\simeq S^0\vee S^3$ identifies
\[
 \Omega^{\mathrm{fr}}_6(S^3)
       =\pi^{\mathrm S}_6\oplus\pi^{\mathrm S}_3.
 \tag{E.B6}\hypertarget{tag-02A2}{}\label{maps:eq:E.B6}
\]
We use Pontryagin--Thom in its form for maps to a space,
\cite[Definition 8.16 and Theorem 8.17]{rootDavisKirk01}.
The first component of $[S^6,e_i]$ is obtained by forgetting
the map, so it is zero for the bounding sphere framing.
The second component is the stable class of $e_i$, equivalently
the framed inverse image of a regular point. It is $\eta^3$
by \SeriesRef{maps:src:B.1.2}. One can also check the splitting directly: the
Pontryagin--Thom collapse of the framed sphere into its
disjoint-basepoint suspension has degree one on the sphere
summand; composing with $e_i$ gives its stable homotopy class.
Thus
\[
                         [S^6,e_i]=(0,\eta^3).
\]
Smashing the two sphere splittings gives
\[
 (S^3\times S^3)_+\simeq
       S^0\vee S^3\vee S^3\vee S^6.
\]
Under this splitting the external product
$[S^6\times S^6,e_1\times e_2]$ has four components.
The first three each contain a factor $[S^6]=0$; the last
is $\eta^3\eta^3=\eta^6=0$. Indeed $\eta^4$ belongs to
$\pi^{\mathrm S}_4=0$, as recorded in
\cite[Corollary~1.3 and Table~1, pp.~112--113]{IsaksenWangXu2023}. Therefore
\[
 [S^6\times S^6,e_1\times e_2]=0
      \quad\hbox{in }\Omega^{\mathrm{fr}}_{12}(S^3\times S^3).
 \tag{E.B7}\hypertarget{tag-02A3}{}\label{maps:eq:E.B7}
\]


\item \emph{The comparison with the classifying map.}
By \eqref{maps:eq:E.B7} there is a compact stably framed thirteen-manifold
$Z$ with boundary $S^6\times S^6$ and a map
$E:Z\to S^3\times S^3$ extending $e_1\times e_2$.
Smooth it relative to the boundary and make it transverse
to $N$ there. Transversality on the boundary already holds:
the derivative of its defining equation is $df_1-df_2$,
which is onto. Relative transversality therefore supplies
the interior perturbation without changing that boundary.

\item The normal bundle of $N\subset S^3\times S^3$ is the
restriction of $h_1^*T\mathbb P^1$ pulled back from the first
sphere. This oriented two-plane bundle is trivial because
$H^2(S^3;\mathbb Z)=0$. Choose a framing on it before
restriction to $N$. Its pullback under $e_1\times e_2$
is a framing of $f_1^*T\mathbb P^1$, agreeing with that
in \ref{maps:src:E.2.8} up to homotopy since $H^1(S^6;\mathbb Z)=0$.
It follows that
\[
                  V=E^{-1}(N),\qquad E|_V:V\longrightarrow N
\]
is a compact stably framed eleven-manifold with boundary
$(P,c)$ and the required boundary framing. This proves \eqref{maps:eq:E.B4}.

\end{enumerate}
\end{proof}


\begin{definition}[Ordinary closed spin]\label{def:V:ordinary-closed-spin}
For a closed
submanifold $L\subset\operatorname{int}D^n$, its ordinary closed
spin is $L\times S^1$ in
$D^n\times S^1\subset\partial(D^n\times D^2)=S^{n+1}$.
\end{definition}

\begin{proposition}[The fibre complement and a spinning obstruction]\hypertarget{tag-02MA}{}\label{maps:src:E.2.10}
\StatementRecord{maps:src:E.2.10}{proved}{02MA}{\SeriesRef{maps:src:E.2.6}, \SeriesRef{maps:src:B.5.2}, \eqref{maps:eq:E.MV4}, \eqref{maps:eq:E.MV6}, \eqref{maps:eq:E.S5}, \eqref{maps:eq:E.S4}; \cite[Section 3, immediately before Theorem 3.4]{mapsWinter08}; \cite[Theorem 4.5]{mapsWinter08}}


Let $F=T^4$ be a regular fibre of the original sphere fibration and
$E=S^6\setminus\operatorname{int}(F\times D^2)$, with the marking of
\ref{maps:src:E.2.6}. Its group and peripheral homomorphism are
\[
 \begin{split}
 K=\pi_1(E)&=\langle a,b,h\mid h\text{ central},\ a^3=h,
                                      \ b^4=h^{-1}\rangle,\\
 g&\longmapsto h,\qquad u,w,d\longmapsto1,
       \qquad \mu\longmapsto ab.
 \end{split}
 \tag{E.S1}\hypertarget{tag-02CC}{}\label{maps:eq:E.S1}
\]
Here $\mu$ is the normal meridian, with the peripheral orientation
chosen so that $[a]=4[\mu]$, $[b]=-3[\mu]$ and $[h]=12[\mu]$.
The image of the boundary group is the embedded copy
$\langle h,\mu\rangle\simeq\mathbb Z^2$. In the unique homotopy
class of normal framings for which every pushed-off loop of $F$
has linking number zero, the image of $g$ is
\[
             \lambda=h\mu^{-12},
 \qquad \ker\{\pi_1(F)\longrightarrow K\}
                 =\mathbb Z\langle u,w,d\rangle.
 \tag{E.S2}\hypertarget{tag-02CN}{}\label{maps:eq:E.S2}
\]
In particular the peripheral structure is that of the $(3,4)$
torus knot, with its regular Seifert fibre rather than its preferred
longitude as the original $g$ coordinate. Nevertheless the embedded
four-torus $F\subset S^6$ is not the ordinary threefold closed spin
of that classical knot, even allowing ambient orientation reversal
and arbitrary changes of its normal framing. More strongly, it is
not a single ordinary closed spin of any embedded $T^3\subset S^5$.

\end{proposition}


\begin{proof}
\begin{enumerate}[label={$\langle1\rangle$\arabic*.},leftmargin=*]
\item Apply the van Kampen presentation in \SeriesRef{maps:src:B.5.2} with the regular base disc
removed. The cusp kills $w,d$ and its peripheral loop. Conjugating
$w$ by the order-three generator also kills $u$. The two finite caps
impose $a^3=h$ and $b^4=h^{-1}$; there is now no relation $ab=1$,
since the outer boundary has not been filled. The outer meridian
is $ab$ in the stated orientation. This proves \eqref{maps:eq:E.S1}, including the peripheral map. Eliminating $h$ gives $\langle a,b\mid a^3b^4=1\rangle$.
Its quotient by $h$ is $C_3*C_4$. The word $ab$ has infinite order
in that free product, so a relation $h^r\mu^s=1$ forces $s=0$.
Abelianization then forces $r=0$. This proves peripheral injectivity.
It also shows that $\lambda$ has infinite order: its image in
$C_3*C_4$ is $(ab)^{-12}$.

\item Write $\theta\in H^1(\partial E;\mathbb Z)$ for the normal-circle
covector and use $(\gamma,u,w,\delta)$ for $H^1(F;\mathbb Z)$.
The generator $\xi\in H^1(E;\mathbb Z)$ restricts to
$\theta+12\gamma$. Homotopy classes of normal framings differ by
$H^1(F;\mathbb Z)$. Consequently there is a unique framing for
which the associated section $s:F\to\partial E$ satisfies
$s^*\xi=0$: it has $s^*\theta=-12\gamma$. Its induced peripheral
map is exactly \eqref{maps:eq:E.S2}. Simultaneously reversing the meridian
conventions changes the displayed signs and preserves all the
following assertions.

\item For this section set
\[
 \mathcal R=\operatorname{im}\{H^2(E;\mathbb Z)
                         \longrightarrow H^2(F;\mathbb Z)\},
 \qquad
 \operatorname{rad}(\mathcal R)=
 \{v\in H_1(F;\mathbb Z):\iota_v\omega=0
                                    \text{ for all }\omega\in\mathcal R\}.
 \tag{E.S3}\hypertarget{tag-02CY}{}\label{maps:eq:E.S3}
\]
This uses the canonical identification of the cohomology of a torus
with the exterior algebra of its first cohomology. Equations
\eqref{maps:eq:E.MV4}--\eqref{maps:eq:E.MV6} give classes whose boundary restrictions are
$\theta u$ and $\theta w$: their unswept terms are zero because
$u\gamma uw=w\gamma uw=0$. Restricting them along $s$ proves
\[
                -12\gamma u,\ -12\gamma w\ \in\mathcal R.
 \tag{E.S4}\hypertarget{tag-02D9}{}\label{maps:eq:E.S4}
\]
The simultaneous contraction kernel of these two forms is
$\mathbb Z d$. Indeed their contractions force successively
$\gamma(v)=u(v)=w(v)=0$. Thus
\[
       \operatorname{rad}(\mathcal R)\subseteq\ker\gamma
        =\ker\{\pi_1(F)\longrightarrow K\}.
 \tag{E.S5}\hypertarget{tag-02DK}{}\label{maps:eq:E.S5}
\]
Both the zero-linking framing and this inclusion are invariants of
the embedded pair, independently of the torus coordinates.

\item Use the ordinary closed spin of Definition~\ref{def:V:ordinary-closed-spin}.
Start with the closed classical knot $K_0\subset\operatorname{int}D^3$
and repeat this operation three times, each time placing the
previous compact pair in the interior of a disc. This is the
ordinary torus spinning construction reviewed in
\cite[Section 3, immediately before Theorem 3.4]{mapsWinter08}.
If $\Sigma$ is a Seifert surface for $K_0$, the same three spins
give the embedded Seifert hypersurface
$\Sigma\times T^3$ with boundary $K_0\times T^3$.
Pushing the latter off using the normal of this hypersurface
gives its zero-linking framing.

\item Let $v$ be the original knot-circle class and let $v_1,v_2,v_3$
be the spinning circles. Each pushed-off two-cycle $v\wedge v_i$
bounds the corresponding copy of $\Sigma\times S^1$ in a parallel
copy of that hypersurface, disjoint from the knot. Hence every
restriction of a complement two-class vanishes on $v\wedge H_1(F)$:
$v$ belongs to $\operatorname{rad}(\mathcal R)$.
On the other hand its pushed-off loop is the original preferred
longitude and is nontrivial for $K_0=T(3,4)$, whereas the spinning
circles are trivial in the complement group. These peripheral
claims also follow directly at each spin: its complement is
\[
 ((D^n\setminus\operatorname{int}\nu L)\times S^1)
       \ \cup_{S^{n-1}\times S^1}\ (S^{n-1}\times D^2),
\]
so van Kampen kills precisely the new spinning circle and preserves
the old complement group. They agree with the primary peripheral
calculation \cite[Theorem 4.5]{mapsWinter08} in the first spin.
Thus the ordinary spun torus violates \eqref{maps:eq:E.S5}. An identification
of embedded pairs would preserve its two sides, giving a
contradiction.

\item For the stronger assertion, let $L=T^3\subset S^5$ be arbitrary
and suppose that its ordinary closed spin were our four-torus.
Write $E_L$ for its exterior. Deleting a ball disjoint from
$L$ changes cohomology only in degrees at least four. In the
displayed spun-complement decomposition, with $n=5$, the cap
and intersection are $S^4\times D^2$ and $S^4\times S^1$.
Mayer--Vietoris and the surjectivity onto the latter first
cohomology therefore give, in degrees two and three, precisely
the cohomology of $(E_L\setminus D^5)\times S^1$.
Let $s$ be the spinning-circle covector and $v_s$ its primitive
homology generator. Put $\mathcal R_L^j=\operatorname{im}
\{H^j(E_L)\to H^j(L)\}$, using the zero-linking framing.
One has $\mathcal R_L^1=0$. Also $\mathcal R_L^3=0$:
the pushed-off $L$ bounds a parallel Seifert hypersurface in
its exterior, so every top-degree restriction evaluates to
zero. The product restrictions consequently satisfy
\[
       \mathcal R^2=\mathcal R_L^2,\qquad
       \mathcal R^3=\mathcal R_L^2\wedge s.
 \tag{E.S12}\hypertarget{tag-02CF}{}\label{maps:eq:E.S12}
\]
In particular $v_s\in\operatorname{rad}(\mathcal R^2)$,
and contraction by $v_s$ is injective on $\mathcal R^3$.
For our four-torus, \eqref{maps:eq:E.S4} forces $v_s=\pm d$.
But \eqref{maps:eq:E.MV4}--\eqref{maps:eq:E.MV6}, now applied to $\alpha=\gamma\delta$
in degree three, show that $\mathcal R^3$ contains
$\pm2\gamma uw$. Its swept term vanishes in the zero-linking
framing, and this is a nonzero class annihilated by $d$.
This contradicts that injectivity. No ordinary closed spin
is possible. Turned or deform-spun constructions insert
additional data and are not excluded by this argument.
\end{enumerate}
\end{proof}

\begin{theorem}[Exact second-homotopy coinvariants]\hypertarget{tag-02MB}{}\label{maps:src:E.2.11}
\StatementRecord{maps:src:E.2.11}{proved}{02MB}{\SeriesRef{maps:src:E.2.9}, \SeriesRef{maps:src:E.2.6}, \eqref{maps:eq:E.S7}, \eqref{maps:eq:E.MV7}, \eqref{maps:eq:E.MV8}, \eqref{maps:eq:E.MV6}, \eqref{maps:eq:E.MV4}, \eqref{maps:eq:E.S9}, \eqref{maps:eq:E.S10}, \eqref{maps:eq:E.S11}}


For the classifying map $c:P\to N=K(G,1)$ of \ref{maps:src:E.2.9}, one has
\[
 c_*:H_3(P;\mathbb Z)\longrightarrow H_3(N;\mathbb Z)=\mathbb Z
       \quad\text{equal to zero},
 \qquad
 \bigl(\pi_2(P)\bigr)_G
       \simeq\mathbb Z^2\oplus(\mathbb Z/12)^6.
 \tag{E.S6}\hypertarget{tag-02DR}{}\label{maps:eq:E.S6}
\]
The pullback of a generator of $H^3(N;\mathbb Z)$ has exact order
twelve in $H^3(P;\mathbb Z)$. In particular $\pi_2(P)$ is not a
projective $\mathbb ZG$-module. The fivefold cannot be diffeomorphic,
or homotopy equivalent, to $N\times S^6$ connected-summed with any
finite collection of $S^r\times S^{10-r}$, $2\leq r\leq8$, and
any homotopy ten-sphere. \end{theorem}


\begin{proof}
\begin{enumerate}[label={$\langle1\rangle$\arabic*.},leftmargin=*]
\item Use the separating base circle in \ref{maps:src:E.2.6} also in the target of
\ref{maps:src:E.2.9}. There $M$ denotes the exterior of a regular orbit of the
weighted Hopf map, namely the torus-knot exterior of \SeriesRef{mapsI:src:A.1.3}. The
two decompositions, and the restriction of $c$, are
\[
 \begin{array}{ccc}
 (E\times F')\cup(F\times E')&\longrightarrow&
                  (M\times S^1)\cup(S^1\times M')=N,\\
 F\times F'\times S^1&\longrightarrow&S^1\times S^1\times S^1.
 \end{array}
 \tag{E.S7}\hypertarget{tag-02DS}{}\label{maps:eq:E.S7}
\]
On the common boundary the map retains exactly the covectors
$\gamma_1,\gamma_2,\theta$. It is the literal quotient there.
Each knot exterior has the integral homology of a circle, so the
third homology and cohomology of each target piece in \eqref{maps:eq:E.S7} vanish.

\item We first prove the homological statement. In the degree-two source of \eqref{maps:eq:E.MV7}--\eqref{maps:eq:E.MV8}
write $t_1,t_2,t_{12}$ for the coefficients of
$\theta\times g_1$, $\theta\times g_2$, and $g_1\times g_2$.
The image of $i_2:H_2(F)\to H_2(E)$ lies entirely in the
$d$-sweep coordinate, by \eqref{maps:eq:E.MV6}. Consequently the $g_1$-sweep
coordinate in the first $H_2(E)$ summand of the target is just
$t_1$, and the corresponding second coordinate is just $t_2$.
Every vector in the kernel of the boundary matrix therefore has
$t_1=t_2=0$. Its first $H_1(E)\otimes\mathbb Zg_2$ coordinate
then forces $12t_{12}=0$, up to an immaterial orientation sign.
These three coordinates are the only ones retained in
$H_2(T^9)\to H_2(T^3)$. The target Mayer--Vietoris sequence
injects $H_3(N)$ into $H_2(T^3)$; its source-piece third groups
are zero. Naturality now proves $c_*H_3(P)=0$.

\item The cohomological calculation determines the extension. On the target boundary write
$v=\gamma_1\gamma_2$. The generators of the two degree-two
piece cohomology groups restrict, with consistent meridian
orientations, to
\[
       \theta\gamma_2+12v,
       \qquad \theta\gamma_1+12v.
 \tag{E.S8}\hypertarget{tag-02DT}{}\label{maps:eq:E.S8}
\]
Their image is a primitive rank-two subgroup of $H^2(T^3)$:
the minor in the last two coordinates is a unit. Its quotient
is freely generated by $v$. Since the degree-three piece groups
are zero, the connecting class $\partial v$ is a generator of
$H^3(N;\mathbb Z)$.

\item In the source decomposition both $\theta\gamma_1$ and
$\theta\gamma_2$ already belong to the image of the degree-two
piece restrictions. This is \eqref{maps:eq:E.MV4} with $\alpha=\gamma$,
since $\phi_2(\gamma)=0$; the second normal circle is reversed,
which only changes its sign. Subtracting one of these classes
from the corresponding product of a degree-one class of $E$
with $\gamma_2$ proves that $12v$ is a source boundary image.
Conversely the coefficient of $v$ in every such boundary image
is divisible by twelve. The pure degree-two terms from a single
factor have no $v$ coefficient, and the mixed terms use the
restriction $\theta+12\varepsilon\gamma$ of the degree-one
generator, with $\varepsilon=\pm1$. It follows that
\[
                  c^*(\partial v)=\partial v
        \quad\text{has exact order }12.
 \tag{E.S9}\hypertarget{tag-02DU}{}\label{maps:eq:E.S9}
\]
This uses the Mayer--Vietoris image.

\item Put $A=(\mathbb Z/12)^7$ and $L=(\pi_2(P))_G$.
The low-degree homology sequence of the universal-cover
fibration, together with \ref{maps:src:E.2.6} and \ref{maps:src:E.2.9}, gives
\[
             0\longrightarrow\mathbb Z z\longrightarrow L
                \longrightarrow\mathbb Z\oplus A\longrightarrow0.
 \tag{E.S10}\hypertarget{tag-02CD}{}\label{maps:eq:E.S10}
\]
The zero on the left uses $c_*H_3(P)=0$, just proved.
To determine this extension, replace $c$ by its mapping cylinder
and use relative groups for $(N,P)$. The induced map on fundamental
groups is an isomorphism and the target has zero second homotopy
group. On universal covers the relative pair is two-connected;
relative Hurewicz gives
$H_3(\widetilde N,\widetilde P)=\pi_2(P)$ and all lower relative
groups vanish. The relative covering homology sequence therefore
identifies $H_3(N,P)=L$ and $H_2(N,P)=0$.
The universal coefficient theorem and the relative cohomology
sequence identify the map
\[
       H^3(N,P)=\operatorname{Hom}(L,\mathbb Z)
         \longrightarrow H^3(N)=\mathbb Z
\]
with evaluation on $z$. Its image is $12\mathbb Z$ by \eqref{maps:eq:E.S9}.

\item The free summand of the quotient in \eqref{maps:eq:E.S10} splits. Lift its
seven order-twelve generators to $x_1,\ldots,x_7$ in the
remaining extension. There are integers $e_i$, defined modulo
twelve, with the complete abelian presentation
\[
        \langle z,x_1,\ldots,x_7\mid
                      12x_i=e_i z\ (1\leq i\leq7)\rangle.
 \tag{E.S11}\hypertarget{tag-02CE}{}\label{maps:eq:E.S11}
\]
Evaluation of homomorphisms to $\mathbb Z$ on $z$ has image
$\bigl(12/\gcd(12,e_1,\ldots,e_7)\bigr)\mathbb Z$.
Thus that gcd is one. The relation matrix in \eqref{maps:eq:E.S11} has
Smith factors $1,12,12,12,12,12,12$, giving
$\mathbb Z\oplus(\mathbb Z/12)^6$. Adding the split free
summand proves \eqref{maps:eq:E.S6}.

\item If $\pi_2(P)$ were projective over $\mathbb ZG$, it would be
a direct summand of a free module. Taking coinvariants preserves
that splitting, so $L$ would be a projective, hence free,
abelian group. The torsion in \eqref{maps:eq:E.S6} contradicts this. Finally
$N\times S^6$ has zero second homotopy group. Connected sum in
dimension ten with a simply connected sphere product adds its
second homotopy as a free induced $\mathbb ZG$-module, as seen
by lifting the connected-sum necks and using their simply
connected boundary spheres. Only $r=2,8$ contribute, with one
free generator each; a homotopy sphere contributes none.
Their second homotopy modules are therefore free, so none is
homotopy equivalent to $P$.

\end{enumerate}
\end{proof}

\begin{theorem}[Central localizations of the homotopy module]\hypertarget{tag-02MC}{}\label{maps:src:E.2.12}
\StatementRecord{maps:src:E.2.12}{proved}{02MC}{\SeriesRef{maps:src:E.2.9}, \SeriesRef{maps:src:E.2.10}, \eqref{maps:eq:E.S17}, \SeriesRef{maps:src:E.2.3}, \eqref{maps:eq:E.S19}, \eqref{maps:eq:E.S21}, \SeriesRef{maps:src:E.2.11}}


Take the central generators $h_1,h_2$ of $G$ from \ref{maps:src:E.2.9} and put
\[
 R=\mathbb ZG,\qquad R_i=R[(h_i-1)^{-1}],\qquad
              R_{12}=R[(h_1-1)^{-1},(h_2-1)^{-1}].
 \tag{E.S13}\hypertarget{tag-02CG}{}\label{maps:eq:E.S13}
\]
These are flat central localizations. As modules over the indicated
localized group rings, the second homotopy group satisfies
\[
 \begin{split}
 R_1\otimes_R\pi_2(P)&\simeq R_1/(h_2-1),\\
 R_2\otimes_R\pi_2(P)&\simeq R_2/(h_1-1),\\
 R_{12}\otimes_R\pi_2(P)&=0.
 \end{split}
 \tag{E.S14}\hypertarget{tag-02CH}{}\label{maps:eq:E.S14}
\]
More generally, $H_j(P;R_1)$ is zero outside $2\leq j\leq7$;
in those six degrees it is respectively a direct sum of
\[
                  1,\ 3,\ 4,\ 4,\ 3,\ 1
 \tag{E.S15}\hypertarget{tag-02CI}{}\label{maps:eq:E.S15}
\]
copies of $R_1/(h_2-1)$. The analogous statement holds with the
indices interchanged. All homology with coefficients in $R_{12}$
vanishes. For the fibre complement of \ref{maps:src:E.2.10}, with
$S=\mathbb ZK[(h-1)^{-1}]$, one likewise has
\[
             S\otimes_{\mathbb ZK}\pi_2(E)\simeq S,
       \qquad H_4(E;S)\simeq S,
 \tag{E.S16}\hypertarget{tag-02CJ}{}\label{maps:eq:E.S16}
\]
and its other localized homology groups vanish. Moreover, if
$q=(h_1-1)(h_2-1)$, then
\[
                  q^2\pi_2(P)=0,
        \qquad \dim_{\mathbb Q}\bigl(\pi_2(P)\otimes\mathbb Q\bigr)
                     =\infty.
 \tag{E.S20}\hypertarget{tag-02CO}{}\label{maps:eq:E.S20}
\]
Consequently the universal cover has no finite homotopy type.
There is no fibration of $P$ over $N$ with compact manifold fibre
and inducing an isomorphism on fundamental groups.

\end{theorem}


\begin{proof}
\begin{enumerate}[label={$\langle1\rangle$\arabic*.},leftmargin=*]
\item \emph{The cusp module with its marked deck transformation.}
In the original cusp, the angular generators $w,d$ vanish and
$\pi_1(N_0)=\mathbb Z\langle g,u\rangle$. In $E$ their images
are respectively $h,1$; in the first factor of $P$ they are
$h_1,1$. Thus the connected cover of this cusp occurring inside
the universal cover is the cover $\widehat W$ of its central
fibre in which the $g$ deck transformation is $z$, whereas $u$
remains a loop. Write $A=\mathbb Z[z,z^{-1}]$.
Its ordinary homology modules are
\[
 H_0(\widehat W)=H_1(\widehat W)=A/(z-1),\quad
 H_2(\widehat W)=A\oplus(A/(z-1))^2,
 \quad H_3(\widehat W)=0,\quad H_4(\widehat W)=A.
 \tag{E.S17}\hypertarget{tag-02CK}{}\label{maps:eq:E.S17}
\]
Here is an integral normalization calculation of this input.
The dual complex of the cover is an infinite cylinder, whose
ordinary homology is $A/(z-1)$ in degrees zero and one and zero
in degree two. The normalization degree-four row is a free
copy of $A$. In degree two, the homology map from the three
double curves to the four component classes is the transpose,
up to the unimodular intersection marking and possible inversion
of $z$, of
\[
 B(z)=\begin{pmatrix}
 -1&-1&-1&-1\\ 1&1&1&z\\-z&-z&-1&-z
 \end{pmatrix}
 \ \sim\
 \begin{pmatrix}1&0&0&0\\0&z-1&0&0\\0&0&z-1&0\end{pmatrix}.
 \tag{E.S18}\hypertarget{tag-02CL}{}\label{maps:eq:E.S18}
\]
For example multiply on the left by
$\left(\begin{smallmatrix}-1&0&0\\1&1&0\\-z&0&1\end{smallmatrix}\right)$,
then use as new columns $e_1,e_4-e_1,e_3-e_1,e_2-e_1$.
Both operations are unimodular over $A$. The transpose is
injective and its cokernel is the displayed degree-two module.
All strata have only even homology, and the normalization
direction has length two, so no higher differential is possible.
There is no degree-two extension from the dual complex, whose
second homology is zero. This proves \eqref{maps:eq:E.S17}. Ordinary chains
in this locally finite normalization have compact support and
use direct sums.

\item \emph{Localization and the vanishing of every other piece.}
The central subgroup $\langle h_1,h_2\rangle$ is an injected
$\mathbb Z^2$ by \ref{maps:src:E.2.9}. The group ring is free over its central
Laurent subring, by choosing coset representatives. In particular
$h_i-1$ and their products are non-zero-divisors. Inverting them
is central localization and is flat; it commutes with homology
of the universal-cover chain complex. The analogous assertions
hold for $h$ in $K$.

\item Use the regular torus bundle over a graph and the six closed
cap neighborhoods in \ref{maps:src:E.2.3}. For $R_1$ coefficients, the first
regular torus has its $g_1$ loop acting by $h_1$. The torus
chain complex is the Koszul complex of its commuting loop
actions. Its $h_1-1$ term is invertible, providing a chain
contraction. Thus every regular fibre has zero homology with
these coefficients. The homology sequence of a bundle over
a graph, and the same calculation over each boundary circle,
show that the whole regular piece and every cap boundary are
acyclic. Each cap of the second factor is a product with a
first-factor regular torus over a disc, so it too is acyclic.

\item At a finite cap of the first factor, the affine generator
$a_j$ has $a_j^{m_j}=h_1^{\pm1}$ in its image in $G$.
The connected cover of the reduced flat surface corresponding
to its map onto $\langle a_j\rangle$ is $T^3\times\mathbb R$.
The $m_j$-th power acts by an integral translation with linear
part $A_j^{m_j}=I$, and therefore acts identically on all of
its homology. Thus $h_1-1$ annihilates every such homology
module before localization. Keeping the other regular torus does not change that assertion, and induction to $G$ commutes
with the central action. Both first-factor finite caps are
therefore acyclic after localization as well. 

\item \emph{The one surviving cap and induction.}
The first cusp is $N_{0,1}\times F_2$, since the two critical
triples are disjoint. Its image in $G$ is precisely
$H=\langle h_1,h_2\rangle$. A connected component of its
preimage in the universal cover is
\[
                \widehat W\times\mathbb R\times T^3
\]
up to deformation retraction: $g_1$ acts on $\widehat W$ by
$z=h_1$, and $g_2=h_2$ translates the real factor. It acts
trivially on the homology of that three-torus. All
components together give induction from $\mathbb ZH$ to
$R$. By \eqref{maps:eq:E.S17}, inverting $h_1-1$ leaves one free
$\mathbb Z[h_1^{\pm1},(h_1-1)^{-1}]$ summand in cusp degrees
two and four and nothing else. The integral product homology
with $T^3$ has no torsion correction. Inducing the result gives
\[
 H_j(N_{0,1}\times F_2;R_1)
   \simeq
 \left(R_1/(h_2-1)\right)^{
           \binom3{j-2}+\binom3{j-4}},
 \tag{E.S19}\hypertarget{tag-02CM}{}\label{maps:eq:E.S19}
\]
where a binomial coefficient is zero outside its usual range.
All other pieces and their intersections are acyclic, as
proved above, so Mayer--Vietoris identifies this with
$H_j(P;R_1)$. This proves \eqref{maps:eq:E.S15} and \eqref{maps:eq:E.S19}, and interchanging
the factors proves the second version. Inverting $h_2-1$
also kills \eqref{maps:eq:E.S19}, proving the double-localization assertion.
Since $\widetilde P$ is simply connected, Hurewicz identifies
its second homology with $\pi_2(P)$; flat localization now
gives \eqref{maps:eq:E.S14}.

\item For $E$ the same decomposition has only one cusp and two
finite caps. Every other piece is acyclic after $h-1$ is
inverted by the identical argument. Its cusp image is
$\langle h\rangle$, so inducing \eqref{maps:eq:E.S17} after localization
gives one free $S$ summand in degrees two and four. This
proves \eqref{maps:eq:E.S16}.

\item \emph{A uniform annihilator without localization.}
Put $I=(h_1-1,h_2-1)\subset R$. Before localization, the homology
of a regular eight-torus with $R$ coefficients is
\[
                 (R/I)\otimes_{\mathbb Z}\bigwedge^j\mathbb Z^6
                      \quad(0\leq j\leq6).
 \tag{E.S21}\hypertarget{tag-02CP}{}\label{maps:eq:E.S21}
\]
Indeed the two nontrivial loop images are the independent central
generators $h_1,h_2$. Their differences from one form a regular
sequence in the central Laurent ring, and $R$ is free over that
ring. The two-variable Koszul complex therefore has only its
degree-zero quotient as homology; the other six loop differentials
are zero. This proves \eqref{maps:eq:E.S21}.
The homology sequence of the regular bundle over a graph has
only base degrees zero and one. Both filtration quotients are
annihilated by $I$, by \eqref{maps:eq:E.S21}, so $I^2$ annihilates the homology
of the regular piece. The same two-stage Wang argument applies
to every boundary mapping torus. 

\item At a finite cap its connected lift is a product of six torus
directions and two real directions up to deformation retraction.
Each central deck generator acts trivially on homology: the
one belonging to that cap is the finite affine power already
checked above, and the other translates the regular real
direction. Hence $I$ annihilates its induced homology. At the
first cusp, $h_2$ translates the real factor of the other
regular torus, so $h_2-1$ annihilates all of its homology;
the second cusp is annihilated by $h_1-1$. Thus $q$ annihilates
the homology of every piece and every intersection, since
$q\in I^2$. The decomposition has a regular piece and disjoint
caps with disjoint boundary collars, so its Mayer--Vietoris
sequence gives in every degree
\[
 0\longrightarrow\operatorname{coker}D_j
   \longrightarrow H_j(P;R)
   \longrightarrow\ker D_{j-1}\longrightarrow0.
\]
Both end groups are killed by $q$, and therefore the middle
group is killed by $q^2$. Since $H_2(P;R)=\pi_2(P)$, this
proves the first assertion of \eqref{maps:eq:E.S20} integrally.

\item \emph{Infinite rational rank.}
If $\pi_2(P)\otimes\mathbb Q$ were finite dimensional,
localizing its linear automorphism $h_1$ at $h_1-1$ would
remain finite dimensional. Explicitly, this removes the
generalized eigenspace with eigenvalue one and inverts an
automorphism on the remaining finite-dimensional quotient.
But \eqref{maps:eq:E.S14}, tensored with $\mathbb Q$, contains a copy of
$\mathbb Q[h_1^{\pm1},(h_1-1)^{-1}]$, as follows by taking
the identity coset in the free module over the central
Laurent subring and setting $h_2=1$. That vector space is
infinite dimensional, a contradiction. This proves the
second assertion and rules out a finite homotopy type of
the universal cover. A fibration over the aspherical $N$
inducing a fundamental-group isomorphism would have simply
connected fibre, and pulling it to $\widetilde N$ would
identify that fibre up to homotopy with $\widetilde P$.
A compact manifold fibre has finite homotopy type, so such
a fibration is impossible.

\item The quotient $R_1/(h_2-1)$ is nonzero, since the two central
generators are independent, and it is annihilated by the
non-zero-divisor $h_2-1$. A projective $R_1$-module is a
direct summand of a free module and has no such nonzero
annihilated element. Thus \eqref{maps:eq:E.S14} also gives an independent
proof of the nonprojectivity in \ref{maps:src:E.2.11}. \end{enumerate}
\end{proof}

\begin{theorem}[A circle fibration and infinite-order return]\hypertarget{tag-02MD}{}\label{maps:src:E.2.13}
\StatementRecord{maps:src:E.2.13}{proved}{02MD}{\SeriesRef{maps:src:E.2.9}, \eqref{maps:eq:E.S26}, \SeriesRef{maps:src:E.2.12}, \SeriesRef{maps:src:B.6.24}; \cite[Theorem 3.6]{rootScott83}}


The shifted fivefold admits a smooth fibre bundle
\[
                  Y^9\longrightarrow P\stackrel{\rho}{\longrightarrow}S^1
 \tag{E.S22}\hypertarget{tag-02CQ}{}\label{maps:eq:E.S22}
\]
with connected closed fibre, inducing the primitive character
\[
 \chi(a_1,a_2,b_1,b_2,h_1,h_2)
                    =(4,-3,-4,3,12,-12)
 \tag{E.S23}\hypertarget{tag-02CR}{}\label{maps:eq:E.S23}
\]
on the presentation of \ref{maps:src:E.2.9}. Its fibre group is the fundamental
group of the oriented principal circle bundle of Euler number
$-1$ over a surface of genus six. Here the base has its complex
orientation and the positive circle is $h_1h_2$; reversing that
circle changes the Euler number to $+1$. In particular
\[
 \pi_1(Y)=\ker\chi,\qquad H_1(Y;\mathbb Z)=\mathbb Z^{12}.
 \tag{E.S24}\hypertarget{tag-02CS}{}\label{maps:eq:E.S24}
\]
The return diffeomorphism of \eqref{maps:eq:E.S22} has infinite order even
among unbased homotopy classes of self-equivalences of $Y$,
although its induced outer automorphism of $\pi_1(Y)$ has
exact order twelve. No connected finite cover of $P$ is
homotopy equivalent to a product $S^1\times Z$.


\end{theorem}


\begin{proof}
\begin{enumerate}[label={$\langle1\rangle$\arabic*.},leftmargin=*]
\item \emph{A closed one-form with a nonzero derivative at every point.}
Use the smooth map $c=(e_1,e_2):P\to N$ constructed in \ref{maps:src:E.2.9}. The manifold $N$ has the locally free $T^2$ action
induced by the two weighted Hopf actions. Take a closed
one-form representing the primitive integral class $\chi$
on $N$ and average it over this compact torus. Its average
$\alpha$ is closed, represents the same integral class,
and is invariant. If $V_i$ is the fundamental vector field
with time-one orbit $h_i$, Cartan's formula gives
$d(\alpha(V_i))=0$. Integrating along an ordinary orbit
therefore gives
\[
                  \alpha(V_1)=12,\qquad \alpha(V_2)=-12.
 \tag{E.S25}\hypertarget{tag-02CT}{}\label{maps:eq:E.S25}
\]
The integrality of the periods integrates $\alpha$ to a
smooth map $\rho_N:N\to\mathbb R/\mathbb Z$.

\item Choose the two cusp replacements defining $e_i$ in disjoint
base discs, themselves disjoint from the other factor's
critical values. At every point of $P$, at least one factor
is a regular torus at which its $e_i$ is the literal character
quotient of \SeriesRef{mapsI:src:A.1.1}. In a local flat marking of that factor,
the constant vertical vector $g_i$ maps to $V_i$. Holding
the other factor fixed gives a tangent vector to $P$, and
\[
                     (c^*\alpha)(g_i)=12\quad\hbox{or }-12.
\]
Thus $c^*\alpha$ is nowhere zero.  The map $\rho=\rho_Nc$ is a
proper submersion. Equivalently, choose a vector field $W$
with $(c^*\alpha)(W)=1$; its complete flow supplies local
trivializations and a return diffeomorphism. The fibre is
connected because $\chi$ is primitive: the image in
$\pi_1(S^1)$ of a connected component of a bundle's total
space has index equal to the number of fibre components.
The homotopy sequence now proves $\pi_1(Y)=\ker\chi$.

\item \emph{The exact surface cover and its circle Euler class.}
Write $\Gamma=G/\langle h_1,h_2\rangle$ for the oriented
orbifold group of $S^2(3,4,3,4)$. The character modulo
twelve descends to the surjection
\[
 \overline\chi:\Gamma\longrightarrow\mathbb Z/12,
       \qquad(a_1,a_2,b_1,b_2)\longmapsto(4,-3,-4,3).
 \tag{E.S26}\hypertarget{tag-02CU}{}\label{maps:eq:E.S26}
\]
Its restriction to every cone group is injective. Hence its
kernel $\Gamma_0$ is torsion free and defines a degree-twelve
ordinary surface cover of the orbifold. Its Euler characteristic
is $12(-5/6)=-10$, so that surface is $\Sigma_6$. The image
of $\ker\chi$ in $\Gamma$ is exactly $\Gamma_0$: a lift of
an element of $\Gamma_0$ has character a multiple of twelve,
which can be cancelled by a power of $h_1$. Its central
kernel is the primitive subgroup
\[
 \ker\chi\cap\langle h_1,h_2\rangle
                =\mathbb Z\langle h_1h_2\rangle.
 \tag{E.S27}\hypertarget{tag-02CV}{}\label{maps:eq:E.S27}
\]

\item We compute the circle Euler number in the following orientation convention. On either weighted Hopf sphere put
$z_1=\cos r\,e^{i\varphi_1}$ and
$z_2=\sin r\,e^{i\varphi_2}$. For the positive action with
weights $(4,3)$ the smooth invariant connection is
\[
 \vartheta=\tfrac14\cos^2r\,d\varphi_1
                    +\tfrac13\sin^2r\,d\varphi_2.
\]
It evaluates to one on the generator with period $2\pi$.
With the boundary orientation of the unit ball in $\mathbb C^2$,
\[
             \frac1{(2\pi)^2}\int_{S^3}\vartheta\wedge d\vartheta
                         =\frac1{12}.
 \tag{E.S28}\hypertarget{tag-02CW}{}\label{maps:eq:E.S28}
\]
For the associated complex line whose unit circle is this
positive circle bundle, the convention agreeing with the
tautological Hopf line is $c_1=-[d\vartheta/(2\pi)]$.
Thus each orbifold circle Euler number is $-1/12$.
Both shifted base identifications preserve the complex
orientation. Pullback to $\Sigma_6$ multiplies the two Euler
numbers by twelve, giving $(-1,-1)$; this is also the
covering formula of \cite[Theorem 3.6]{rootScott83}.
In the integral fibre basis
\[
                  v=(1,1),\qquad w=(0,1)
\]
the determinant is one, the Euler vector is $-v$, and
$\chi(v)=0$, $\chi(w)=-12$. Consequently the circle kernel
has Euler number $-1$ with generator $h_1h_2$.

\item One can see this bundle directly in $N$. Invariance of
$\alpha$ implies
$\rho_N((s_1,s_2)x)=\rho_N(x)+12(s_1-s_2)$, where
$s_i\in\mathbb R/\mathbb Z$. A level contains twelve
diagonal circles over an ordinary base point. The peripheral
permutations \eqref{maps:eq:E.S26} give exactly the connected cover
$\Sigma_6$; they have no stabilizers at the cone points.
The diagonal circle acts freely, since at every cone only
one of the two Hopf actions has a nontrivial stabilizer.
The level is therefore the circle bundle just computed.
This identifies its fundamental group with $\ker\chi$
and gives a three-dimensional aspherical model.
The homology sequence of a circle bundle of Euler number
$\pm1$ gives its first homology $\mathbb Z^{12}$, proving
the second assertion in \eqref{maps:eq:E.S24}.

\item \emph{Faithful central action and the return obstruction.}
The centre of $G$ is exactly $\langle h_1,h_2\rangle$:
the quotient is a nonelementary oriented hyperbolic orbifold
group with trivial centre. This centre acts faithfully on
$\pi_2(P)$. Indeed if $h_1^a h_2^b$ acts as the identity,
localize first at $h_1-1$ and apply \ref{maps:src:E.2.12}. On
$R_1/(h_2-1)$ its action is multiplication by $h_1^a$.
This module is free over the central ring
$\mathbb Z[h_1^{\pm1},(h_1-1)^{-1}]$ after choosing coset
representatives, so this is the identity only if $a=0$.
The symmetric localization then gives $b=0$.

\item If $H<G$ has finite index, its image in $\Gamma$ is again
nonelementary. Its centralizer in $\Gamma$ is trivial:
commuting with two hyperbolic elements with distinct axes
fixes at least three boundary points and hence is the
identity. It follows that
$Z(H)=H\cap Z(G)$. A homotopy product $S^1\times Z$ for
the corresponding finite cover would supply an infinite
central element acting trivially on $\pi_2$, namely the
circle factor; this contradicts the faithful action just
proved. If a positive power of the return of \eqref{maps:eq:E.S22} were
homotopic to the identity, its mapping torus would be
homotopy equivalent to $S^1\times Y$. That mapping torus
is the corresponding finite cyclic cover of $P$, so this
too is impossible.

\item Finally take $t=a_1a_2$, for which $\chi(t)=1$, as a
stable letter. The image of $\ker\chi$ in $\Gamma$ is
$\Gamma_0$, so the same centralizer argument gives
$C_G(\ker\chi)=Z(G)$. The automorphism
$\operatorname{Ad}(t^n)|_{\ker\chi}$ is inner precisely
when $t^n$ differs from an element of $\ker\chi$ by a
central element. Since $\chi(Z(G))=12\mathbb Z$, this
holds precisely for $12\mid n$. Thus the outer action has
order twelve, while the return has infinite homotopy
order. This return is the circle fibration of the shifted
ten-manifold; it is distinct from the six-sphere parameter
returns computed in \SeriesRef{maps:src:B.6.24}.

\end{enumerate}
\end{proof}

\begin{theorem}[The cyclic-cover splice and its separating torus twist]\hypertarget{tag-02ME}{}\label{maps:src:E.2.14}
\StatementRecord{maps:src:E.2.14}{proved}{02ME}{\SeriesRef{maps:src:E.2.13}, \SeriesRef{maps:src:E.2.6}, \SeriesRef{maps:src:E.2.10}, \eqref{maps:eq:E.S34}, \eqref{maps:eq:E.S35}, \eqref{maps:eq:E.S37}, \eqref{maps:eq:E.S39}, \eqref{maps:eq:E.S38}, \eqref{maps:eq:E.S41}, \eqref{maps:eq:E.S43}, \eqref{maps:eq:E.S36}, \eqref{maps:eq:E.S40}, \eqref{maps:eq:E.S44}, \eqref{maps:eq:E.S45}, \SeriesRef{maps:src:E.2.7}, \eqref{maps:eq:E.S47}, \eqref{maps:eq:E.S46}}


Let $E_{12}\to E$ be the connected cyclic cover defined by
reduction modulo twelve of the positive meridian character.
The circle fibration of \ref{maps:src:E.2.13} may be chosen with fibre
\[
 Y=(E_{12}\times T^3)\ \cup_{T^8}\ (T^3\times E_{12}),
 \tag{E.S29}\hypertarget{tag-02CX}{}\label{maps:eq:E.S29}
\]
with the following exact boundary marking. Write
$\theta_i'$ for the lifted normal circle and use the
two regular-torus markings. Then
\[
       \theta_1'=\gamma_2-\gamma_1,\qquad
       \theta_2'=\gamma_1-\gamma_2.
 \tag{E.S30}\hypertarget{tag-02CZ}{}\label{maps:eq:E.S30}
\]
There is a return representative $\psi$ whose twelfth power
is the separating torus twist
\[
 \begin{gathered}
 T(x,s)=(x+\eta(s)(g_1+g_2),s)\quad\hbox{on }T^8\times[-1,1],
 \\
 T=\operatorname{id}\text{ elsewhere}.
 \end{gathered}
 \tag{E.S31}\hypertarget{tag-02D0}{}\label{maps:eq:E.S31}
\]
Here $\eta$ is zero near $-1$ and one near $1$, and circle
coordinates have period one. In particular \eqref{maps:eq:E.S31} is the
identity near both ends of its collar. Its rational
cohomological variation satisfies
\[
 (T^*-I)^2=0,\qquad
 \operatorname{rank}(T^*-I)|_{H^j(Y;\mathbb Q)}=
 \begin{cases}1&j=4,5,\\0&j\ne4,5.\end{cases}
 \tag{E.S32}\hypertarget{tag-02D1}{}\label{maps:eq:E.S32}
\]
The rational Betti numbers of the fibre are
\[
       (b_0,\ldots,b_9)(Y)
          =(1,12,121,407,693,693,407,121,12,1).
 \tag{E.S33}\hypertarget{tag-02D2}{}\label{maps:eq:E.S33}
\]
Every eigenvalue of $\psi^*$ is a twelfth root of unity.
At eigenvalue one there is exactly one Jordan block of
size two in each of degrees four and five; all its other
blocks, and all blocks at other eigenvalues, have size one.
Thus the infinite order is already visible on ordinary
rational cohomology. \end{theorem}


\begin{proof}
\begin{enumerate}[label={$\langle1\rangle$\arabic*.},leftmargin=*]
\item \emph{Adapting the fibration to the splice.}
Cut the base along the circle used in \ref{maps:src:E.2.6}, with the two
critical triples on opposite sides. Over the first disc the
second factor is a regular marked $T^4$, and over the second
disc the first factor is regular. Let
$\xi_E:E\to\mathbb R/\mathbb Z$ represent the positive
meridian class. Its restriction to the boundary can be
chosen exactly
\[
                         \xi_E=\theta+12\gamma.
 \tag{E.S34}\hypertarget{tag-02D3}{}\label{maps:eq:E.S34}
\]
Indeed this is its integral restriction by \ref{maps:src:E.2.10}; homotopy
extension makes the representative equal to that linear
map on a collar. The two maps
\[
        \rho_A=\xi_{E_1}-12\gamma_2,
             \qquad \rho_B=12\gamma_1-\xi_{E_2}
 \tag{E.S35}\hypertarget{tag-02D4}{}\label{maps:eq:E.S35}
\]
agree exactly on the common boundary, where
$\theta_2=-\theta_1$. Smooth them with this common product
collar. The regular-factor circle derivative is $-12$ on
the first piece and $12$ on the second, so they give a
proper submersion throughout, including both cusp pieces.
These are also permissible choices in \cref{maps:src:E.2.13}: on the
quotient $N$, an invariant closed one-form differs from
these representatives only by the differential of a
single-valued function on each of the two pieces. Absorb
that function by translation in the unused regular circle.
Thus this construction uses the same quotient and
primitive character.

\item A level of $\xi_E-12\gamma_2$ is the pullback of the
degree-twelve circle cover along $\xi_E$, times the three
remaining regular circles. The pullback is $E_{12}$.
Its boundary is connected because the normal meridian has
character one. The covering of that boundary has
$\theta=12\theta'$ and unchanged $\gamma,u,w,\delta$.
Equation \eqref{maps:eq:E.S34} in the level condition gives precisely
\eqref{maps:eq:E.S30}. This proves \eqref{maps:eq:E.S29} with its smooth
boundary maps.

\item Choose the vector field integrating \eqref{maps:eq:E.S35} to be
$-g_2/12$ on the first piece off its collar and $g_1/12$
on the second. On the common regular collar interpolate by
\[
                     W_s=\frac{\eta(s)g_1-(1-\eta(s))g_2}{12}.
 \tag{E.S36}\hypertarget{tag-02D5}{}\label{maps:eq:E.S36}
\]
It has $d\rho(W_s)=1$, and it preserves the collar
coordinate $s$ and the original base point. Its flow
therefore defines a return $\psi$. At time twelve it is
the identity off the collar, since the two indicated
regular circles have period one. On the collar it
translates by $\eta(s)g_1-(1-\eta(s))g_2$; removing the
integral translation $-g_2$ gives \eqref{maps:eq:E.S31}. This comparison uses a globally regular factor on each piece.

\item \emph{The variation map and the invariant boundary calculation.}
Put $D=T^8$ and let $r:H^j(Y;\mathbb Q)\to H^j(D;\mathbb Q)$
be restriction. Denote the Mayer--Vietoris connecting map
by $\partial:H^{j-1}(D;\mathbb Q)\to H^j(Y;\mathbb Q)$.
Writing $v=g_1+g_2$, the collar pullback of a translation
invariant form is
\[
       T^*\omega=\omega+d\eta\wedge\iota_v\omega.
\]
Average its restriction to the torus and change the global
representative by an exact form to make this equality
applicable on the collar. The difference represents the
connecting class, and hence, with the connecting sign
chosen by the increasing collar coordinate,
\[
                          T^*-I=\partial\,\iota_v\,r.
 \tag{E.S37}\hypertarget{tag-02D6}{}\label{maps:eq:E.S37}
\]
As $r\partial=0$, its square is zero.

\item The computation of \eqref{maps:eq:E.S37} needs no further information on $H^*(E_{12})$. Its cyclic deck action is a translation on the
boundary torus, hence acts trivially on boundary cohomology.
Every noninvariant deck eigenspace therefore restricts to
zero. Rational transfer identifies the invariant part
with $H^*(E)$, whose marked restriction is \ref{maps:src:E.2.6}. Its
pullback to \eqref{maps:eq:E.S30} is
\[
 R_{k,1}(a)=\phi_k(a)+12(\gamma_2-\gamma_1)a,
 \qquad
 R_{k,2}(b)=\phi_k(b)+12(\gamma_1-\gamma_2)b.
 \tag{E.S38}\hypertarget{tag-02D7}{}\label{maps:eq:E.S38}
\]
Use the other factor's three covectors $u,w,\delta$,
and take exterior products to obtain the two subspaces
$U_j^k\subset\bigwedge^k\mathbb Q^8$.
The image of $r$ is exactly $U_1^k\cap U_2^k$, and the
kernel of $\partial$ is $U_1^{k-1}+U_2^{k-1}$.
Consequently the rank in \eqref{maps:eq:E.S32} is exactly the rank of
\[
 U_1^k\cap U_2^k\ \stackrel{\iota_v}{\longrightarrow}
 \bigwedge^{k-1}\mathbb Q^8/(U_1^{k-1}+U_2^{k-1}).
 \tag{E.S39}\hypertarget{tag-02D8}{}\label{maps:eq:E.S39}
\]

\item Here are the full dimensions for this small rational
calculation. Let $d_k$ be twice the degree-$k$ dimension
of $H^*(E)\otimes H^*(T^3)$ and let
$r_k=\dim(U_1^k+U_2^k)$:
\[
 \begin{array}{c|rrrrrrrrr}
 k&0&1&2&3&4&5&6&7&8\\ \hline
 d_k&2&8&20&44&70&68&36&8&0\\
 r_k&1&8&19&34&47&46&27&8&0
 \end{array}
 \tag{E.S40}\hypertarget{tag-02DA}{}\label{maps:eq:E.S40}
\]
Contraction in \eqref{maps:eq:E.S39} has rank one only for
$k=4,5$ and is zero in every other degree. These entries result from exterior products of \eqref{maps:eq:E.S38}
and row reduction of the induced contraction maps. One can normalize the remaining sign in
$\phi_2$ by the ring structure: with
$r_1(1)=\theta+12\gamma$, the equality
$r_1(1)r_2(\delta)=r_3(a)$ forces the coefficient
$c_2=-2$, and then $c_3=2$, in the notation of \ref{maps:src:E.2.6}.



\item \emph{The finite-cover cohomology needed for the full Betti numbers.}
The full rational cohomology of $E_{12}$ is computed as follows.
Let $\mathcal L_\lambda$ be the rank-one system on $E$
whose meridian acts by $\lambda$, with $\lambda^{12}=1$.
On a regular torus it is trivial. Its three base transports
in vertical degree $q$ are
\[
 \lambda^4\bigwedge^q A_1^{-t},\qquad
 \lambda^{-3}\bigwedge^q A_2^{-t},\qquad
                         \bigwedge^q M_0^{-t}.
 \tag{E.S41}\hypertarget{tag-02DB}{}\label{maps:eq:E.S41}
\]
The exponents carry the affine meridian character.

\item At a finite cap, averaging on its finite root cover
identifies the cap cohomology with the invariant part of
\eqref{maps:eq:E.S41}. At the unchanged cusp the system is trivial,
because $g$ has character twelve and $u$ has character zero.
Its rational specialization is an isomorphism onto the
monodromy invariants in every degree.  The cusp Betti numbers are
$(1,2,4,2,1)$ by its normalization calculation. In degree
two, the integral surjectivity of the specialization on homology (proof of \cref{maps:src:E.2.5}) gives injectivity on cohomology;
the monodromy-invariant subspace has the same dimension
four. Its image therefore includes $uw$ and $\gamma\delta$.
In degree one the image is $\langle\gamma,u\rangle$;
their products with those two degree-two classes give the
two independent invariant classes in degree three and
the top class in degree four. Thus all the specialization
maps have the asserted ranks, with no additional
skyscraper summand. The proof of \cref{maps:src:E.2.5} gives these specialization assertions integrally in all degrees.

\item It follows that $R^q f_*\mathcal L_\lambda=j_*L_q$ over
the disc containing the three critical values. On this
disc $H^p(j_*L_q)=0$ for $p\geq2$. To check this directly,
retract the punctured disc to a free graph on three loops.
Restriction of a graph cocycle to the three local
coinvariant spaces is onto: its values on the three free
generators can be chosen independently. The exact
restriction sequence therefore kills $H^2(j_*L_q)$.
If $d_q=\binom4q$, $g_q$ is the common fixed dimension of
\eqref{maps:eq:E.S41}, and $k_{i,q}$ are their three individual fixed
dimensions, this same graph calculation gives
\[
 \dim H^0(j_*L_q)=g_q,\qquad
 \dim H^1(j_*L_q)=2d_q+g_q-\sum_i k_{i,q}.
 \tag{E.S42}\hypertarget{tag-02DC}{}\label{maps:eq:E.S42}
\]
The Leray sequence has only these two base degrees and
degenerates. Exact exterior-matrix ranks give the following
table; a row applies to each character of the indicated
order:
\[
 \begin{array}{c|rrrrrr}
 \operatorname{ord}(\lambda)&b_0&b_1&b_2&b_3&b_4&b_5\\ \hline
 1 &1&1&4&6&4&0\\
 2 &0&0&4&6&4&0\\
 3 &0&0&3&4&3&0\\
 4 &0&0&3&4&3&0\\
 6 &0&1&5&6&5&1\\
 12&0&1&4&4&4&1
 \end{array}
 \tag{E.S43}\hypertarget{tag-02DD}{}\label{maps:eq:E.S43}
\]
Here $b_j=\dim H^j(E;\mathcal L_\lambda)$ and $b_6=0$.
For the nontrivial characters all $g_q$ are zero; for the
trivial one they are all one. For \eqref{maps:eq:E.S41}--\eqref{maps:eq:E.S43}, represent the
cyclotomic fields by rational companion matrices and compute each local
fixed space by exact row reduction.
Summing the rows with multiplicities $1,1,2,2,2,4$ gives
\[
             (b_0,\ldots,b_6)(E_{12})=(1,7,46,56,46,6,0).
 \tag{E.S44}\hypertarget{tag-02DE}{}\label{maps:eq:E.S44}
\]
Both sides have Euler characteristic $24$, as required by
the degree-twelve covering of $E$.

\item \emph{The Betti and Jordan calculations on the fibre.}
The return on each of the two pieces of \eqref{maps:eq:E.S29}, up to
isotopy allowing its boundary to move, is its cyclic deck
transformation times the identity of $T^3$. This follows
also by moving the interpolation in \eqref{maps:eq:E.S36} entirely into
the other piece's collar. It acts trivially on the
cohomology of $D$. The generalized eigenvalue-one part of
the Mayer--Vietoris sequence therefore uses just the
invariant part $H^*(E)$ in \eqref{maps:eq:E.S40}. Its dimensions are
\[
       (1,0,1,19,45,45,19,1,0,1).
 \tag{E.S45}\hypertarget{tag-02DF}{}\label{maps:eq:E.S45}
\]
Every other eigenspace is the direct sum of the two piece
eigenspaces, because the boundary has none. In particular
those eigenspaces are semisimple. Their combined Poincar\'e
polynomial is
\[
     2(1+t)^3(6t+42t^2+50t^3+42t^4+6t^5),
 \tag{E.S46}\hypertarget{tag-02DG}{}\label{maps:eq:E.S46}
\]
by \eqref{maps:eq:E.S43}--\eqref{maps:eq:E.S44}. Adding \eqref{maps:eq:E.S45} proves \eqref{maps:eq:E.S33}.

\item Finally the Wang sequence for $P=M_\psi$ and the rational
Betti numbers in \ref{maps:src:E.2.7} give the numbers of Jordan blocks
at eigenvalue one as
\[
          (1,0,1,19,44,44,19,1,0,1).
 \tag{E.S47}\hypertarget{tag-02DH}{}\label{maps:eq:E.S47}
\]
The only possible block sizes are one and two: the
twelfth power is \eqref{maps:eq:E.S31}, and \eqref{maps:eq:E.S37} squares to zero.
Subtracting \eqref{maps:eq:E.S47} from \eqref{maps:eq:E.S45} gives exactly one block
of size two in each of degrees four and five. This agrees with the direct variation ranks \eqref{maps:eq:E.S39}.
All other eigenvalues and blocks are already described
by \eqref{maps:eq:E.S43} and \eqref{maps:eq:E.S46}. \end{enumerate}
\end{proof}

\begin{theorem}[Integral fibre groups and the coefficient of the torus twist]\hypertarget{tag-02MF}{}\label{maps:src:E.2.15}
\StatementRecord{maps:src:E.2.15}{proved}{02MF}{\eqref{maps:eq:E.S44}, \eqref{maps:eq:E.S33}, \eqref{maps:eq:E.S42}, \eqref{maps:eq:E.S53}, \eqref{maps:eq:E.S29}, \eqref{maps:eq:E.S38}, \eqref{maps:eq:E.S40}, \eqref{maps:eq:E.S37}, \eqref{maps:eq:E.S55}}


The groups $H^*(E_{12};\mathbb Z)$ are torsion free, with
the ranks in \eqref{maps:eq:E.S44}. For the fibre $Y$, the only
torsion in cohomology is
\[
       H^3(Y;\mathbb Z)_{\rm tors}
       =H^7(Y;\mathbb Z)_{\rm tors}=(\mathbb Z/6)^2.
 \tag{E.S48}\hypertarget{tag-02DI}{}\label{maps:eq:E.S48}
\]
Equivalently, its integral homology has the ranks \eqref{maps:eq:E.S33}
and torsion $(\mathbb Z/6)^2$ exactly in degrees two and six.
The twelfth return $T=\psi^{12}$ acts as the identity on
all integral cohomology except in degrees four and five.
In each of those two free groups it has, in an integral
basis, the form
\[
             \begin{pmatrix}1&72\\0&1\end{pmatrix}
                                  \oplus I_{691}.
 \tag{E.S49}\hypertarget{tag-02DJ}{}\label{maps:eq:E.S49}
\]

For $n\geq1$, let $P_{12n}$ be the connected cyclic cover
of $P$ of degree $12n$ defined by $\chi$. Its Betti numbers,
independent of $n$, are
\[
       (1,13,133,528,1099,1384,1099,528,133,13,1).
 \tag{E.S50}\hypertarget{tag-02DL}{}\label{maps:eq:E.S50}
\]
Its integral homology torsion is exactly
\[
 H_j(P_{12n};\mathbb Z)_{\rm tors}=
 \begin{cases}
       (\mathbb Z/6)^2,&j=2,3,6,7,\\
       \mathbb Z/(72n),&j=4,5,\\
       0,&\text{otherwise}.
 \end{cases}
 \tag{E.S51}\hypertarget{tag-02DM}{}\label{maps:eq:E.S51}
\]
In particular these covers are pairwise nonhomotopy-equivalent.
\end{theorem}


\begin{proof}
\begin{enumerate}[label={$\langle1\rangle$\arabic*.},leftmargin=*]
\item \emph{The integral cohomology of the cover.}
Apply the proof of \eqref{maps:eq:E.S42} with the coefficient lattice
$\mathbb Z[C_{12}]$, using its regular permutation matrix
$P_{12}$. In vertical degree $q$ set
\[
 \begin{aligned}
 T_{1,q}=P_{12}^4\otimes\bigwedge^q A_1^{-t},\quad
 &T_{2,q}=P_{12}^{-3}\otimes\bigwedge^q A_2^{-t},\\
 T_{0,q}=I_{12}\otimes\bigwedge^q M_0^{-t},
 &V_q=\mathbb Z^{12\binom4q}.
 \end{aligned}
 \tag{E.S52}\hypertarget{tag-02DN}{}\label{maps:eq:E.S52}
\]
This coefficient system computes $H^*(E_{12};\mathbb Z)$.
 At the first finite
cap the cover has four components, each the ordinary
degree-three torus root cover; at the second there are
three ordinary degree-four torus root covers. An invariant
vector is determined by its value on one sheet of such
an orbit, with the other values its integral monodromy
translates. This identifies cap cohomology with the full
primitive invariant lattice. The cusp has twelve unchanged
components and the integral specialization of the proof of \cref{maps:src:E.2.5}.
Consequently its direct-image sheaf is exactly the integral
$j_*L_q$, including at every critical point.

\item The three graph-cocycle values again surject onto the
three local coinvariant groups, integrally. Thus $H^2(j_*L_q)=0$,
and its remaining cohomology is the kernel and cokernel
of the exact lattice map
\[
 \begin{gathered}
 V_q\longrightarrow
        \bigoplus_{i=0}^2\operatorname{im}(T_{i,q}-I),\\
           v\longmapsto((T_{i,q}-I)v)_i.
 \end{gathered}
 \tag{E.S53}\hypertarget{tag-02DO}{}\label{maps:eq:E.S53}
\]
Use Hermite column bases for the three image lattices. The source and target ranks,
and the rank of \eqref{maps:eq:E.S53}, are respectively
\[
 \begin{array}{c|rrrrr}
 q&0&1&2&3&4\\ \hline
 \operatorname{rank}V_q&12&48&72&48&12\\
 \operatorname{rank}\bigoplus\operatorname{im}(T_{i,q}-I)
                      &17&92&126&92&17\\
 \operatorname{rank}(\text{E.S53})&11&47&71&47&11
 \end{array}
 \tag{E.S54}\hypertarget{tag-02DP}{}\label{maps:eq:E.S54}
\]
Every nonzero Smith entry is one. Hence $H^0(j_*L_q)$
has rank one and $H^1(j_*L_q)$ is free with ranks
$(6,45,55,45,6)$. The two-column Leray filtration has
free end groups and therefore splits as abelian groups.
This proves the torsion-free assertion and \eqref{maps:eq:E.S44}
integrally. The integer matrices are determined by the Hermite image bases in
\eqref{maps:eq:E.S53}; their Smith reductions give the displayed invariants.

\item \emph{The integral images.}
Each piece $A=E_{12}\times T^3$ in \eqref{maps:eq:E.S29} has free
integral homology. By Poincar\'e--Lefschetz duality,
$H^{k+1}(A,\partial A;\mathbb Z)$ is free. The cokernel
of $H^k(A;\mathbb Z)\to H^k(\partial A;\mathbb Z)$
embeds in that free group, so the boundary image is a
primitive lattice. Its rational span has already been
computed in \eqref{maps:eq:E.S38}. Therefore the integral
image is precisely the saturation of that displayed
image in $\bigwedge^k\mathbb Z^8$.

\item Choose saturated column bases $\overline A_k,\overline B_k$
for the two images. The combined matrix
$[\overline A_k\ \overline B_k]$ has the ranks in the
second row of \eqref{maps:eq:E.S40}. All its nonzero Smith entries are
one except that, in each of degrees two and six, exactly
two entries are six. The Mayer--Vietoris exact sequence
has a free kernel term on its right, since all piece
cohomology is free. Thus its abelian extension splits,
and these Smith entries give exactly \eqref{maps:eq:E.S48}. The free
ranks remain \eqref{maps:eq:E.S33}.

\item \emph{The integral variation, including possible torsion components.}
Formula \eqref{maps:eq:E.S37} holds integrally. One can use the singular
prism of the loop of translations instead of differential
forms: the action map $S^1\times T^8\to T^8$ has
cohomology formula
\[
 a\longmapsto 1\otimes a+[S^1]^*\otimes\iota_v a.
\]
The difference between the two collar homotopies is
exactly this loop, so its connecting class is
$\partial\iota_v r(a)$. This proves the integral formula
also on torsion and still gives $(T^*-I)^2=0$.

\item For complete lattice control, take an integral kernel
basis of $[\overline A_k\ -\overline B_k]$. Its first
block, multiplied by $\overline A_k$, is a basis of the
common boundary lattice. Contract it by $v$ and reduce
modulo $[\overline A_{k-1}\ \overline B_{k-1}]$ using
the integral Smith transforms. In degrees four and five
the result has rank one and unique nonzero Smith entry
$72$. In degree three its two possible components
modulo six are both zero; in degree seven the common
boundary image is zero. All the remaining reductions
are zero. The reductions use the integral unimodular Smith transforms.
In degrees four and five the connecting cokernel is
free and embeds as a primitive summand of $H^k(Y)$,
because the following Mayer--Vietoris kernel is free.
Thus the divisibility in the whole cohomology group is $72$.

\item A rank-one integral map of Smith entry $72$ has the
form $72u\ell$, where $u$ and $\ell$ are primitive.
Its square is zero precisely when $\ell(u)=0$.
Extend $u$ to a basis of $\ker\ell$ and choose a vector
on which $\ell$ is one. This gives \eqref{maps:eq:E.S49}. The same
calculation proves identity on every other integral
cohomology group, including its torsion.

\item \emph{The finite-cover groups and the two potentially nontrivial extensions.}
The cover $P_{12n}$ is the mapping torus of $T^n$.
Since $(T^*-I)^2=0$, its variation is $n(T^*-I)$,
with Smith entry $72n$ in degrees four and five.
For every $k$ the integral Wang sequence gives
\[
 0\longrightarrow\operatorname{coker}(T^{n*}-I)_{k-1}
 \longrightarrow H^k(P_{12n};\mathbb Z)
 \longrightarrow\ker(T^{n*}-I)_k\longrightarrow0.
 \tag{E.S55}\hypertarget{tag-02DQ}{}\label{maps:eq:E.S55}
\]
The ranks immediately give \eqref{maps:eq:E.S50}. Whenever the right
group is free, \eqref{maps:eq:E.S55} splits as an abelian sequence.
In particular degrees four and eight have torsion
$(\mathbb Z/6)^2$, contributed by $H^3(Y)$ and $H^7(Y)$
respectively, while degrees five and six have torsion
$\mathbb Z/(72n)$. Every other such degree has no torsion.

\item The only right groups in \eqref{maps:eq:E.S55} that are not free
occur in degrees three and seven.  Instead use the
nondegenerate integral torsion linking pairing of the
closed oriented ten-manifold $P_{12n}$: cohomological
torsion in degree $k$ is Pontryagin dual to that in
degree $11-k$. The already fixed torsion in degree
eight therefore forces the degree-three torsion to
be $(\mathbb Z/6)^2$, and the fixed degree-four torsion
forces the same answer in degree seven. The free
ranks are already known. This resolves precisely
the two possibly absorbing extensions. The universal
coefficient theorem now translates these cohomology
groups to \eqref{maps:eq:E.S51}. The order $72n$ in the middle
groups distinguishes the covers for distinct $n$.
\end{enumerate}
\end{proof}

\begin{theorem}[The full complement module and its peripheral map]\hypertarget{tag-02MG}{}\label{maps:src:E.2.16}
\StatementRecord{maps:src:E.2.16}{proved}{02MG}{\eqref{maps:eq:E.S1}, \SeriesRef{maps:src:E.2.17}, \SeriesRef{maps:src:E.2.21}, \SeriesRef{maps:src:E.2.12}, \eqref{maps:eq:E.C5}, \eqref{maps:eq:E.C8}, \eqref{maps:eq:E.C11}, \eqref{maps:eq:E.C7}, \eqref{maps:eq:E.C14}, \eqref{maps:eq:E.C12}, \eqref{maps:eq:E.C16}, \eqref{maps:eq:E.C17}, \SeriesRef{maps:src:E.2.6}}


Let $K=\pi_1(E)$ have the marking \eqref{maps:eq:E.S1}, and put
\[
 R=\mathbb ZK,\qquad Q=K/\langle h\rangle=C_3*C_4,
 \qquad \overline R=\mathbb ZQ.
\]
Use right deck modules; conversion to left modules applies the
involution $g\mapsto g^{-1}$ to every coefficient. Then
\[
 \pi_2(E)\simeq R\oplus\overline R^{\,2}\oplus L,
 \qquad
 L=\langle x,t\mid t(a-1)=t(b-1)=0,
                         \ x(1+b)+3t=0\rangle_{\overline R}.
 \tag{E.C1}\hypertarget{tag-02A4}{}\label{maps:eq:E.C1}
\]
Equivalently, replacing $x$ by $x+t$ and eliminating $t$ gives
\[
 L\simeq\overline R\big/
       \big((b^2-1)\overline R+(1+b)(a-1)\overline R\big).
 \tag{E.C2}\hypertarget{tag-02AG}{}\label{maps:eq:E.C2}
\]
There is a nonsplit exact sequence
\[
 0\longrightarrow\mathbb Z\langle t\rangle
  \longrightarrow L\longrightarrow
  \mathbb Z_{-}\otimes_{\mathbb ZC_4}\overline R
  \longrightarrow0,
 \tag{E.C3}\hypertarget{tag-02AJ}{}\label{maps:eq:E.C3}
\]
where $b$ acts as $-1$ on $\mathbb Z_{-}$ and the first term
has trivial $Q$ action. Its extension class is the nonzero
element of an order-two extension group. In the connected
lift of the peripheral $T^5$ to the universal cover of $E$,
the degree-two classes $uw,ud,wd$ map respectively to
\[
                           t,\quad0,\quad0.
 \tag{E.C4}\hypertarget{tag-02AK}{}\label{maps:eq:E.C4}
\]
Thus the central torsion summand of the module is annihilated
by $h-1$ itself. Formula \eqref{maps:eq:E.C1} is a full integral module
calculation, including the central value $h=1$.  There is also a chain-level peripheral comparison.
For the splitting constructed below, let $e$ generate the
free $R$ summand and let $y,z$ generate the two copies of
$\overline R$. The projection of the peripheral degree-two
chain map to \eqref{maps:eq:E.C1} is
\[
 \begin{array}{c|rrrrrr}
 \text{cell}&gu&gw&gd&uw&ud&wd\\ \hline
 \text{value}&0&0&6t&t&0&0
 \end{array}
 \tag{E.C19a}\hypertarget{tag-02AE}{}\label{maps:eq:E.C19a}
\]
and, writing $\Theta$ for the outer base circle with
group image $\mu=ab$, use the cone convention: the
notation $v\Theta$ denotes the suspended cell
$\Theta\times v$, with the interval factor first.
In this convention,
\[
 \begin{array}{c|l}
  g\Theta&eh-2x-6yb+z(5-2b)\\
  u\Theta&x-yb+zb\\
  w\Theta&-y(1+b)+z\\
  d\Theta&x+2t+z(b-1).
 \end{array}
 \tag{E.C19b}\hypertarget{tag-02AF}{}\label{maps:eq:E.C19b}
\]
These are values on chains, not assertions that all
the displayed cells are cycles in the cover. In
particular the free coefficient $h$ on $g\Theta$
is nonzero even though the free summand receives
zero from the lifted peripheral homology in \eqref{maps:eq:E.C4}.


\end{theorem}


\begin{proof}
\begin{enumerate}[label={$\langle1\rangle$\arabic*.},leftmargin=*]
\item We first give the local chain comparison; the homology maps alone do not determine \eqref{maps:eq:E.C1}.

\item \emph{The cusp boundary and its derived attachment.}
Write $A=\mathbb Z[h,h^{-1}]$ and $q=h-1$. Let
$\widehat W$ be the gamma cover used in \ref{maps:src:E.2.12}, and let
$\widehat M$ be the corresponding boundary cover. Ordinary
chains with the deck action give finite free $A$ chain
models. The homology of $\widehat W$ is
\[
 H_0=H_1=A/(q),\quad H_2=A\oplus(A/(q))^2,
 \quad H_3=0,\quad H_4=A.
 \tag{E.C5}\hypertarget{tag-02AL}{}\label{maps:eq:E.C5}
\]
For the boundary, start with the torus chain complex
$C= A\otimes\bigwedge\langle g,u,w,d\rangle$,
whose differential contracts $g$ and multiplies by $q$.
The cellular monodromy has degree-one matrix
\[
 F(M_0)g=g-hd,\quad F(M_0)u=u+w,
 \quad F(M_0)w=w,\quad F(M_0)d=d.
 \tag{E.C6}\hypertarget{tag-02AM}{}\label{maps:eq:E.C6}
\]
Its higher matrices are exterior powers. The mapping-torus
complex is the cone of $I-F(M_0)$, with differential
\[
 d(c,e)=\big(dc+(I-F(M_0))e,-de\big).
 \tag{E.C7}\hypertarget{tag-02AN}{}\label{maps:eq:E.C7}
\]
For example $d(g\theta)=hd-q\theta$ and
$d(u\theta)=-w$. Cancelling these two unit pairs leaves
the degree-one classes $u$ and $\theta$, with resolutions
$q$ and a unit times $q^2$. In degree two the four surviving
classes are $uw,ud,w\theta,d\theta$, each with resolution
$q$. In particular
\[
 H_0(\widehat M)=A/(q),\quad
 H_1(\widehat M)=A/(q^2)\oplus A/(q),\quad
 H_2(\widehat M)=(A/(q))^4.
 \tag{E.C8}\hypertarget{tag-02AO}{}\label{maps:eq:E.C8}
\]
Geometrically this is the Heisenberg three-manifold times
a circle: the deck transformation shears the meridian
by the $d$ circle. The homology map to \eqref{maps:eq:E.C5} is the
identity in degree zero, projection to $u$ in degree one,
and the primitive projection
$(uw,ud,w\theta,d\theta)\mapsto(U,D,0,0)$ in degree two.
The last assertion uses the integral specialization of the proof of \cref{maps:src:E.2.5} and the collapsed meridian tori.

\item The extension is determined by a relative intersection calculation. Let $S_i$ denote the normalized hexagon
component with gamma index $i$. In its standard Picard
basis $H,E_1,E_2,E_3$, the six boundary curves are
\[
 E_1,\ H-E_1-E_2,\ E_2,\ H-E_2-E_3,
                       \ E_3,\ H-E_1-E_3.
 \tag{E.C9}\hypertarget{tag-02AP}{}\label{maps:eq:E.C9}
\]
The first opposite pair is self-glued in the surviving
$u$ direction; the other two pairs join $S_i$ to its
two gamma neighbours. Put $C_0=-H+E_1$, the primitive
generator of the free kernel of the normalization matrix
$B(h)$ in \ref{maps:src:E.2.12}. The divisor $S_i$ restricts on its
own normalization to $2C_0$: the total central divisor
is principal, and the two self-gluing preimages cancel
their terms in the normal restriction. It restricts
to $-C_0$ on each neighbouring component and to zero
on every other component. Thus the free part of the
equivariant intersection map is
\[
 H_4(\widehat W)=A\longrightarrow
 H_4(\widehat W,\widehat M)\longrightarrow A,
 \qquad 2-h-h^{-1}=-h^{-1}q^2.
 \tag{E.C10}\hypertarget{tag-02A5}{}\label{maps:eq:E.C10}
\]
Poincar\'e--Lefschetz duality over $A$, with its usual
involution, gives the same coefficient, up to a unit,
for the free part of $H_2(\widehat W)\to
H_2(\widehat W,\widehat M)$. The latter relative group
is $A$: the universal-coefficient sequence in degree
four has $\operatorname{Hom}_A(H_4,A)=A$, zero $H_3$,
and no higher extension terms. Indeed every module
in \eqref{maps:eq:E.C5} has a projective resolution of length at
most one. The relative sequence therefore includes
\[
 0\longrightarrow A\xrightarrow{\,\text{unit}\ q^2\,}A
  \longrightarrow A/(q^2)\longrightarrow0.
 \tag{E.C11}\hypertarget{tag-02A6}{}\label{maps:eq:E.C11}
\]
Consequently the boundary map has a unit component in
$\operatorname{Ext}^1_A(A/(q^2),A)=A/(q^2)$.
Every unit of $A/(q^2)=\mathbb Z[\varepsilon]/(\varepsilon^2)$
lifts to a unit $\pm h^n$ of $A$.

\item This determines the low-degree chain map, as follows. All the
modules in \eqref{maps:eq:E.C5} and \eqref{maps:eq:E.C8} have projective dimension
at most one. Their complexes split in the derived
category into their two-term projective resolutions:
the obstructions to splitting consecutive truncations
are extension groups in degree at least two and vanish.
In degrees relevant to $H_2$ of a mapping cone, a map
between these split complexes consists of the displayed
homology maps and extension components from $H_i$ to
$H_{i+1}$. Normalize the component in \eqref{maps:eq:E.C11} by a unit
on the free $H_2$ summand of the target. Its components
to the two torsion $H_2$ summands can then be removed
by target maps $A\to(A/(q))^2$: reduction of the unit
modulo $q$ is $\pm1$. The remaining components from
the $u$ summand are removed by target degree-one
extension maps from $H_1(\widehat W)$ to
$H_2(\widehat W)$. Finally the degree-zero component
is removed using the isomorphism on $H_0$ and a target
extension map from $H_0$ to $H_1$.

\item Each change just described is a unit on the diagonal
of the split complex and is triangular with extension
entries off the diagonal. It is therefore invertible.
Replacing the split resolutions by the original finite
free complexes represents it by a chain homotopy
equivalence.  The change
is applied to the cusp target and its entire incoming
boundary map simultaneously. The other vertices and
all maps into the regular core are unchanged. It thus
gives a chain homotopy equivalence of the global
mapping-cone complex that is the identity on the
external peripheral complex. Tensoring with $R$ preserves these comparisons,
since $R$ is free over $A$.

\item One resulting cusp model in degrees zero through three
has ranks $(1,2,4,2)$, generators
\[
 (1);\quad(g,u);\quad(gu,E,U,D);\quad(R_U,R_D),
\]
and differentials $dg=q$, $d(gu)=qu$,
$dR_U=qU$, $dR_D=qD$. The map \eqref{maps:eq:E.C7} to this model
is, through degree two,
\[
 \begin{array}{c|ccccc}
  &g&u&w&d&\theta\\ \hline
  i_1&g&u&0&0&0
 \end{array},\qquad
 \begin{array}{c|rrrrrr}
  &gu&gw&gd&uw&ud&wd\\ \hline
  i_2&gu&-hD&E&U&D&0
 \end{array},
 \qquad i_2(v\theta)=0.
 \tag{E.C12}\hypertarget{tag-02A7}{}\label{maps:eq:E.C12}
\]
The signs and the factor $h$ satisfy the chain equation
directly. At $h=1$, choose
$E=-\overline C_1-\overline C_2$,
$U=\overline C_2+\overline C_3$, and
$D=\overline C_2$.

\item \emph{Finite caps and the global complex.}
For any marked monodromy
$A=\left(\begin{smallmatrix}1&0\\t&D\end{smallmatrix}\right)$
use the chain lift
$F(A)=\left(\begin{smallmatrix}1&0\\ht&D\end{smallmatrix}\right)$
and its exterior powers. These lifts multiply exactly
and commute with contraction by $qg^*$.

\item This exterior lift represents the equivariant
torus map. The gamma cover is
$\mathbb R_g\times T^3$, and the marked linear map is
$(g,v)\mapsto(g,Dv+gt)$. At every integer value of $g$
the translation is integral, so the vertex torus map
is exactly $D$. In a single unit strip, concentrate
the loop of translations $v\mapsto v+st$ at its
terminal endpoint by an equivariant cellular
approximation relative to the two vertex tori.
The translation prism on the minimal cellular complex
of $T^3$ is exterior multiplication by $t$: its value
on each product torus is the homology class of the
literal translation sweep, and all differentials
in that minimal torus complex are zero. The terminal
vertex has deck coefficient $h$. Consequently
\[
 v\longmapsto\bigwedge D(v),\qquad
 g\wedge v\longmapsto
 g\wedge\bigwedge D(v)+ht\wedge\bigwedge D(v),
 \tag{E.C20}\hypertarget{tag-02AH}{}\label{maps:eq:E.C20}
\]
which is exactly the exterior lift above. There is
only this one strip parameter. Composing two such
maps adds the translation loops by
$t_A+D_At_B$, the same strict multiplication law
as $F(A)F(B)=F(AB)$. 
A regular
boundary with deck image $s$ is the cone of $sI-F(A)$.
This convention gives the group relation
$s\lambda s^{-1}=A\lambda$.

\item At finite cap $j$, take its adapted first fibre vector
$n_jv_j$, where
\[
 (m_1,n_1,v_1)=(3,1,(1,2,-4,0)),\qquad
 (m_2,n_2,v_2)=(4,-1,(-1,-3,3,0)).
 \tag{E.C13}\hypertarget{tag-02A8}{}\label{maps:eq:E.C13}
\]
It has gamma coordinate one and is fixed by $A_j$.
The other three coordinates form $L_0=\langle u,w,d\rangle$,
with monodromy $D_j=A_j|_{L_0}$. The cap is the
mapping torus of this three-torus, and its complex has
differential $sI-\bigwedge^rD_j$ on the swept cells.
The fibre-to-cap map is the identity on $\bigwedge L_0$;
on the adapted long-circle factor it is the integral
norm matrix $N_{m_jn_j}(s,D_j)$ determined by
\[
 (sI-D_j)N_{m_jn_j}(s,D_j)=(h-1)I,
 \qquad s=a\text{ or }b.
 \tag{E.C14}\hypertarget{tag-02A9}{}\label{maps:eq:E.C14}
\]
For positive $n_j$ the norm is
$\sum_{r=0}^{m_j-1}s^rD_j^{m_j-1-r}$, and for
$n_j=-1$ it is minus $s^{-m_j}$ times that sum.
In phase degree $r$, use
$N_{m_jn_j}(s,\bigwedge^rD_j)$, rather than the
exterior power of the norm matrix.
The boundary base-circle generator maps to the cap circle, while
the cell containing both that meridian and the adapted
long circle maps to zero. These are winding
maps in the mapping-torus coordinates. The linear
three-torus return fixes its origin; hence there is
no additional translation prism in this comparison.

\item Form the core over a bouquet of three circles with
deck images $a,b,1$ and monodromies $A_1,A_2,M_0$.
Attach the two finite complexes by \eqref{maps:eq:E.C14} and the
cusp by \eqref{maps:eq:E.C12}, using their double mapping cylinders.
Through degree three the resulting free $R$ complex
has ranks $(4,20,49,62)$. Write its group-ring entries in the integral normal form $h^r$ times
an alternating word in $a,a^2$ and $b,b^2,b^3$; exact unit
cancellations use multiplication in that normal form.

\item The final reduction is as follows. Cancelling degree-zero
and degree-one unit pairs leaves ranks $(1,2,34,62)$,
with $d_1=(a-1,b-1)$. The only nonzero columns of
$d_2$ are $F$ and $F(b-1)$, where
\[
 F=\begin{pmatrix}1+a+a^2\\h(1+b+b^2+b^3)\end{pmatrix}.
 \tag{E.C15}\hypertarget{tag-02AA}{}\label{maps:eq:E.C15}
\]
The first component of $F$ is injective on $R$:
multiplication by $a-1$ converts it to $h-1$, and
$R$ is free over the Laurent ring $A$. Thus removing
the $F$ column after subtracting its multiple gives
a free basis of $\ker d_2$ of rank33. The corresponding
row of $d_3$ is zero by the chain equation. Unit
reductions of the remaining presentation leave five
generators $(x,t,y,z,e)$. One generator $e$ occurs
in no relation. Three pure $t$ relations are
\[
 tq=0,\qquad t(1-a^{-1})=0,
                  \qquad t(b-a^{-1})=0.
 \tag{E.C16}\hypertarget{tag-02AB}{}\label{maps:eq:E.C16}
\]
They generate exactly the augmentation relations on
$t$: multiply the second by $a$ to get $t(a-1)=0$,
and subtract it from the third to get $t(b-1)=0$.
Every other pure $t$ coefficient has augmentation
zero and is consequently redundant. Reducing the
remaining $t$ coefficients by augmentation, three
relations give successively $yq=0,zq=0,xq=0$.
The final three relations are
\[
 x(1+b)+3t=0,\quad x(b+b^2)+3t=0,
                   \quad x(1+b+b^2+b^3)+6t=0.
 \tag{E.C17}\hypertarget{tag-02AC}{}\label{maps:eq:E.C17}
\]
The last two are respectively the first multiplied
by $b$, and the first multiplied by $1+b^2$.
These identities prove \eqref{maps:eq:E.C1}.

\item The invariant subgroup generated by $t$ injects in
$L$. One way to see this is to start with the exact
sequence of right $\mathbb ZC_4$ modules
\[
 0\longrightarrow\mathbb Z\longrightarrow
 \mathbb Z[C_4/C_2]\longrightarrow\mathbb Z_-\longrightarrow0,
\]
where the first map is the sum of the two cosets.
Induce to $Q$ and push out its first term by the
augmentation to the trivial module $\mathbb Z$.
The resulting presentation is \eqref{maps:eq:E.C1}, after the
change $x\mapsto x+t$ which replaces the coefficient
$3$ by $1$. Pushout preserves the displayed injection.
Restriction to $C_4$ shows nonsplitting: an invariant
correction of the lift $x$ changes this odd coefficient
by an even number. Equivalently the standard cyclic
resolution computes
$\operatorname{Ext}^1_{\mathbb ZC_4}(\mathbb Z_-,\mathbb Z)
=\mathbb Z/2$. Induction--restriction gives the same
extension group for \eqref{maps:eq:E.C3}.

\item Finally insert the fibre cycles $uw,ud,wd$ in
the core before any cancellation and carry their
coordinate columns through the same basis changes.
The five-generator presentation gives respectively
\[
 t,\qquad t(a^{-1}-a),\qquad t(a^{-1}-1).
 \tag{E.C18}\hypertarget{tag-02AD}{}\label{maps:eq:E.C18}
\]
Relations \eqref{maps:eq:E.C16} reduce these to \eqref{maps:eq:E.C4}. The external
boundary cover has homotopy type $T^3$ and these are
its three degree-two generators. This also checks
the ordinary coefficient independently: after taking
$K$ coinvariants, $2x+3t=0$, so $t$ is twice a
primitive class, exactly the coefficient $2$ in \ref{maps:src:E.2.6}.

\item For the remaining peripheral data, use the
boundary word $abc$, where the cusp generator $c$
has deck image $1$. On swept fibre chains its map to
the three core edges has blocks
\[
                 bI,\qquad F(A_1),\qquad F(A_1)F(A_2).
 \tag{E.C21}\hypertarget{tag-02AI}{}\label{maps:eq:E.C21}
\]
The identity
\[
 (aI-F(A_1))b+F(A_1)(bI-F(A_2))
       +F(A_1)F(A_2)(I-F(M_0))=(ab-1)I
\]
checks its chain equation. This is the usual prism
subdivision of the peripheral word, with
its given ribbon marking. Tracking all ten
degree-two columns through the same reductions,
and using only $t(a-1)=t(b-1)=0$ and
$xq=yq=zq=0$, gives \eqref{maps:eq:E.C19a}--\eqref{maps:eq:E.C19b}.

\end{enumerate}
\end{proof}

\begin{theorem}[The full shifted homotopy module]\hypertarget{tag-02MH}{}\label{maps:src:E.2.17}
\StatementRecord{maps:src:E.2.17}{proved}{02MH}{\SeriesRef{maps:src:E.2.16}, \SeriesRef{maps:src:E.2.13}, \eqref{maps:eq:E.P7}, \eqref{maps:eq:E.P8}, \eqref{maps:eq:E.P11}, \SeriesRef{maps:src:E.2.9}, \SeriesRef{maps:src:E.2.10}, \eqref{maps:eq:E.P9}, \eqref{maps:eq:E.P14}, \eqref{maps:eq:E.P16}, \SeriesRef{maps:src:E.2.11}}


Use the right-module convention of \ref{maps:src:E.2.16}. Write
\[
 \begin{gathered}
 R=\mathbb ZG,\qquad q_i=h_i-1,\qquad S=\mathbb Z\Gamma,\\
 \Gamma=G/\langle h_1,h_2\rangle
  =\langle a,b,c,d\mid a^3=b^4=c^3=d^4=abcd=1\rangle.
 \end{gathered}
 \tag{E.P1}\hypertarget{tag-02BJ}{}\label{maps:eq:E.P1}
\]
where $a,b$ belong to the first complement, $c,d$ to
the second, and $\mu=ab=(cd)^{-1}$. The full integral
second-homotopy module is
\[
 \pi_2(P)\simeq R/q_2R\ \oplus\ R/q_1R
                    \ \oplus\ T_L\oplus T_R.
 \tag{E.P2}\hypertarget{tag-02BS}{}\label{maps:eq:E.P2}
\]
The last two modules are inflated from the following
explicit six-generator presentations over $S$.
Put $N_a=1+a+a^2$, $N_b=1+b+b^2+b^3$, and define
$N_c,N_d$ similarly.

For $T_L$ use generators $p_i,r_i$ for $0\leq i\leq2$,
put $v_i=p_i-r_i$, and define
\[
 \begin{split}
 \Omega_L={}&v_0(1+d-d^2-d^3)
       +v_1(d^3-d)+v_2(d^2+d^3),\\
 p_iN_a={}&0,\qquad r_iN_b=0\quad(0\leq i\leq2),\\
 \Omega_L(a-1)={}&\Omega_L(b-1)
       =\Omega_L(c-1)=\Omega_L(d-1)=0.
 \end{split}
 \tag{E.P3}\hypertarget{tag-02BT}{}\label{maps:eq:E.P3}
\]
For $T_R$ use generators $B_i,C_i$, put $w_i=B_id+C_i$,
and define
\[
 \begin{split}
 \Omega_R={}&w_0a(1-b-b^2+b^3)
              +w_1a(b^2-1)-w_2a(1+b^3),\\
 B_iN_c={}&0,\qquad C_iN_d=0\quad(0\leq i\leq2),\\
 \Omega_R(a-1)={}&\Omega_R(b-1)
       =\Omega_R(c-1)=\Omega_R(d-1)=0.
 \end{split}
 \tag{E.P4}\hypertarget{tag-02BU}{}\label{maps:eq:E.P4}
\]
These are complete presentations: six generators and
ten relations in each block. In particular
\[
 q_1q_2\pi_2(P)=0,\qquad
 \ker(q_1:\pi_2(P)\to\pi_2(P))
  \cap\ker(q_2:\pi_2(P)\to\pi_2(P))=T_L\oplus T_R.
 \tag{E.P5}\hypertarget{tag-02BV}{}\label{maps:eq:E.P5}
\]
This sharpens the annihilator of \cref{maps:src:E.2.12} and determines the central-intersection summand.
Its coinvariants are $(\mathbb Z/12)^6$, so \eqref{maps:eq:E.P2}
recovers $\pi_2(P)_G=\mathbb Z^2\oplus(\mathbb Z/12)^6$.

Let $\Pi=\pi_1(Y)$ and let $\Sigma\simeq\pi_1(\Sigma_6)$
be its quotient by $z=h_1h_2$, as in \ref{maps:src:E.2.13}. Then
\[
 \pi_2(Y)\simeq(\mathbb Z\Pi)^{24}
          \oplus\operatorname{Res}^{\Gamma}_{\Sigma}(T_L\oplus T_R),
 \qquad
 \pi_2(Y)_\Pi\simeq\mathbb Z^{110}\oplus(\mathbb Z/6)^2.
 \tag{E.P6}\hypertarget{tag-02BW}{}\label{maps:eq:E.P6}
\]
The surface-group summand is inflated to $\Pi$, so
$z$ acts trivially on it. \end{theorem}


\begin{proof}
\begin{enumerate}[label={$\langle1\rangle$\arabic*.},leftmargin=*]
\item \emph{A small chain presentation with the peripheral map.}
The calculation in \ref{maps:src:E.2.16} permits the following model
through degree three for the first complement. Its
bases have ranks $(1,2,6,6)$, with degree-one basis
$(a,b)$ and degree-two basis $(r,x,t,y,z,e)$.
The first two differentials are
\[
 d_1=(a-1,b-1),\qquad
 d_2(r)=\begin{pmatrix}N_a\\h_1N_b\end{pmatrix},
 \quad d_2(x)=d_2(t)=d_2(y)=d_2(z)=d_2(e)=0.
 \tag{E.P7}\hypertarget{tag-02BX}{}\label{maps:eq:E.P7}
\]
The six degree-three columns are the six relations
\[
 t(a-1),\quad t(b-1),\quad xq_1,\quad yq_1,\quad zq_1,
                         \quad x(1+b)+3t.
 \tag{E.P8}\hypertarget{tag-02BY}{}\label{maps:eq:E.P8}
\]
The degree-two norm $N_a$ in \eqref{maps:eq:E.P7} is written with
the same expression as in \eqref{maps:eq:E.P3}, but here $a^3=h_1$,
not $1$. The second complement is the identical
model with $(a,b,h_1)$ replaced by $(c,d,h_2)$.

\item The ordinary one-relator complex $C_K$ is the subcomplex
on degrees zero and one and the cell $r$. It is an
aspherical model: it is the double mapping cylinder
of the degree-three and degree-four maps of circles,
and its universal cover is a tree of contractible
lines and strips. Equivalently it is the standard
Seifert spine of the torus-knot exterior. Projection
of \eqref{maps:eq:E.P7}--\eqref{maps:eq:E.P8} onto $C_K$ retains the peripheral map. In degree one this map is
\[
 g\longmapsto(N_a,0),\qquad
 u,w,d\longmapsto0,\qquad
 \Theta\longmapsto(b,1).
 \tag{E.P9}\hypertarget{tag-02BZ}{}\label{maps:eq:E.P9}
\]
Its degree-two $r$ coefficient is $b-1$ on the swept
cell $g\Theta$ and zero on the other nine cells.
Its remaining components are exactly \eqref{maps:eq:E.C19a}--\eqref{maps:eq:E.C19b}.
\item Here and below the small complex is only a
presentation through the degree needed for $H_2$.
Replacing redundant degree-three columns by the six
generators of their image in \eqref{maps:eq:E.P8} changes no
degree-two cycle or boundary. Reducing a peripheral
degree-two column modulo that image changes it by
a chain homotopy with values in degree three.
Products with a torus and the mapping cones below
therefore retain the same degree-two homology.
For the comparison with an aspherical target one
may first keep the original, full complex: its
projection to $C_K$ extends by zero in degrees at
least three, since the Fox column in \eqref{maps:eq:E.P7} is
injective and hence the $r$ row of $d_3$ is zero.
Only after taking the kernel is its image in degree
two shortened as in \eqref{maps:eq:E.P8}. 

\item \emph{The exact kernel complex for the splice.}
Let $T(s)$ be the four-torus chain complex whose
first circle has deck image $s$ and whose other
three circles have trivial image. Its differential
contracts that first circle and multiplies by $s-1$.
The two vertex complexes for the splice are
\[
 V_1=C_{E_1}\otimes T(h_2),\qquad
 V_2=C_{E_2}\otimes T(h_1).
 \tag{E.P10}\hypertarget{tag-02BK}{}\label{maps:eq:E.P10}
\]
We have reversed the two product factors in the
second vertex only for notation; the Koszul signs
below account for that reversal. The boundary
complex $B$ is the nine-torus complex with ordered
coordinates
\[
 (g_1,u_1,w_1,d_1,g_2,u_2,w_2,d_2,\Theta),
 \quad\text{deck images }(h_1,1,1,1,h_2,1,1,1,\mu).
 \tag{E.P11}\hypertarget{tag-02BL}{}\label{maps:eq:E.P11}
\]
Its differential is the usual exterior contraction
by those deck images minus one.

\item The source signs are as follows. Write
a boundary product cell as $\alpha\wedge\beta$
with an optional final $\Theta$, where $\alpha$
uses the first four coordinates and $\beta$ the
second four. Put $\epsilon=1$ when $\Theta$ occurs
and $0$ otherwise. Its first vertex map applies
the \ref{maps:src:E.2.16} peripheral map to $(\alpha,\Theta)$,
keeps $\beta$, and has sign
$(-1)^{\epsilon(|\alpha|+|\beta|)}$.
The second map applies that peripheral map to
$(\beta,\Theta_2)$, keeps $\alpha$, and has sign
$(-1)^{|\alpha|(|\beta|+\epsilon)+\epsilon|\beta|}$.
If $\epsilon=1$, multiply its coefficient on the
right by $-\mu$. This is the exact cellular map
of $\Theta_1\mapsto\Theta_2^{-1}$, since
$\mu_2=cd=\mu^{-1}$. The suspended cells in (E.C19)
have the interval first; these signs convert from
the ordered exterior cells in \eqref{maps:eq:E.P11}.
Denote the two maps by $j_1,j_2$.

\item The chain complex of the splice is the cone of
$(j_1,-j_2):B\to V_1\oplus V_2$.
Project each $C_E$ to $C_K$ and each $T(s)$ to its
first circle, and project $B$ to its three coordinates
$g_1,g_2,\Theta$. The result is the complex of the
aspherical graph of spaces
\[
 (K(K_1,1)\times S^1)\ \cup_{T^3}\
                         (K(K_2,1)\times S^1).
 \tag{E.P12}\hypertarget{tag-02BM}{}\label{maps:eq:E.P12}
\]
It is the model $N$ of \ref{maps:src:E.2.9}. Both peripheral maps
on groups are injective by \ref{maps:src:E.2.10}; their common
group is $\langle h_1,h_2,\mu\rangle=\mathbb Z^3$.
Consequently its universal cover is a Bass--Serre
tree of contractible vertex and edge spaces and
is contractible. Its chain complex has ranks
$(2,7,9,5,1)$, including its degree-four cell.

\item The indicated projection is a degreewise split
surjection of the full chain complexes. The
peripheral Fox coefficient in \eqref{maps:eq:E.P9} makes the
boundary square commute; higher degrees commute
as well by the product construction and the
zero higher groups of $C_K$. Let $D$ be its kernel.
The target has zero positive-degree homology, so
the long exact sequence gives
$H_2(D)\simeq H_2(\widetilde P)=\pi_2(P)$.
Using the small models for the relevant images,
the ranks of $D$ through degree three are
\[
                         (0,6,40,123).
 \tag{E.P13}\hypertarget{tag-02BN}{}\label{maps:eq:E.P13}
\]
Its six degree-one generators are the unused
regular-factor circles at the two vertices.
The six shifted boundary circles $u_i,w_i,d_i$
map to them by unit entries, because each circle
maps to zero in its own complement and identically
to the opposite regular torus. Cancel these six
unit pairs. This leaves a presentation of $H_2(D)$
with34 generators and123 relations.

\item Here is also an explicit block definition of that
matrix. In each $V_i$ keep the cells killed by
projection to $C_{K_i}\otimes C(S^1)$; their ranks
in degrees two and three are17 and45. In $B_2$
keep the33 cells containing at least one of
$u_1,w_1,d_1,u_2,w_2,d_2$. In those ordered bases
the matrix is
\[
 \mathcal M=
 \begin{pmatrix}
 (d_{V_1})_3&0&(j_1)_2\\
 0&(d_{V_2})_3&-(j_2)_2
 \end{pmatrix}:
 R^{45}\oplus R^{45}\oplus R^{33}
             \longrightarrow R^{17}\oplus R^{17}.
 \tag{E.P14}\hypertarget{tag-02BO}{}\label{maps:eq:E.P14}
\]
Every entry is specified by \eqref{maps:eq:E.P7}--\eqref{maps:eq:E.P11} and
(E.C19); \eqref{maps:eq:E.P14} itself is a complete finite
presentation of the desired module. The formula
includes the derived peripheral comparison.


\item \emph{Reduction to the compact presentation.}
The peripheral chain equations and $d_2d_3=0$ hold by multiplication
of their group-ring matrices. Cancelling the indicated unit pairs reduces
\eqref{maps:eq:E.P14} to $16$ generators and $52$ nonzero columns.
The calculation takes place in the amalgam
\[
 G=(K_1\times\mathbb Z)\mathop{*}_{\mathbb Z^3}
                                      (\mathbb Z\times K_2),
 \tag{E.P15}\hypertarget{tag-02BP}{}\label{maps:eq:E.P15}
\]
whose normal form is specified as follows: the torus-knot factor uses $h^r$ times an
alternating word in $a,a^2$ and $b,b^2,b^3$.
Modulo the peripheral subgroup, strip initial
$ab$ and $b^3a^2=(ab)^{-1}$, allowing their last
letter to be a partial power; choose $b^3$ in
the sole length-one tie with $a$. Inverting these
left coset representatives gives right coset
representatives. The usual alternating amalgam
normal form, followed by a final element of
$\mathbb Z^3$, is unique. These word operations track both central exponents.

\item In the resulting16-generator presentation two
generators occur only in their respective single
relations $q_2e_1=0$ and $q_1e_2=0$. They give the
first two direct summands of \eqref{maps:eq:E.P2}. For the other
14 generators, two generators $t_0,t_1$ have
trivial action of $a,b,c,d$. There are three
generators initially killed by $q_2$ whose pure
relations include $h_1N_b$; the identity
$h_1N_b(b-1)=1-h_1$ kills them by $q_1$ as well.
There are three initially killed by $q_1$ with
pure $N_c$ relations, and three more with pure
$h_2N_d$ relations. They are killed by $q_2$.
The remaining generators already have both
central relations. Thus the whole14-generator
summand is a module over $S$, without imposing
an additional quotient. The pure norm columns
are exact differences of the recorded columns
and are checked before setting either central
element equal to1.

\item For a readable final reduction, the14 generators
separate into two seven-generator blocks. The
first has generators $U_0,U_1,U_2,X,Y,Z,t_1$,
with $t_1$ invariant, $U_iN_b=0$, and
\[
 \begin{split}
 (U_0+X-Yd+Zd)N_a&=0,\\
 (U_1-Y(1+d)+Z)N_a&=0,\\
 (U_2+X+2t_1+Z(d-1))N_a&=0,\\
 X(1+d)+3t_1&=0.
 \end{split}
 \tag{E.P16}\hypertarget{tag-02BQ}{}\label{maps:eq:E.P16}
\]
The second has generators $B_i,C_i,t_0$, with
$B_iN_c=C_iN_d=0$, $t_0$ invariant, and the single
additional relation $t_0=\Omega_R$. These formulas
follow by collecting the displayed16-generator
columns.

\item In the first block put
\label{layout:V:1}
\[
\begin{gathered}
v_0=X-Yd+Zd,\quad v_1=-Y(1+d)+Z,
 \\
\quad v_2=X+2t_1+Z(d-1),\qquad
 r_i=U_i,\quad p_i=U_i+v_i.
\end{gathered}
\tag{E.P17}\hypertarget{tag-02BR}{}\label{maps:eq:E.P17}
\]
The linear part of this change on $(X,Y,Z)$ is
\[
 \begin{pmatrix}1&0&1\\-d&-1-d&0\\d&1&d-1\end{pmatrix},
 \quad\det=1,
\]
over $\mathbb Z[d,d^{-1}]$. It is therefore an
integral group-ring basis change. Its inverse gives
\[
 X=v_0(1-d^2)+v_1(d^2-d)+(v_2-2t_1)d^2.
\]
Substitute in the last relation of \eqref{maps:eq:E.P16}, using
the invariance of $t_1$. The result is
$t_1=\Omega_L$. Eliminate $t_1$ with this unit
coefficient. Its four invariance relations become
the four relations in \eqref{maps:eq:E.P3}. Eliminating $t_0$
in the second block gives \eqref{maps:eq:E.P4}. This proves the
compact presentation in \eqref{maps:eq:E.P2}--\eqref{maps:eq:E.P4}, including
the integral central-intersection information.

\item For the two cyclic axis summands, the other
central element is a non-zero-divisor: their
group rings are free over its Laurent polynomial
ring. Thus neither contributes to the common
kernel of $q_1,q_2$. This proves \eqref{maps:eq:E.P5}. On
coinvariants, the four invariance relations in
each block vanish, while the norm relations
become three copies of multiplication by3 and
three copies of multiplication by4. Each block
therefore has coinvariants $(\mathbb Z/12)^3$.
This is an independent recovery of \ref{maps:src:E.2.11}.

\item \emph{Restriction to the circle fibre.}
The universal covers of $P$ and $Y\times\mathbb R$
are $\Pi$-equivariantly homotopy equivalent by
\ref{maps:src:E.2.13}. The map $\Pi\to G/\langle h_i\rangle$
is injective, since $\chi(h_i)=\pm12$, and has
index12. Hence each of the first two group-ring
summands in \eqref{maps:eq:E.P2} is free of rank12 over
$\mathbb Z\Pi$. The last two factor through the
index12 surface subgroup $\Sigma<\Gamma$ and
have trivial $z$ action. This proves the module
formula in \eqref{maps:eq:E.P6}.

\item For its exact coinvariants use the regular
12-sheet permutation representation with
exponents $(4,-3,-4,3)$ on $(a,b,c,d)$.
Evaluating the compact12-by20 presentation gives
an integral144-by240 matrix. Its Smith form has
56 unit entries and two entries6; its free
cokernel rank is86. Adding the24 free
coinvariants proves the last formula in \eqref{maps:eq:E.P6}.

\end{enumerate}
\end{proof}


\begin{definition}[Whitehead quadratic group]\label{def:V:whitehead-quadratic}
For an abelian group $A$, let $\Gamma(A)$ be Whitehead's universal
quadratic group, with quadratic generator $q(a)$ and symmetric
bilinear cross effect
$[a,b]=q(a+b)-q(a)-q(b)$.  In particular $[a,a]=2q(a)$.
\end{definition}

\begin{proposition}[The cusp quadratic attachment]\hypertarget{tag-02MI}{}\label{maps:src:E.2.18}
\StatementRecord{maps:src:E.2.18}{proved}{02MI}{\SeriesRef{maps:src:E.2.16}, \eqref{maps:eq:E.P6}, \SeriesRef{maps:src:D.1.9}, \SeriesRef{maps:src:E.2.12}, \eqref{maps:eq:E.C12}, \eqref{maps:eq:E.Q4}, \eqref{maps:eq:E.Q5}, \eqref{maps:eq:E.Q3}, \eqref{maps:eq:E.Q7}, \eqref{maps:eq:E.Q6}, \eqref{maps:eq:E.C3}, \eqref{maps:eq:E.Q8}, \SeriesRef{maps:src:E.2.17}, \eqref{maps:eq:E.P2}, \eqref{maps:eq:E.P3}, \eqref{maps:eq:E.P4}, \eqref{maps:eq:E.Q11}, \eqref{maps:eq:E.Q10}, \eqref{maps:eq:E.Q12}; \cite[Theorems 1 and 5]{mapsWhitehead50}}


Keep the right-module presentation of \ref{maps:src:E.2.16} and write
\[
 R=\mathbb ZK,\qquad M=\pi_2(E)=eR\oplus T,
 \qquad T=\overline R^{\,2}\oplus L.
\]
There is a lift of the normalized cusp component
$S=\mathbb {CP}^2\#3\overline{\mathbb {CP}}^{\,2}$ to the
universal cover $\widetilde E$ whose fundamental class has image
\[
 \pm\bigl([e(h-1),t]-12q(t)\bigr)
       \quad\hbox{under}\quad
 H_4(\widetilde E;\mathbb Z)\longrightarrow\Gamma(M).
 \tag{E.Q1}\hypertarget{tag-02C0}{}\label{maps:eq:E.Q1}
\]
The sign is solely the orientation chosen on the normalized
component.  Its mixed component in
\[
 \Gamma(M)=\Gamma(eR)\oplus(eR\otimes_{\mathbb Z}T)
                                      \oplus\Gamma(T)
 \tag{E.Q2}\hypertarget{tag-02C4}{}\label{maps:eq:E.Q2}
\]
is nonzero.  Consequently $E$ is not homotopy equivalent to
$E_0\#(S^2\times S^4)$ for any connected compact six-manifold
$E_0$, even without requiring the equivalence to preserve the
boundary.  Thus the free $R$ summand of its second homotopy group
does not come from a sphere-product connected summand.
Likewise the nine-dimensional fibre $Y$ is not
homotopy equivalent to $Y_0\#(S^2\times S^7)$ for any
connected closed nine-manifold $Y_0$, despite the twenty-four
free $\mathbb Z\Pi$ summands in \eqref{maps:eq:E.P6}.

\end{proposition}


\begin{proof}
\begin{enumerate}[label={$\langle1\rangle$\arabic*.},leftmargin=*]
\item We use the natural, deck-equivariant part of Whitehead's exact
sequence
\[
 H_4(\widetilde E)\xrightarrow{\omega_E}\Gamma(\pi_2(E))
       \longrightarrow\pi_3(E)\longrightarrow
 H_3(\widetilde E)\longrightarrow0.
 \tag{E.Q3}\hypertarget{tag-02C5}{}\label{maps:eq:E.Q3}
\]
The exactness and identification of its quadratic term are
\cite[Theorems 1 and 5]{mapsWhitehead50}; the direct-sum rule
\eqref{maps:eq:E.Q2} is proved there in Section~2.  These statements apply
to the simply connected universal cover, with no finiteness
requirement.  Naturality specifies the element in \eqref{maps:eq:E.Q1}.

\item For the normalization calculation take the negatives of
the complex-oriented Picard homology generators and denote
them by $(H,E_1,E_2,E_3)$.  Their intersection matrix is
still $\operatorname{diag}(1,-1,-1,-1)$.  This common reversal matches the phase-cell orientations of the cusp basis: the toric
line with class $-(uw+6\gamma\delta)$ has complex curve
degrees $(-6,0,1)$, whereas its values on the corresponding
phase classes are $(6,0,-1)$; the support function $rs$
similarly fixes the middle sign.  The exact wall-degree
calculation in \SeriesRef{maps:src:D.1.9} gives this comparison.  With these
orientations the six conductor classes are
\[
 E_1,\ H-E_1-E_2,\ E_2,\ H-E_2-E_3,
                                      E_3,\ H-E_1-E_3.
\]
The first opposite pair is self-glued in the surviving $u$
direction.  In the gamma cover its relation is therefore
$E_1=H-E_2-E_3$.  The other two pairs connect adjacent gamma
translates.  Their relations, after eliminating $H$, are
$(h-1)E_2=(h-1)E_3=0$.  Thus the image of normalization
homology is, with these curve generators,
\[
 A E_1\oplus(A/(h-1))E_2\oplus(A/(h-1))E_3,
 \qquad A=\mathbb Z[h,h^{-1}],\qquad H=E_1+E_2+E_3.
 \tag{E.Q4}\hypertarget{tag-02C6}{}\label{maps:eq:E.Q4}
\]
This also follows directly from the transpose of the matrix
$B(h)$ of \ref{maps:src:E.2.12}, with the intersection signs of the Picard basis.  It is the whole $H_2$ of the gamma cover:
the degree-zero normalization row is the cellular complex
of an infinite cylinder, whose second homology is zero.

\item We fix the comparison with the cusp basis of \eqref{maps:eq:E.C12}.  In the ordinary
central fibre the first conductor class $\overline C_1$ is the
self-$u$ pair; the other two are the images of $E_2,E_3$,
in either order.  The chosen cusp basis is
\[
 E=-\overline C_1-\overline C_2,\qquad
 U=\overline C_2+\overline C_3,\qquad D=\overline C_2.
\]
The two torsion coordinates in \eqref{maps:eq:E.Q4} are determined by
this ordinary comparison, because each is annihilated by
$h-1$.  The normalization map is onto the whole group in
\eqref{maps:eq:E.Q4}, and $E_2,E_3$ generate its entire $A$-torsion
submodule.  Consequently the coefficient of $E_1$ in the
free quotient generates $A$ and hence is a unit $\pm h^n$.  The ordinary comparison
fixes its sign.  Translating the entire chosen normalization
component by $h^{-n}$ makes this coefficient $-1$ and leaves
both torsion coordinates fixed.  The common Picard sign
convention above changes every generator simultaneously.
The resulting
normalization map in the basis used for the full peripheral
comparison is
\[
 E_1\longmapsto-E-D,\qquad
 E_2\longmapsto D,\qquad E_3\longmapsto U-D,
 \qquad H\longmapsto-E+U-D.
 \tag{E.Q5}\hypertarget{tag-02C7}{}\label{maps:eq:E.Q5}
\]
Exchanging $E_2,E_3$ has no effect below.  \item The intersection form of $S$ is
$\operatorname{diag}(1,-1,-1,-1)$.  Its usual four-cell
description, or the Hopf invariant of its two-dimensional
classes, therefore gives
\[
 \omega_S[S]=q(H)-q(E_1)-q(E_2)-q(E_3).
\]
Substitute $H=E_1+E_2+E_3$ and then \eqref{maps:eq:E.Q5}.  Bilinearity
of the cross effect gives
\[
 \omega_S[S]\longmapsto
 [E_1,E_2]+[E_1,E_3]+[E_2,E_3]
                      =-[E,U]-2q(D).
 \tag{E.Q6}\hypertarget{tag-02C8}{}\label{maps:eq:E.Q6}
\]
The normalization map $S\to E$ lifts to $\widetilde E$
because $S$ is simply connected.  Naturality of \eqref{maps:eq:E.Q3}
shows that the last expression is the image of its four-dimensional fundamental class.  \item The global cancellation also tracks the three cusp two-cycles in addition to the peripheral chains.  It gives, in the generators $(x,t,y,z,e)$ of
\ref{maps:src:E.2.16},
\[
 E\longmapsto(1-h)e+6t,\qquad
 U\longmapsto t,\qquad D\longmapsto0.
 \tag{E.Q7}\hypertarget{tag-02C9}{}\label{maps:eq:E.Q7}
\]
For clarity, before the last relation reductions its
$t$ coefficient for $E$ is
$-3h+4ha+4ha^2+h^2$, whose augmentation is six;
its free coefficient is $1-h$.  Its coefficient for $D$
is $2h^{-1}a^2-1-a$, which has augmentation zero and hence
annihilates the invariant generator $t$.  The other coordinates vanish. These identities follow by applying
the integral group-ring cancellations of \ref{maps:src:E.2.16} to
the displayed cusp-cycle columns.  Applying $\Gamma$ to \eqref{maps:eq:E.Q7} in \eqref{maps:eq:E.Q6} gives
\[
 -[(1-h)e+6t,t]
                =[e(h-1),t]-12q(t),
\]
as asserted.  In ordinary coinvariants $t=2v$ and
$E=12v$ for a primitive class $v$.  The quadratic image
is then $-48q(v)$, equal to the degree-two cup square evaluated
on the normalization.

\item We finish with the connected-sum obstruction, allowing
every possible free summand.  Multiplication by $h-1$
is injective on $R$: a finitely supported group-ring
element constant along an infinite $h$ orbit is zero.
It annihilates $T$, so
\[
 \operatorname{Ann}_M(h-1)=T,
                       \qquad\operatorname{Hom}_R(T,R)=0.
 \tag{E.Q8}\hypertarget{tag-02CA}{}\label{maps:eq:E.Q8}
\]
Suppose that a copy $fR$ is a direct summand of $M$.
Write $f=eu+t_0$ and let a retraction send $e$ to $v$.
It vanishes on $T$, and hence $vu=1$.

\item The ring $\mathbb ZK$ is directly finite.  One elementary
verification uses the index-twelve subgroup of $K$:
its quotient by the central $\mathbb Z$ is the free group
of the punctured genus-three surface, so it is
$F_6\times\mathbb Z$.  This subgroup, and hence $K$, is
residually finite.  In every finite-group quotient the
regular representation is finite dimensional over
$\mathbb Q$, so $vu=1$ implies $uv=1$.  Residual
finiteness separates the finite support of $uv-1$ and
proves the same equality in $R$.  Thus $u$ is a unit,
the kernel of the retraction is exactly $T$, and every
free splitting is obtained by replacing $e$ by
$eu+t_0$.

\item Under such a change the mixed component of \eqref{maps:eq:E.Q1} is
$f u^{-1}(h-1)\otimes t$.  It is nonzero: $t$ has
infinite order by \eqref{maps:eq:E.C3}, and $u^{-1}(h-1)$ is nonzero.
The central generator may also be reversed by a group
automorphism, which multiplies $h-1$ by a unit and
does not affect the argument.

\item If $E\simeq E_0\#(S^2\times S^4)$, its universal cover
has, through dimension four, the homotopy type of
$\widetilde E_0$ wedged with one copy of $S^2\vee S^4$
for each deck element.  This follows by lifting the
connected-sum balls; deleting or regluing their boundary
five-spheres does not alter the four-skeleton.  The new
two-spheres generate a free $R$ summand.  The four-spheres
have trivial attaching maps, and every remaining
four-dimensional Whitehead image lies in
$\Gamma(\pi_2(E_0))$.  There can therefore be no mixed
component between this free summand and its complement.
This contradicts \eqref{maps:eq:E.Q1}--\eqref{maps:eq:E.Q8} for every choice of the
summand, and proves the stated obstruction.

\item For the assertion about $Y$, first take the maps from the two complements to the splice.  Carry the four cycles $(t_0,e_0,t_1,e_1)$ through the unit
cancellations of \ref{maps:src:E.2.17}. In the resulting
sixteen-generator presentation they are four unit vectors.  The seven-generator block changes therefore give
\[
 t_0=\Omega_R,\qquad t_1=\Omega_L,
 \qquad e_0\in R/q_2R,\quad e_1\in R/q_1R
 \tag{E.Q9}\hypertarget{tag-02CB}{}\label{maps:eq:E.Q9}
\]
with exactly the axes in \eqref{maps:eq:E.P2}.  
Both $t_i$ are invariant under $G$: the four displayed
relations for each $\Omega$ in \eqref{maps:eq:E.P3}--\eqref{maps:eq:E.P4} annihilate
it by $a-1,b-1,c-1,d-1$.  On $\Gamma$-coinvariants
their respective images are
\[
 -2(B_2+C_2),\qquad 2(p_2-r_2),
 \tag{E.Q10}\hypertarget{tag-02C1}{}\label{maps:eq:E.Q10}
\]
and each has order exactly six: the relevant two coordinates
have independent orders three and four in \eqref{maps:eq:E.P3}--\eqref{maps:eq:E.P4}.
In particular an integer annihilating either $t_i$ in the
full module must be divisible by six.

\item Put $A_\Pi=\mathbb Z\Pi$ and $z=h_1h_2$.  Choose the twelve
coset representatives giving the free restriction of each
axis in \eqref{maps:eq:E.P6}, and denote the resulting twenty-four generators
by $f_j$.  The normalized cusp surfaces and all their chosen
deck translates are maps to $\widetilde P$.  Transfer
them to $\widetilde Y$ using the $\Pi$-equivariant homotopy
equivalence $\widetilde P\simeq\widetilde Y\times\mathbb R$.
On the first axis $h_2$ is the identity, so multiplication
by $h_1-1$ equals multiplication by $z-1$; the analogous
statement holds on the second axis.  Naturality of \eqref{maps:eq:E.Q3}
therefore supplies, for every $j$, an element in the image
of the Whitehead map of $\widetilde Y$ of the form
\[
 \pm\bigl([f_j(z-1),t_{i(j)}]-12q(t_{i(j)})\bigr),
 \tag{E.Q11}\hypertarget{tag-02C2}{}\label{maps:eq:E.Q11}
\]
where $i(j)=0$ on the first twelve axes and $1$ on the
remaining twelve.  Invariance of $t_i$ under $G$ removes
any extra coset action in this formula.

\item Write $\pi_2(Y)=F\oplus T_Y$, with $F=A_\Pi^{24}$ and
$(z-1)T_Y=0$ as in \eqref{maps:eq:E.P6}.  Suppose a free copy $fA_\Pi$
is a direct summand, with a retraction $r:\pi_2(Y)\to A_\Pi$
satisfying $r(f)=1$.  Multiplication by $z-1$ is injective
on $A_\Pi$, so $r(T_Y)=0$.  Write $r(f_j)=v_j$ and
$f=\sum_j f_j u_j+t$; then
\[
                         \sum_j v_j u_j=1.
 \tag{E.Q12}\hypertarget{tag-02C3}{}\label{maps:eq:E.Q12}
\]
In the splitting $fA_\Pi\oplus\ker r$, the mixed component
of \eqref{maps:eq:E.Q11} is $f v_j(z-1)\otimes t_{i(j)}$.  If all these
mixed components were zero, every integer coefficient of
$v_j(z-1)$ would be divisible by six, by \eqref{maps:eq:E.Q10}.  The
group-ring basis makes this implication immediate even if
$t_i$ has larger or infinite order.  Multiplication by
$z-1$ is also injective on $(\mathbb Z/6)[\Pi]$: along an
infinite $z$ orbit a finitely supported coefficient function
with zero difference is zero.  Hence every $v_j$ belongs
to $6A_\Pi$, contradicting \eqref{maps:eq:E.Q12}.

\item Finally, a connected sum with $S^2\times S^7$ has, through
dimension four in its universal cover, the original space
wedged with a free deck orbit of two-spheres.  All images
of its four-dimensional Whitehead map consequently have
zero mixed component with that free summand.  The preceding
argument rules this out for every free summand of the module.  This proves the assertion for $Y$.
\end{enumerate}
\end{proof}




\begin{proposition}[The relative smooth cusp model]\hypertarget{tag-02MJ}{}\label{maps:src:E.2.19}
\StatementRecord{maps:src:E.2.19}{proved}{02MJ}{\SeriesRef{maps:src:E.2.21}, \SeriesRef{maps:src:E.2.16}, \eqref{maps:eq:E.T3}, \eqref{maps:eq:E.T4}, \eqref{maps:eq:E.T8}}


For the original marked cusp, let
\[
 Z(q)=\frac{\log q}{2\pi i}B_0+C(q),\qquad
 B_0=\begin{pmatrix}0&1\\-1&0\end{pmatrix},
\]
where $C(q)$ is holomorphic and bounded on a disc about zero.
For a sufficiently small closed disc the six-dimensional
cusp tube is orientation-preservingly diffeomorphic, relative to
its full marked torus-bundle boundary, to the unitless toric
quotient defined by $C=0$.  The latter has the following finite
normal-crossings plumbing description.

Start with the canonical complex line bundle of
$S=\mathbb {CP}^2\#3\overline{\mathbb {CP}}^{\,2}$ and its six
toric curves in cyclic order
\[
 \begin{aligned}
 C_0&=E_1,& C_1&=H-E_1-E_2,& C_2&=E_2,\\
 C_3&=H-E_2-E_3,& C_4&=E_3,& C_5&=H-E_1-E_3.
 \end{aligned}
 \tag{E.T1}\hypertarget{tag-02DV}{}\label{maps:eq:E.T1}
\]
Pair $C_i$ with $C_{i+3}$ for $i=0,1,2$.  Both the surface-normal
line and the canonical line restricted to every $C_i$ have degree
$-1$.  In the resulting $\mathcal O(-1)\oplus\mathcal O(-1)$
coordinates $(z,a,b)$, with $z=0$ at the endpoint meeting
$C_{i+1}$ and $b$ the surface-normal coordinate, the pairing is
\[
             (z,a,b)\longmapsto(z^{-1},zb,za).
 \tag{E.T2}\hypertarget{tag-02DW}{}\label{maps:eq:E.T2}
\]
The function giving the cusp parameter is $q=zab$ in this
coordinate chart.  The full coordinate map \eqref{maps:eq:E.T2} specifies the normal
framing of each plumbing.
The two triple points are the endpoint orbits
$\{0,2,4\}$ and $\{1,3,5\}$, where endpoint $i$ means
$C_i\cap C_{i+1}$.  Ordered local complex triples undergo a
cyclic permutation around each orbit, with threefold composition
the identity.  The plumbing is formed with collars and rounded
corners; its distinguished outer boundary is $|q|=\eta$.

\end{proposition}


\begin{proof}
\begin{enumerate}[label={$\langle1\rangle$\arabic*.},leftmargin=*]
\item We first construct the comparison by a smooth family.  Set
$C_s(q)=(1-s)C(q)$ for $0\leq s\leq1$.  The affine toric cover
is fixed.  On its dense torus the deck translation for
$\lambda\in\mathbb Z^2$ becomes
\[
 (x,q,s)\longmapsto
 \bigl(q^{B_0\lambda}
       \exp(2\pi i C_s(q)\lambda)x,\ q,\ s\bigr).
 \tag{E.T3}\hypertarget{tag-02DX}{}\label{maps:eq:E.T3}
\]
The exponential is componentwise.  The powers of $q$ permute
the toric charts by translation of the $A_2$ triangulation,
while the exponential factors are nowhere-zero analytic units
on every chart.  Thus \eqref{maps:eq:E.T3} defines a smooth family of
holomorphic deck actions on a fixed smooth toric covering
space.  On a central stratum the lattice translation permutes
its bounded fan cells without fixed points.  Off the central
fibre, freeness and properness are exactly the nondegeneracy
of the real period lattice.

\item These assertions hold uniformly after shrinking $\eta$.
Indeed put $T=-\log|q|/(2\pi)$, and let $K$ bound
$\|\operatorname{Im}C_s(q)\|$ on the compact parameter interval
and closed disc.  For $T>K+1$ the matrix
$TB_0+\operatorname{Im}C_s(q)$ is invertible, since $B_0$ is
orthogonal.  A fundamental parallelogram has logarithmic
position
$B_0a+T^{-1}\operatorname{Im}C_s(q)a$ with
$a\in[0,1]^2$; it meets a uniformly bounded number of toric
charts.  More explicitly choose a triangle containing the
unperturbed point $B_0a$ and let $P$ be its unimodular vertex
matrix.  Its barycentric coordinates $t_j$ are nonnegative,
and the corresponding chart coordinates satisfy
$\log|z_j|=t_j\log|q|+
O(2\pi\|P^{-1}\|K\sqrt2)$.
There are only finitely many such triangles, so with
$c=2\pi\sqrt2\max\|P^{-1}\|$ every chosen chart coordinate
has modulus at most $\exp(cK)$, uniformly in $q,s$ and
all compact phases.  Their closed polydiscs have compact
closure in the affine charts.  These finitely many compact
chart images prove
compactness of the quotient over $[0,1]$ and $|q|\leq\eta$.
The central charts and their incidence relations are fixed.
Projection of this quotient to $s$ is consequently a proper
smooth submersion, including on every normal-crossings stratum.
Its only boundary is $|q|=\eta$ together with the parameter
end faces, and its boundary restriction is also a submersion.

\item On $|q|=\eta$ prescribe the marked real coordinates
$x=\exp(2\pi i(Z_s(q)a+b))$, $(a,b)\in T^4$.
They give a smooth product trivialization in $s$ and precisely
the prescribed mapping-torus marking.  Extend its horizontal
vector field over a collar, and then over the compact family.
Local toric product charts provide extensions tangent to
every central stratum; a partition of unity preserves that
property and the lift of $\partial_s$.  Integrating this
field proves the asserted relative diffeomorphism.  Its
orientation is the complex orientation on each fibre.
This is relative Ehresmann transport in a smooth
family.

\item It remains to identify the endpoint $C=0$ with the framed
plumbing data.  In the lattice $\mathbb Z^2$ use cyclic rays
\[
 (1,0),\ (0,1),\ (-1,1),\ (-1,0),\ (0,-1),\ (1,-1).
 \tag{E.T4}\hypertarget{tag-02DY}{}\label{maps:eq:E.T4}
\]
The star of the vertex zero in the $A_2$ triangulation consists
of the six triangles spanned by zero and consecutive rays.
The corresponding three-dimensional cones have primitive rays
$(v,1)$ and determinant one.  Their common divisor at the
central ray is the toric surface $S$.  A principal central
fibre has normal line
\[
             \mathcal O_S(-C_0-\cdots-C_5)=K_S.
 \tag{E.T5}\hypertarget{tag-02DZ}{}\label{maps:eq:E.T5}
\]
The cyclic surface fan in \eqref{maps:eq:E.T4} also proves this identification
directly: it is the fan of the three torus-fixed blow-ups of
$\mathbb {CP}^2$, and each neighboring-ray relation gives
$C_i^2=-1$.  In the Picard basis \eqref{maps:eq:E.T1},
$K_S=-3H+E_1+E_2+E_3$, so $K_S\cdot C_i=-1$.

\item Near a double curve the two adjacent unimodular charts give
the total space of $\mathcal O(-1)\oplus\mathcal O(-1)$.
Its two base charts have coordinates
\[
 (z,a,b),\qquad
 (w,A,B)=(z^{-1},za,zb),\qquad q=zab=wAB.
 \tag{E.T6}\hypertarget{tag-02E0}{}\label{maps:eq:E.T6}
\]
The period translation taking a neighboring component back
to the chosen component exchanges the two normal directions.
It reverses the chosen endpoint coordinate on the opposite
toric curve.  In the first chart this is exactly \eqref{maps:eq:E.T2};
in the target infinity chart it is simply $(z,b,a)$, so it
extends smoothly over the endpoint.  The map is involutive,
preserves $q$, and its complex Jacobian determinant is $1$.
These equalities specify the normal-disk attachment, including
its sign and angular twisting.

\item The endpoint $C_i\cap C_{i+1}$ is sent to
$C_{i+2}\cap C_{i+3}$.  At the first endpoint write the
ordered ambient coordinates as the surface-normal coordinate,
the $C_i$-normal coordinate and the $C_{i+1}$-normal coordinate.
The next endpoint has their order cyclically permuted.
The threefold composition is the identity on the ordered coordinates.  Thus the two endpoint orbits give
ordinary triple charts $q=z_0z_1z_2$, with no additional
clutching loop at either triple point.

\item Equivalently, for a triangle with vertex matrix $P$ whose
columns are $(v_i,1)$, translation by $n\in\mathbb Z^2$
replaces it by $L_nP$, where
\[
 L_n=\begin{pmatrix}1&0&n_1\\0&1&n_2\\0&0&1\end{pmatrix}.
 \tag{E.T7}\hypertarget{tag-02E1}{}\label{maps:eq:E.T7}
\]
It acts as the identity in the corresponding translated
unimodular coordinates, since
$(L_nP)^{-1}L_nP=I$.  This verifies the preceding plumbing
and triple identifications in the toric charts.
More explicitly, order the star matrix as
$P_i=((0,1),(r_i,1),(r_{i+1},1))$ and reverse its
columns when using the coordinate order $(z,a,b)$.
For every $i$ the matrix of translation by $-r_i$ from
this chart to the opposite following-endpoint chart is
\[
 \begin{pmatrix}-1&0&0\\1&0&1\\1&1&0\end{pmatrix}.
 \tag{E.T9}\hypertarget{tag-02E3}{}\label{maps:eq:E.T9}
\]
Its monomials are exactly \eqref{maps:eq:E.T2}.  In normal-first order,
$P_{i+2}^{-1}L_{-r_i}P_i$ is the cyclic permutation
matrix.  Thus the framing and triple-cycle assertions hold for all six fan identifications.
Taking collars of the component core, three double-curve
annuli and two triple charts, and rounding their outer
corners, gives the asserted finite plumbing.  These are
subcharts of the smooth toric quotient, so the construction
is Hausdorff and the collar identifications are the indicated
smooth ones.

\item Finally the period translations for $g,u$ are $-e_2,e_1$.
On $q=\eta e^{-2\pi i\theta}$ the marked expression is
$x=\exp(2\pi i((-\theta+iT)B_0a+b))$.
At $\theta=1$ it agrees with $\theta=0$ and
$b$ replaced by $b-B_0a$.  This is exactly the inverse
endpoint convention for
\[
 M_0=\begin{pmatrix}
 1&0&0&0\\0&1&0&0\\0&1&1&0\\-1&0&0&1
 \end{pmatrix}.
 \tag{E.T8}\hypertarget{tag-02E2}{}\label{maps:eq:E.T8}
\]
The comparison therefore preserves the full outer marking.

\end{enumerate}
\end{proof}


\begin{proposition}[The explicit smooth plumbing presentation]\hypertarget{tag-02ML}{}\label{maps:src:E.2.20}
\StatementRecord{maps:src:E.2.20}{proved}{02ML}{\SeriesRef{maps:src:E.2.21}, \SeriesRef{maps:src:E.2.19}, \eqref{maps:eq:E.S1}, \eqref{maps:eq:E.T2}, \eqref{maps:eq:E.U7}, \eqref{maps:eq:E.U10}, \eqref{maps:eq:E.U8}, \eqref{maps:eq:E.U9}, \SeriesRef{maps:src:B.5.1}, \SeriesRef{maps:src:E.2.6}}


The following finite smooth plumbing presentation uses only standard
pieces and specified attaching maps. Its conversion into ordinary
framed handles and the relative smooth identification are given in \ref{maps:src:E.2.21}.

Let $F=\mathbb R^4/\mathbb Z^4$, with basis $(g,u,w,d)$, and use the
two matrices $A_j$ and fixed vectors $v_j$ of \SeriesRef{mapsI:src:A.1.1}.  Set
\[
 (m_1,m_2)=(3,4),\quad(\kappa_1,\kappa_2)=(1,-1),\quad
 w_j=\kappa_jv_j,\qquad L=\ker\gamma,
\]
so that $\Lambda=\mathbb Zw_j\oplus L$ and $A_jw_j=w_j$.
On $L$ the two matrices are
\[
 B_1=\begin{pmatrix}0&1&0\\-1&-1&0\\1&0&1\end{pmatrix},
 \qquad
 B_2=\begin{pmatrix}0&-1&0\\1&0&0\\0&1&1\end{pmatrix}.
 \tag{E.U1}\hypertarget{tag-02E4}{}\label{maps:eq:E.U1}
\]
Define the standard smooth four-manifolds
\[
 \mathcal S_j=(T^3\times\mathbb R)/
       ((l,\sigma+1)\sim(B_j^{-1}l,\sigma)),
 \qquad \mathcal N_j=\mathcal S_j\times D^2_r.
 \tag{E.U2}\hypertarget{tag-02E7}{}\label{maps:eq:E.U2}
\]


The cap boundary is identified with $M_{A_j}$ by the following
explicit map.  Write $x=\phi w_j+l$ and use the interval marking
$(x,1)\sim(A_j^{-1}x,0)$ on $M_{A_j}$.  Then
\[
 [x,\theta]\longmapsto
 \left([l,\sigma=m_j\kappa_j\phi+\theta],
                         \ r=e^{2\pi i\kappa_j\phi}\right).
 \tag{E.U3}\hypertarget{tag-02E8}{}\label{maps:eq:E.U3}
\]
All angles on the right are taken in their stated mapping-torus
quotients.  In the ordered angular variables $(\phi,\theta)$
and $(\sigma,\alpha)$, where $r=e^{2\pi i\alpha}$, the map is
\[
 \begin{pmatrix}\sigma\\\alpha\end{pmatrix}
 =
 \begin{pmatrix}m_j\kappa_j&1\\\kappa_j&0\end{pmatrix}
 \begin{pmatrix}\phi\\\theta\end{pmatrix},
 \qquad
 \begin{pmatrix}\phi\\\theta\end{pmatrix}
 =
 \begin{pmatrix}0&\kappa_j\\1&-m_j\end{pmatrix}
 \begin{pmatrix}\sigma\\\alpha\end{pmatrix}.
 \tag{E.U4}\hypertarget{tag-02E9}{}\label{maps:eq:E.U4}
\]
These are integral unimodular matrices.

To define a compact six-manifold $\mathcal E_{\rm pl}$, take the
flat $F$-bundle over an oriented disc with three open interior
discs removed.  Its three marked inner loop monodromies are
$A_1,A_2,M_0$, with $A_1A_2M_0=I$.  Equivalently cut the planar
surface along three disjoint paths to its outer boundary, take
the product with $F$, and reglue the two banks by these three
integral linear diffeomorphisms.  This specifies smooth band
attachments, including their product normal coordinates.
Fill the first two boundary components by \eqref{maps:eq:E.U2}--\eqref{maps:eq:E.U4}, and
the third by the finite canonical-bundle plumbing of \ref{maps:src:E.2.19},
using its complete marked $M_0$ boundary map.  The remaining
boundary is marked as $F\times S^1_\mu$, with the outer meridian
orientation of \eqref{maps:eq:E.S1}.  The core uses its zero section.

For the regular-fibre complement there is an
orientation-preserving diffeomorphism
\[
                  E\ \cong\ \mathcal E_{\rm pl}
                 \quad\hbox{relative to }F\times S^1_\mu .
 \tag{E.U5}\hypertarget{tag-02EA}{}\label{maps:eq:E.U5}
\]
Consequently the shifted fivefold has the explicit smooth
presentation
\[
 P\ \cong\
 (\mathcal E_{\rm pl}\times T^4)
 \ \cup_{\ (x_1,x_2,\theta)\mapsto(x_1,x_2,-\theta)\ }\
 (T^4\times\mathcal E_{\rm pl}).
 \tag{E.U6}\hypertarget{tag-02EB}{}\label{maps:eq:E.U6}
\]
Every piece in this presentation is a standard torus product,
a linear three-torus mapping torus times an oriented disk, or
the canonical disk bundle of the explicitly specified toric
blow-up with the normal-disk plumbings \eqref{maps:eq:E.T2}.  All attachment
maps have been specified on their normal coordinates.

\end{proposition}


\begin{proof}
\begin{enumerate}[label={$\langle1\rangle$\arabic*.},leftmargin=*]
\item We verify first that \eqref{maps:eq:E.U2}--\eqref{maps:eq:E.U4} are the finite caps,
including the normal framing.  The original cap is
\[
 (F\times D^2_s)/
 \bigl((x,s)\sim(A_jx+v_j/m_j,e^{-2\pi i/m_j}s)\bigr).
 \tag{E.U7}\hypertarget{tag-02EC}{}\label{maps:eq:E.U7}
\]
The good-reduction real flat marking is smooth on the entire
root disc, so these are smooth coordinates even at $s=0$.
Because $\gamma A_j=\gamma$ and $\gamma(v_j)=\kappa_j=\pm1$,
\[
                  r=s\,e^{2\pi i\kappa_j\gamma(x)}
 \tag{E.U8}\hypertarget{tag-02ED}{}\label{maps:eq:E.U8}
\]
is invariant under the displayed action.  It is single-valued
on $F$, has the same absolute value as $s$, and is complex
linear on every normal disk.  The induced map
$[x,s]\mapsto([x],r)$ identifies \eqref{maps:eq:E.U7} smoothly with
$S_j\times D^2$.  An inverse is obtained by choosing any
representative $x$ and setting
$s=r e^{-2\pi i\kappa_j\gamma(x)}$; invariance makes this
independent of the representative in the quotient.

\item In the integral basis $(w_j,u,w,d)$ the affine generator is
\[
                  (\phi,l)\longmapsto
                     (\phi+\kappa_j/m_j,B_jl).
 \tag{E.U9}\hypertarget{tag-02EE}{}\label{maps:eq:E.U9}
\]
Thus $\sigma=m_j\kappa_j\phi$ gives exactly the mapping torus
in \eqref{maps:eq:E.U2}.  This also supplies a direct inverse to the product
description: choose a representative $(l,\sigma)$, put
$x=(\kappa_j\sigma/m_j)w_j+l$ and
$s=r e^{-2\pi i\sigma/m_j}$.  Replacing that representative
by the defining mapping-torus relation changes $(x,s)$ by
the finite affine generator and hence not its class.

\item The boundary parametrization of the finite cap is
\[
 [x,\theta]\longmapsto
       [x+\theta v_j/m_j,e^{-2\pi i\theta/m_j}].
 \tag{E.U10}\hypertarget{tag-02E5}{}\label{maps:eq:E.U10}
\]
Substitute \eqref{maps:eq:E.U10} into \eqref{maps:eq:E.U8}--\eqref{maps:eq:E.U9}.  The coefficient of
$\theta$ in the normal angle is
$-1/m_j+\kappa_j^2/m_j=0$, while the base angle is
$m_j\kappa_j\phi+\theta$.  This proves \eqref{maps:eq:E.U3} and \eqref{maps:eq:E.U4}.
At $\theta=1$ that base angle advances by one, precisely
replacing $l$ by $B_j^{-1}l$; this checks the mapping-torus
seam.  A lattice change of $\phi$ changes the base angle by
$m_j\kappa_j$, and $B_j^{m_j}=I$, so it introduces no further
ambiguity.

\item The function of the cap that gives its original base parameter
is, in these smooth coordinates,
\[
                         q=r^{m_j}e^{-2\pi i\sigma}.
 \tag{E.U11}\hypertarget{tag-02E6}{}\label{maps:eq:E.U11}
\]
On the marked boundary it has angle
$m_j\alpha-\sigma=-\theta$, the original clockwise base
convention.  A positive normal meridian at fixed
$(l,\sigma)$ has $(\Delta\phi,\Delta\theta)=(\kappa_j,-m_j)$.
Since $\kappa_jw_j=v_j$, it is exactly $t_{v_j}\rho_j^{-m_j}$,
recovering the filling relation
$\rho_j^{m_j}=t_{v_j}$.

\item Over the punctured planar core, the family has the
smooth real flat marking with linear monodromies $A_j$ and
its zero section.  The logarithmic finite gluing has exactly
the drift in \eqref{maps:eq:E.U10}; after subtracting it, the remaining
translation has a single-valued real lift and is removed by
the collar isotopy already proved in \SeriesRef{maps:src:B.5.1}.    At the cusp use the relative smooth diffeomorphism of \ref{maps:src:E.2.19}.  Each comparison
has precisely the prescribed boundary value and can be
made product on a smaller collar.  They therefore glue to
the relative diffeomorphism \eqref{maps:eq:E.U5}.  In particular the
normal framing of the unfilled ordinary fibre is unchanged.

\item Finally use the two hemispheres in the shifted
fibre-product decomposition of \ref{maps:src:E.2.6}.  On each hemisphere
the other factor is a regular marked torus bundle over a
disc, hence a product preserving its zero section and flat
marking.  Their common boundary has unchanged torus
coordinates and opposite normal angles.  Taking the product
of \eqref{maps:eq:E.U5} with those regular factors and gluing gives
\eqref{maps:eq:E.U6}, with its induced complex orientation.  
\end{enumerate}
\end{proof}



\begin{theorem}[Ordinary framed handles and the smooth type]\hypertarget{tag-02MM}{}\label{maps:src:E.2.21}
\StatementRecord{maps:src:E.2.21}{proved}{02MM}{\eqref{maps:eq:E.T4}, \SeriesRef{maps:src:E.2.20}, \eqref{maps:eq:E.U3}, \eqref{maps:eq:E.U4}, \SeriesRef{maps:src:E.2.19}, \eqref{maps:eq:E.T6}, \eqref{maps:eq:E.H13}; \cite[Section~4.1, equation~(16), p.~103]{mapsVogel09}; \cite{mapsMilnor65}}


The following finite ordinary framed-handle prescription defines a
compact oriented six-manifold $\mathcal H_6$ with marked boundary
$T^4\times S^1$, and a closed oriented ten-manifold $\mathcal H_{10}$.
There are orientation-preserving diffeomorphisms
\[
 E\cong\mathcal H_6\quad\hbox{relative to the full marked boundary},
 \qquad P\cong\mathcal H_{10}.
 \tag{E.H0}\hypertarget{tag-02AQ}{}\label{maps:eq:E.H0}
\]
The respective numbers of ordinary handles, in increasing index, are
\[
 (2,5,14,22,21,10,2),\qquad
 (4,31,119,290,485,574,485,290,119,31,4).
\]
The attaching embeddings and their normal coordinates are given below.  
The balls and complex normal charts carry their complex orientations;
orient the product blocks by their displayed coordinate embeddings.
For example the dense block below has orientation
$-d\arg X\wedge d\arg Y\wedge d\mu_1\wedge d\mu_2
\wedge d\operatorname{Re}q\wedge d\operatorname{Im}q$.
This fixes the orientation convention even when the chosen order of
angle and moment coordinates differs from the complex coordinate order.

\emph{1. The two triple-point balls.}
Use two complex coordinate balls, denoted $V_+$ and $V_-$, with
ordered coordinates $Z=(Z_0,Z_1,Z_2)$ and $W=(W_0,W_1,W_2)$.
The boundary faces are rounded coordinate faces of
\[
 |Z_i|\leq\epsilon_v,\quad |Z_0Z_1Z_2|\leq\eta,
 \qquad
 |W_i|\leq\epsilon_v,\quad |W_0W_1W_2|\leq\eta.
 \tag{E.H1}\hypertarget{tag-02AR}{}\label{maps:eq:E.H1}
\]
They are smooth six-balls.  Throughout the prescription, rounding is
in the nonnegative radii, preserves the angular coordinates, and is
fixed on the attaching collars.  Positive changes of the radii only
give the corresponding collar isotopy.  The complex products in
\eqref{maps:eq:E.H1} are the same function $q$.
The coordinate formulas may first be read on manifolds with corners;
apply the same radial collar reparametrization on the two sides of
each attaching face before smoothing their union.  This convention
also specifies the maps at a corner meeting $|q|=\eta$.

\emph{2. The three annular blocks, with their endpoint framings.}
For $i=0,1,2$ use a block $S^1\times D^1\times D^4$ with coordinates
\[
 z=\exp(s+2\pi i\theta),\quad -R\leq s\leq R,
 \qquad (a,b)\in D^4,\qquad q=zab.
 \tag{E.H2}\hypertarget{tag-02AW}{}\label{maps:eq:E.H2}
\]
The normal four-ball here is a rounded radial domain
$|a|,|b|\leq\epsilon_e$, $|ab|\leq\eta/|z|$;
its positive radial trivialization identifies it with $D^4$.
The ends attach to framed tubes about the following coordinate
circles.  A table entry $(Z_j;Z_k,Z_l)$ means that the curve
coordinate is $Z_j$ and the two ordered complex normal coordinates
are $Z_k,Z_l$; the same convention applies to $W$.
\[
\begin{array}{c|c|c}
 i&z=0\hbox{ end: }(z;a,b)&z=\infty\hbox{ end: }(w;A,B)\\ \hline
0&(Z_2;Z_1,Z_0)&(W_0;W_1,W_2)\\
1&(W_2;W_1,W_0)&(Z_1;Z_2,Z_0)\\
2&(Z_0;Z_2,Z_1)&(W_1;W_2,W_0)
\end{array}
\tag{E.H3}\hypertarget{tag-02AX}{}\label{maps:eq:E.H3}
\]
The full-dimensional endpoint map between the two end charts is
\[
                 (w,A,B)=(z^{-1},za,zb).
 \tag{E.H4}\hypertarget{tag-02AY}{}\label{maps:eq:E.H4}
\]
In particular the second attaching circle has angle $-\theta$,
and its two complex normal vectors each rotate by $+2\pi\theta$.
The factors $|z|$ are the positive radial rescalings.
  Every coordinate circle
of the two balls occurs exactly once in \eqref{maps:eq:E.H3}.  The maps are
embeddings of $S^1\times\partial D^1\times D^4$.

\emph{3. The dense torus block and its full sector maps.}
Attach $T^2\times D^2\times D^2_q$ along
$T^2\times\partial D^2\times D^2_q$.  Here is a completely explicit
choice of its base disk and the attaching map.  Use the rays
$r_i$ of \eqref{maps:eq:E.T4} and the hexagon
\[
 \mathcal P=\{v\in\mathbb R^2:\langle v,r_i\rangle\geq-1\},
 \quad
 \operatorname{Vert}(\mathcal P)=
 \{(-1,-1),(0,-1),(1,0),(1,1),(0,1),(-1,0)\}.
 \tag{E.H5}\hypertarget{tag-02AZ}{}\label{maps:eq:E.H5}
\]
On the positive real part of $(\mathbb C^*)^2$, use the coordinates
given by the toric moment map
\[
 \mu(X,Y)=
 \frac{\displaystyle\sum_{v\in\operatorname{Vert}(\mathcal P)\cup\{0\}}
                 |X^{v_1}Y^{v_2}|^2v}
      {\displaystyle\sum_{v\in\operatorname{Vert}(\mathcal P)\cup\{0\}}
                 |X^{v_1}Y^{v_2}|^2}.
 \tag{E.H6}\hypertarget{tag-02B0}{}\label{maps:eq:E.H6}
\]
Choose a small inward parallel hexagon, with its corners rounded
symmetrically, $\mathcal P_\epsilon\subset\operatorname{int}\mathcal P$.
The inverse image is $T^2\times\mathcal P_\epsilon$; its two circle
coordinates are $\arg X,\arg Y$.  On the boundary disk circle put
$(X,Y)$ equal to the positive inverse of $\mu$, multiplied by these
two phases.  On the six sectors of that boundary, the attaching
map into the preceding blocks is the following monomial map, with
the whole complex normal coordinate $q$ included:
\[
\begin{array}{c|ccc}
 i& n_i&a_i&a_{i+1}\\ \hline
0&q/(XY)&X&Y\\
1&q/Y&XY&X^{-1}\\
2&qX&Y&(XY)^{-1}\\
3&qXY&X^{-1}&Y^{-1}\\
4&qY&(XY)^{-1}&X\\
5&q/X&Y^{-1}&XY
\end{array}
\qquad q=n_i a_i a_{i+1}.
\tag{E.H7}\hypertarget{tag-02B1}{}\label{maps:eq:E.H7}
\]
Use the cyclic identifications
$(n_2,a_2,a_3)=(Z_1,Z_2,Z_0)$,
$(n_4,a_4,a_5)=(Z_2,Z_0,Z_1)$ and the analogous two cyclic
permutations of $W$ for vertices $3,5$.  Vertices $0,1$ are $Z,W$.
The chart at an edge is $(z,a,b)=(a_{i+1},a_i,n_i)$, with \eqref{maps:eq:E.H4}
at its other end.  These rules specify \eqref{maps:eq:E.H7} also where a sector
meets an annular block.  The transition on a common sector is the
monomial map $P_{i+1}^{-1}P_i$, where
$P_i=((0,1),(r_i,1),(r_{i+1},1))$; the sum of its exponent rows
is $(1,1,1)$, so it preserves $n_i a_i a_{i+1}$.  Thus the sector maps specify the
normal framing as well as the core map.  

\emph{4. Expansion into ordinary handles.}
We use the following fixed elementary rule on every parameterized
block in this prescription:
\[
 T^r\times D^k\times D^{n-r-k}
 \quad\hbox{attached on}\quad
 T^r\times\partial D^k\times D^{n-r-k}.
 \tag{E.H8}\hypertarget{tag-02B2}{}\label{maps:eq:E.H8}
\]
Cut each ordered parameter circle into its two closed semicircles,
with small endpoint collars.  The first semicircle is its zero
handle and the second is its one handle.  For a bit vector
$\varepsilon\in\{0,1\}^r$, the resulting ordinary handle is of
index $k+|\varepsilon|$.  Attach by increasing $|\varepsilon|$,
then lexicographically.  On the original attaching face use
the displayed full-dimensional attaching embedding of the block;
on a newly exposed parameter face use the identity endpoint map
in that circle.  The normal coordinates are the remaining ordered
semicircle coordinates followed by the original normal coordinates.
Corner rounding and the permutation putting negative coordinates
first complete an ordinary framed attaching embedding
$S^{k+|\varepsilon|-1}\times D^{n-k-|\varepsilon|}$.
  For one parameter
it is the literal two-piece decomposition in
\cite[Section~4.1, equation~(16), p.~103]{mapsVogel09}.


Apply \eqref{maps:eq:E.H8} first to \eqref{maps:eq:E.H2}, with $(r,k,n)=(1,1,6)$, and then
to the core, with $(r,k,n)=(2,2,6)$.  Together with the two balls
this gives cusp handle counts $(2,3,4,2,1)$.

\emph{5. The torus bands and the finite caps.}
Take the marked flat torus bundle of \ref{maps:src:E.2.20} over the disc with
three holes, regarding the cusp boundary as the incoming boundary.
Cut along the two marked disjoint arcs from the two finite boundary
circles to the outer circle.  The cut surface is an annulus;
its torus bundle is the collar of the already constructed cusp
boundary.  Restoring the two strips attaches two blocks
\[
                 T^4\times D^1\times D^1.
 \tag{E.H9}\hypertarget{tag-02B3}{}\label{maps:eq:E.H9}
\]
To fix their maps without a planar drawing, use the oriented ribbon
polygon whose successive based inner boundary words are
$A_1,A_2,M_0$ and whose outer word is their product.  The negative
bank of the $j$th cut has coordinate $(x,t)$, the positive bank
has $(A_jx,t)$, and the strip uses the product normal coordinate.
Reversing a bank orientation replaces that matrix by its inverse.
This is the cut-polygon definition of the flat bundle in \ref{maps:src:E.2.20};
the identity $A_1A_2M_0=I$ specifies the unfilled outer boundary.
There is no affine term.  Thus \eqref{maps:eq:E.H9} has its full attaching
embedding, including the transverse interval, fixed by these
linear maps.  Expand each by \eqref{maps:eq:E.H8}, with $(r,k,n)=(4,1,6)$.
The unused boundary circles are precisely the two finite boundaries
and the marked ordinary boundary.
For clarity, the initial bank coordinates on the already attached
cusp boundary are the explicit coordinates of \cref{maps:src:E.2.19}:
\[
 \begin{gathered}
 q=\eta e^{-2\pi i\theta},\qquad T=-\log\eta/(2\pi),\\
 (X,Y)=\exp\bigl(2\pi i((-\theta+iT)B_0a+b)\bigr),
 \qquad a=(g,u),\ b=(w,d).
 \end{gathered}
 \tag{E.H9a}\hypertarget{tag-02B4}{}\label{maps:eq:E.H9a}
\]
The exponential is componentwise; at an interval seam use
$b\mapsto b-B_0a$.  Equations \eqref{maps:eq:E.H7} and \eqref{maps:eq:E.H9a} therefore
give the bank embeddings into the existing handle
boundary, including their derivatives.  

For the $j$th finite cap use its already explicit
$\mathcal S_j\times D^2_r$ and the entire boundary embedding
\eqref{maps:eq:E.U3}.  The mapping torus $\mathcal S_j$ has the following product
block construction: a block $T^3\times D^1$ over the first
semicircle of $\sigma$, followed by a block $T^3\times D^1$ over
the second semicircle, with endpoint maps $I$ and $B_j^{-1}$.
The first is a three-parameter zero handle, the second a
three-parameter one handle.  As a filling relative to
$\mathcal S_j\times S^1_r$, the normal disk supplies two negative
directions.  Consequently the relative cap consists of
\[
 \begin{array}{c|c|c}
 \hbox{base-circle block}&(r,k,n)&\hbox{full attaching data}\\ \hline
 0&(3,2,6)&\hbox{\eqref{maps:eq:E.U3}, the first semicircle, and the $r$-disk}\\
 1&(3,3,6)&\hbox{\eqref{maps:eq:E.U3}, endpoint maps $I,B_j^{-1}$, and the $r$-disk}.
 \end{array}
 \tag{E.H10}\hypertarget{tag-02AS}{}\label{maps:eq:E.H10}
\]
More explicitly an attaching face in the normal direction is sent
to \eqref{maps:eq:E.U3} at $r=e^{2\pi i\alpha}$, and an attaching face in a
base direction is the identical face of the preceding base block,
crossed with $D^2_r$.  At the base seam use $B_j^{-1}$ on $l$ and
the identity on $r$.  The inverse angular formula \eqref{maps:eq:E.U4} supplies
the full map back to the boundary being filled.  In particular
the normal meridian is $v_j-m_j\rho_j$.
Apply \eqref{maps:eq:E.H8} to the two blocks in \eqref{maps:eq:E.H10}.  Each cap adds the
ordinary relative counts $(1,4,6,4,1)$ in indices $2,\ldots,6$.
This completes the definition of $\mathcal H_6$.

\emph{6. The closed ten-dimensional handlebody.}
For each handle of $\mathcal H_6$, take its product with the four
ordered torus circles and apply \eqref{maps:eq:E.H8}.  This gives
$\mathcal H_6\times T^4$ with its full boundary framing.  On the
other side take the dual relative handle decomposition of
$T^4\times\mathcal H_6$: reverse the handle order and exchange
the negative and positive disks of every handle.  A handle of
index $j$ becomes one of index $10-j$, and its attaching map is
the belt embedding of the original handle.  Glue its
initial boundary by the explicit map
\[
                 (x_1,x_2,\theta)\longmapsto
                                    (x_1,x_2,-\theta).
 \tag{E.H11}\hypertarget{tag-02AT}{}\label{maps:eq:E.H11}
\]
All subsequent dual attaching maps are obtained by this fixed
exchange of the two coordinate disks, in reverse order, on the
previously specified embeddings.  Use the factor order $E,T^4$
on the first half and $T^4,E$ on the second; their interchange
has sign $(-1)^{24}=1$.  The boundary-angle reversal in \eqref{maps:eq:E.H11}
is the orientation-reversing boundary gluing.  This defines
$\mathcal H_{10}$ entirely by ordinary handles.  If
\[
 c(t)=2+5t+14t^2+22t^3+21t^4+10t^5+2t^6,
 \qquad
       c_{10}(t)=(c(t)+t^6c(t^{-1}))(1+t)^4,
 \tag{E.H12}\hypertarget{tag-02AU}{}\label{maps:eq:E.H12}
\]
the coefficients are the counts in the statement.  Expanding these polynomials gives $76$ handles in dimension six
and $2432$ in dimension ten.

\end{theorem}


\begin{proof}
\begin{enumerate}[label={$\langle1\rangle$\arabic*.},leftmargin=*]
\item The main geometric point is that the first three steps give
\emph{smooth blocks of the cusp tube}.  We give the collar
argument explicitly.  Work in the smooth toric quotient of \ref{maps:src:E.2.19}.
Choose disjoint small coordinate neighborhoods of the two triple
points.  The three incident coordinate axes in each neighborhood
are exactly the six ends in \eqref{maps:eq:E.H3}.  Remove smaller endpoint disks
from each normalized double curve.  Its remaining annulus has
the two complex normal lines $\mathcal O(-1)|_{\rm annulus}$,
which are trivialized by $(a,b)$ in \eqref{maps:eq:E.T6}.  The resulting normal
domain in the tube is exactly the domain in \eqref{maps:eq:E.H2}.
Every ray in $\mathbb C^2$ meets its rounded boundary once and
transversely.  Indeed the three defining radius functions have
positive radial derivative wherever the corresponding boundary
face is active.  Positive radial rescaling therefore gives a
smooth family of four-balls as $z$ ranges over this compact
annulus.  Choose the radial change to be constant linear scaling
on a common sufficiently small ball about the origin and interpolate
monotonically only away from that ball; it is then smooth at the
origin as well as at the boundary and preserves the stated normal
frames.  The same argument for the coordinate and product
radius functions in \eqref{maps:eq:E.H1} gives six-balls at the triple points.
Rounding can be chosen constant on the end attaching collars.

\item   At
vertex $i$ the normal-first coordinates are $(n_i,a_i,a_{i+1})$.
At vertex $i+2$ the deck identification is their cyclic
permutation, and the threefold composition is the identity by
\ref{maps:src:E.2.19}.  At the two ends of $C_i$ the coordinates are
$(z,a,b)=(a_{i+1},a_i,n_i)$ and
$(w,A,B)=(a_{i-1},a_i,n_{i-1})$.  The ordinary
$\mathcal O(-1)^2$ transition is \eqref{maps:eq:E.H4}.  Substituting the two
cyclic permutations gives exactly \eqref{maps:eq:E.H3}.  This checks all six ends.

\item We next choose the stratum collars so that their remaining
interface is the entire attaching face of the dense block.
On the normalized toric surface take a small collar of its
hexagonal toric boundary, choosing the corner neighborhoods
first and the edge strips second.  Choose the strips narrower
than the corner collars and invariant under the opposite-edge
inversion and the two cyclic triple identifications.  The
uncovered moment region is the rounded interior disk
$\mathcal P_\epsilon$.  This can be done directly with six
inward parallel affine inequalities in \eqref{maps:eq:E.H5}, smoothing each
corner inside its already chosen corner neighborhood.  Thus the
common frontier in the smooth part of the normalization is
exactly $T^2\times\partial\mathcal P_\epsilon$.

\item The map \eqref{maps:eq:E.H6} gives the claimed coordinates: in logarithmic
positive coordinates it is the gradient of a log-sum-exponential
function.  Its Hessian is twice the covariance matrix of the
seven affinely spanning exponent vectors, hence positive
definite.  Strict convexity and the growth in each supporting
direction give a diffeomorphism onto
$\operatorname{int}\mathcal P$.  The inverse on the compact
disk and its frontier is smooth.  The canonical line on this
free torus locus is trivial; more specifically the toric
function $q$ is a nonzero toric factor times its normal
coordinate.  Hence \emph{$q$ itself}, including its phase,
trivializes that normal line.  The portion of the truncated
tube over this disk is therefore exactly
\[
 (T^2\times\mathcal P_\epsilon)\times D^2_q,
 \quad\hbox{with interface}\quad
 (T^2\times\partial\mathcal P_\epsilon)\times D^2_q.
 \tag{E.H13}\hypertarget{tag-02AV}{}\label{maps:eq:E.H13}
\]
In extending the preceding moment cuts off the zero fibre,
hold their equations constant in this $q$ coordinate on an
open neighborhood of the compact frontier.  On their ends
use the product cuts of the annular and triple charts.  A
collar interpolation between these choices is an interpolation
of positive radial cut functions, fixed on each already chosen
end collar.  Their radial derivatives remain positive, so its
flow is a collar isotopy.  The cuts can consequently be chosen
to agree with \eqref{maps:eq:E.H13} exactly.  In particular no twisting of
the $q$ disk is introduced on the six sector seams.

\item These are blocks in a sufficiently small truncated
neighborhood of the compact central fibre.  Their central
parts cover that fibre: two triple points, the three open
double annuli and the remaining normalized surface core.
Compactness from \ref{maps:src:E.2.19} permits one common sufficiently small
$\eta$: the union of their open chart neighborhoods contains
the compact zero fibre, so properness implies that it contains
the entire inverse image of a sufficiently small closed
$q$-disc.  The nested cuts partition this entire inverse image.
If a larger originally chosen tube is used, the
difference is an exterior collar: on the compact annular
region without critical points choose a vector field lifting
the radial direction of $q$, and integrate it.  On the
prescribed outer collar choose the lift in the marked real
torus coordinates of \ref{maps:src:E.2.19}.  It is then the identity on all
those marked coordinates during radial collar transport.


\item The embedding in \eqref{maps:eq:E.H13} is \eqref{maps:eq:E.H7}.  To check this,
the exponent matrix for a unimodular chart is $P_i^{-1}$,
applied to $(X,Y,q)$.  This gives each of the six displayed
rows.  Consecutive changes are $P_{i+1}^{-1}P_i$ and telescope
to the identity.  The translations identifying the two
triple orbits are exactly those already used in \eqref{maps:eq:E.H3}.
All six normal coordinate multipliers in \eqref{maps:eq:E.H7} are thereby
fixed, and the sector seams are equal maps on overlaps.  Positive rescaling of the chosen
radial domains leaves these monomials unchanged before
the collar reparametrization.  This completes the smooth,
fully framed identification of the cusp with steps 1--3.

\item The handle expansion rule follows from the product decomposition
of the parameter circles.  For one circle, write $S^1=I_0\cup I_1$.
First attach $I_0\times D^k\times D^{n-k-1}$; after rounding it
is a $k$-handle.  The remaining block is
$(I_1\times D^k)\times D^{n-k-1}$, a $(k+1)$-handle.
Its attaching face consists of the original negative face
and its two endpoint faces in the first handle.  The stated
maps on these three faces agree on their corners because
they are restrictions of the same product coordinates.
Iterate this argument on the ordered circle factors.
It proves \eqref{maps:eq:E.H8}, including all newly created attaching
faces and their normal framings.

\item The planar core relative to its incoming circle has two
ordinary surface one-handles.  Its flat bundle is a product
over each strip, so restoring the two cut strips is exactly
\eqref{maps:eq:E.H9}.  The zero section and the full lattice marking fix
the normal coordinates on these attachments.  Each cap in
\eqref{maps:eq:E.H10} is the smooth product of \ref{maps:src:E.2.20}.  Its relative
handle claim can also be checked by the function
$-|r|^2+\varepsilon f(l,\sigma)$, with a cutoff near the
incoming boundary: choose $f$ by a two-critical-point
function of $\sigma$, perturbed on its two critical tori
by the product Morse function on $T^3$.  On the remaining
arcs its base derivative dominates the perturbation, so
there are no further critical points.  The indices are
two plus the base indices, as in \eqref{maps:eq:E.H10}.    Summing the blocks
gives
\[
 (2,3,4,2,1,0,0)+2(0,1,4,6,4,1,0)
                     +2(0,0,1,4,6,4,1)
       =(2,5,14,22,21,10,2).
\]
Every comparison has its prescribed boundary value, so the
relative diffeomorphism of \ref{maps:src:E.2.20} gives the first assertion
of \eqref{maps:eq:E.H0}.  Taking the products, the dual relative handles
and the boundary flip of \ref{maps:src:E.2.20} gives the second.
The handle polynomials evaluated at $-1$ give Euler numbers $2$ and $0$.

\item The following relative uniqueness theorems determine the smooth
type from the full attaching embeddings.  The relevant relative uniqueness
results are Milnor's Theorem~3.13, which identifies an elementary
cobordism from its full characteristic embedding, Theorem~1.4,
which makes smooth composition unique with the marked boundary
fixed, and the isotopy-extension Theorem~5.8
\cite{mapsMilnor65}.  In the present notation the characteristic
embedding is the whole map
$S^{j-1}\times D^{n-j}\longrightarrow\partial W$.  Its normal coordinates
are part of \eqref{maps:eq:E.H3}--\eqref{maps:eq:E.H11}.  The elementary consequence needed
here can also be seen directly.  If $\phi_0,\phi_1$ are joined
by the isotopy of full attaching embeddings, extend
the velocity field on their images to a slightly larger boundary
collar and multiply it by a cutoff which is zero on the protected
marked boundary.  Its flow extends the attaching isotopy to an
ambient boundary isotopy.  Extending this flow across a product
collar of $W$, and using the identity on the new coordinate handle,
identifies $W\cup_{\phi_0}(D^j\times D^{n-j})$ with
$W\cup_{\phi_1}(D^j\times D^{n-j})$.  The same cutoff argument
applies to the small radial corner changes above.  This derives
the relative, framed version from the stated isotopy-extension
theorem.

\item Apply this comparison successively in the displayed block and
sector order.  The initial ball comparisons are their coordinate
identifications.  At each next stage the attaching map
and the prescribed attaching map agree in the computed normal
coordinates, or differ by the explicitly described positive
collar isotopy.  The comparison of that stage is consequently
the identity on the new coordinate handle and the product-collar
extension on the previous stages.  Theorems~3.13 and~1.4 give an
inductive relative diffeomorphism of the whole six-manifold.
The final ordinary boundary coordinates are unchanged.  The
same argument, with the belt embeddings and the boundary flip, applies to every handle on the dual ten-dimensional
side.  In particular a top handle keeps its entire attaching map, so no homotopy-sphere ambiguity arises.  The computed
ordered full framed attaching data, under this relative uniqueness
theorem, are the classifying data used for \eqref{maps:eq:E.H0}.
\end{enumerate}
\end{proof}

\clearpage
